\section{Theorem S II: carriers and Fourier decay}\label{sec:pf:thmS-carrier}

\providecommand{\thmScarrierPt}{\mathrm{Pt}}
\providecommand{\thmScarrierNst}{N^{\mathrm{st}}}
\providecommand{\thmScarrierNz}{N^{0}}
\providecommand{\thmScarrierNhat}{\widehat{N}}
\providecommand{\thmScarrierMaj}{\operatorname{maj}}
\providecommand{\thmScarrierWt}{\operatorname{wt}}
\providecommand{\thmScarrierCnt}{\operatorname{cnt}}
\providecommand{\thmScarrierNtop}{N_{\mathrm{top}}}
\providecommand{\thmScarrierWb}{\mathcal{W}}
\providecommand{\thmScarrierEn}{\mathcal{E}}
\providecommand{\thmScarrierld}{\operatorname{ld}}
\providecommand{\thmScarrierDat}{\mathfrak{D}}

This section proves the \emph{carrier input} $\sfname{CarrierP}$ of Theorem~S (\cref{thm:theorem-S}), the lower bound for
``Poincar\'e amplitudes'' used in the pointwise form of Theorem~S (\cref{sec:pf:thmS}). There the off-diagonal leading
coefficient $K_0(u,u')$ is bounded through an abstract recursion over towers of \emph{points}, which controls sums
$\sum_\rho\|\sum_pc_\rho(p)V(p)\|^2$ of Hilbert--Schmidt operators $V(p)$ attached to points $p$ by the energy $j$ of an
archimedean test function, up to a factor $1+o(1)$. The recursion needs a list of \emph{real inputs}: positivity of certain
Gram kernels, a Fourier decay estimate (a ``plane-wave bound'') for the Fourier transforms of these kernels in their last
variable, two families of off-diagonal estimates (``lift'' and ``block'' estimates), a Poisson-summation datum, a limit
formula for the energy at level $0$, and the identification of the Gram kernel at the top level with the Hilbert--Schmidt
inner products of the $V(p)$. All of these are archimedean statements: the Gram kernels are products, over the complex
places $v$ of a CM field $E$, of local kernels built from Bessel transforms of test functions on a unitary group
$\Uni(H_{k,r})\subset\GL_{2r+k+1}(\C)$.

\Cref{pf:thmS-carrier:main} states that for all global data satisfying the size, count and pinning hypotheses of
Theorem~S there is an archimedean \emph{Hida setting} attached to the data, whose energy datum satisfies the real inputs
with a constant $C_4$ that depends on nothing, with energy equal to the energy constant $j_X$ of the mass data,
and with the operators at point representatives as top vectors. The proof has three layers.
\begin{enumerate}[label=(\alph*)]
\item \emph{Assembly} (\cref{sec:pf:thmS-carrier:assembly,sec:pf:thmS-carrier:parts}). The real inputs are products over
  places of local statements; positivity is propagated by the Schur product theorem, the decay estimates by a
  combinatorial ``big places'' argument, the energy by a product of local limits. This reduces everything to local
  statements at \emph{built places} (\cref{pf:thmS-carrier:compact}).
\item \emph{Local analysis at built places} (\cref{sec:pf:thmS-carrier:built,sec:pf:thmS-carrier:model}). The local
  estimates are first proved for compactly supported test functions and then passed to the limit
  (\cref{pf:thmS-carrier:core}); for compactly supported coefficients they come from Lagrange extraction of top coefficients
  of polynomial dependences on the letters, Loewner-order estimates, Fourier inversion with third differences, and Gram
  comparisons of germs (\cref{sec:pf:thmS-carrier:leaves}). The local energy is identified with the energy of a \emph{model
  function} on a domain in $\C^r\times\C$ by a gauge transfer (\cref{pf:thmS-carrier:modelplace}), whose properties are
  derived from the realisation of the model operators by differential operators, the transfer of adjoints, and a
  polynomiality statement for the radial shift operators.
\item \emph{The energy limit} (\cref{sec:pf:thmS-carrier:LE}). For the model functions, an abstract regularised energy
  identity (\cref{pf:thmS-carrier:LE}) shows that the cut-off energy of a word $w$ tends to $\mathrm{size}(w)\cdot j$; its
  proof is an a priori estimate by absorption, combining a Leibniz formula, derivative bounds obtained from Cauchy's
  integral formula, and Markov's inequality for polynomials.
\end{enumerate}
The input on the global side (\cref{sec:pf:thmS-carrier:X}) is the construction of the Hida setting from the spectral
data of the global datum; it uses the uniqueness of Bessel models for discrete series with large gap and the existence of
$K$-finite vectors with non-zero Bessel functional.

the part of the formal proof treated here is large (about $67\,000$ lines). The arguments of layers~(a) and~(c) are given
in full. For layer~(b) we prove the model-place theorem from six transfer statements and three of these from deeper local
identities, and we prove the built bounds from the local wave, lift, top and germ statements, of which the degree, top,
uniform-decay, germ-rate and coset-sum statements are proved here; the remaining local identities are cited by their formal
names and are converted elsewhere.

\paragraph{Conventions.}
All Hilbert spaces are complex, with inner products linear in the first variable. A kernel $\kappa$ on a finite set $Y$
is \emph{positive semidefinite} (PSD) if the matrix $(\kappa(y,y'))_{y,y'\in Y}$ is Hermitian positive semidefinite. A
\emph{probability vector} on a finite set $U$ is a function $\theta\ge0$ vanishing outside $U$ with $\sum_U\theta=1$.
``For $\eps\to0^+$'' refers to the filter of right neighbourhoods of $0$.

% ---------------------------------------------------------------------------
\subsection{The objects}\label{sec:pf:thmS-carrier:objects}

\subsubsection{Points and energy data}
Let $E$ be a field. A \emph{point} is a finite sequence $p=((s_1,t_1),\dots,(s_m,t_m))$ of pairs $(s_j,t_j)\in\N\times
E$; its \emph{level} is $|p|=\sum_js_j$. We write $\thmScarrierPt(E)$ for the set of points, $\emptyset$ for the empty
point and $p\cdot(s,t)$ for the point obtained by appending $(s,t)$. The point $p$ is a \emph{point of level $i$} if
$|p|=i$, $s_j\ge1$ and $t_j\ne0$ for all $j$.

\begin{definition}[Energy data]\label{pf:thmS-carrier:energydatum}
An \emph{energy datum} $\thmScarrierDat$ over $E$ with word set $W$ consists of integers $r,f,K\ge0$; functions
$|\cdot|_v\colon E\to\R$ ($v<f$); weights $\thmScarrierWt\colon W\to\R$ and counts $\thmScarrierCnt\colon W\to\N$; a
finite set $W_{\mathrm{all}}\subset W$ containing an element $\mathrm{emp}$ with $\thmScarrierCnt(\mathrm{emp})=0$ and
$\thmScarrierWt(\mathrm{emp})=1$; \emph{Gram kernels} $\gamma_{\eps,i,w}(p,q)\in\C$ ($\eps\in\R$, $i\in\N$, $w\in W$,
$p,q\in\thmScarrierPt(E)$); kernels $\mathcal K_i(p,q)$ with $\mathcal K_0(\emptyset,\emptyset)=1$, theta kernels
$\Theta_{i,s,t}(p,q)$ and numbers $e_i\in\R$; additive subgroups $\Lambda_{b,s},\Lambda'_{b,s},\Lambda^\vee_b$ of $E$;
\emph{Fourier coefficients} $\thmScarrierNhat_{\eps,i,w}(Y,\kappa;l)\in\C$ for finite $Y\subset\thmScarrierPt(E)$,
kernels $\kappa$ and $l\in E$; and \emph{sizes} $\rho^D_{i,v},\rho^L_{b,s,v},V_{b,s}\in\R$. Put
$\thmScarrierNtop=4K(r+1)$, $\thmScarrierWb(n)=\{w\in W_{\mathrm{all}}:\thmScarrierCnt(w)\le n\}$,
\begin{gather*}
\thmScarrierNst_{\eps,i,s,w}(Y,\kappa;t,t')=\frac1{\thmScarrierWt(w)}\sum_{h,h'\in Y}\kappa(h,h')\,
  \gamma_{\eps,i,w}\bigl(h\cdot(s,t),h'\cdot(s,t')\bigr),\\
\thmScarrierNz_{\eps,k,w}(Y,\kappa)=\frac1{\thmScarrierWt(w)}\Re\sum_{h,h'\in Y}\kappa(h,h')\gamma_{\eps,k,w}(h,h'),\\
\thmScarrierMaj(\rho,l)=\prod_{v<f}\min\Bigl(1,\bigl(\rho_v/|l|_v\bigr)^K\Bigr),\qquad
\kappa^{(s,t,t')}(h,h')=\kappa\bigl(h\cdot(s,t),h'\cdot(s,t')\bigr).
\end{gather*}
\formal{Formalism.FourierDecayAlt.TB.Data}
\end{definition}

The kernels $\mathcal K_i$, $\Theta_{i,s,t}$ and the numbers $e_i$ enter the abstract recursion of \cref{sec:pf:thmS}
but none of the following conditions.

\begin{definition}[Real inputs]\label{pf:thmS-carrier:realinputs}
Let $\thmScarrierDat$ be an energy datum, $C_4,j\in\R$ and $V\colon\thmScarrierPt(E)\to H$ a map to a Hilbert space.
$\thmScarrierDat$ satisfies the \emph{real inputs} with constants $(C_4,j)$ and top vectors $V$ if the following eight
conditions hold.
\begin{enumerate}[label=(R\arabic*)]
\item\label{pf:thmS-carrier:R1} (Gram positivity) For $\eps>0$, all $i,w$ and all finite $Y$, the kernel
  $\gamma_{\eps,i,w}$ is PSD on $Y$.
\item\label{pf:thmS-carrier:R2} (Words) $\thmScarrierWt(w)>0$ for all $w\in W$.
\item\label{pf:thmS-carrier:R3} (Wave estimate) For $\eps>0$, $1\le i\le r+1$ and $w\in W_{\mathrm{all}}$ with
  $\thmScarrierCnt(w)+K\le\thmScarrierNtop$ there are probability vectors $\theta_A$ on
  $\thmScarrierWb(\thmScarrierCnt(w)+K)$, one for each set $A$ of places, such that for all finite $Y$ of points of
  level $|h|=i-1$ and all PSD kernels $\kappa$ on $Y$:
  $\thmScarrierNhat_{\eps,i,w}(Y,\kappa;0)=\thmScarrierNz_{\eps,i-1,w}(Y,\kappa)$, and for every
  $l\in\Lambda^\vee_i\smallsetminus\{0\}$ there is $A$ with
  \[
  \bigl|\thmScarrierNhat_{\eps,i,w}(Y,\kappa;l)\bigr|\le\thmScarrierMaj(\rho^D_i,l)\sum_{w'}\theta_A(w')\,
  \thmScarrierNz_{\eps,i-1,w'}(Y,\kappa).
  \]
\item\label{pf:thmS-carrier:R4} (Poisson input) There are a finite-dimensional real inner product space $V_\R$, a finite
  set $I$ and $\R$-linear maps $\pi_\iota\colon V_\R\to\C$ such that $c\mapsto(\pi_\iota(c))_\iota$ is bijective, full
  lattices $\Lambda_i\subset V_\R$, additive maps $\mathrm{emb}_i\colon E\to V_\R$ and maps $\mathrm{dm}_i\colon E\to
  V_\R$ with:
  (\emph{frames}) for $1\le i\le r+1$, $\mathrm{emb}_i$ is injective, $\mathrm{emb}_i(\Lambda_{i,1})\subset\Lambda_i$,
  $\mathrm{dm}_i(0)=0$, and every element of the dual lattice $\Lambda_i^*$ is $\mathrm{dm}_i(l)$ for some
  $l\in\Lambda^\vee_i$;
  (\emph{wave functions}) for $\eps>0$, $1\le i\le r+1$, $w$, $Y$ of level $i-1$ and PSD $\kappa$ there are a continuous
  $\varphi\colon V_\R\to[0,\infty)$ and $C_\varphi$ with $\varphi(c)\le C_\varphi\prod_\iota(1+|\pi_\iota c|)^{-3}$,
  $\varphi(\mathrm{emb}_i t)=\Re\thmScarrierNst_{\eps,i,1,w}(Y,\kappa;t,t)$ for $t\ne0$, and
  $\thmScarrierNhat_{\eps,i,w}(Y,\kappa;l)=\covol(\Lambda_i)^{-1}\widehat\varphi(\mathrm{dm}_i l)$ for all $l$, where
  $\widehat\varphi$ is the Fourier transform.
\item\label{pf:thmS-carrier:R5} (Lift estimate) For $\eps>0$, $1\le s\le i\le r+1$ with $(i,s)\ne(r+1,1)$ and
  $w\in W_{\mathrm{all}}$ with $\thmScarrierCnt(w)+K\le\thmScarrierNtop$ there are probability vectors $\theta_\pi$ on
  $\thmScarrierWb(\thmScarrierCnt(w)+K)^2$, one for each $\pi\colon\{v<f\}\to\{0,1,2\}$, such that for all finite $Y$ of
  level $i-s$, finite $S\subset E$, kernels $\kappa$ PSD on $\{h\cdot(s,t):h\in Y,t\in S\}$ and $t,t'\in
  S\smallsetminus\{0\}$ with $t^{-1}\ne t'^{-1}$ there is $\pi$ with
  \begin{multline*}
  \bigl|\thmScarrierNst_{\eps,i,s,w}(Y,\kappa^{(s,t,t')};t,t')\bigr|\le
  \thmScarrierMaj\bigl(\rho^L_{i-s+1,s},t^{-1}-t'^{-1}\bigr)\sum_{(x_1,x_2)}\theta_\pi(x_1,x_2)\\
  \times\frac12\Bigl(\Re\thmScarrierNst_{\eps,i,s,x_1}(Y,\kappa^{(s,t,t)};t,t)
  +\Re\thmScarrierNst_{\eps,i,s,x_2}(Y,\kappa^{(s,t',t')};t',t')\Bigr).
  \end{multline*}
\item\label{pf:thmS-carrier:R6} (Block estimate) For $\eps>0$, $2\le s\le i\le r+1$ and $w\in W_{\mathrm{all}}$ with
  $\thmScarrierCnt(w)+(s+1)K\le\thmScarrierNtop$ and $V_{i-s+1,s}\le1$ there is a probability vector $\theta$ on
  $\thmScarrierWb(\thmScarrierCnt(w)+(s+1)K)$ such that for all finite $Y$ of level $i-s$, PSD $\kappa$ on $Y$, $t_0\in
  E$ and finite $S\subset E\smallsetminus\{0\}$ with $S\subset t_0+\Lambda_{i-s+1,s}$,
  \[
  \sum_{t\in S}\Re\thmScarrierNst_{\eps,i,s,w}(Y,\kappa;t,t)\le C_4^{fs}\prod_{1\le s'<s}V_{i-s+1,s'}\sum_{w'}
  \theta(w')\thmScarrierNz_{\eps,i-s,w'}(Y,\kappa).
  \]
\item\label{pf:thmS-carrier:R7} (Energy) For every $w\in\thmScarrierWb(\thmScarrierNtop)$,
  $\thmScarrierWt(w)^{-1}\Re\gamma_{\eps,0,w}(\emptyset,\emptyset)\to j$ as $\eps\to0^+$.
\item\label{pf:thmS-carrier:R8} (Cut-off) For every finite set $Y$ of points of level $r+1$, for all sufficiently small
  $\eps>0$: $\gamma_{\eps,r+1,\mathrm{emp}}(p,q)=\langle V(q),V(p)\rangle$ for all $p,q\in Y$ with
  $\mathcal K_{r+1}(p,q)\ne0$.
\end{enumerate}
\formal{Formalism.FourierDecayAlt.PB.RealInputs}
\formal{Formalism.FourierDecayAlt.J.Frames}
\formal{Formalism.FourierDecayAlt.J.WaveFn}
\end{definition}

\subsubsection{Hida settings}
\begin{definition}[Archimedean coordinates; real place data]\label{pf:thmS-carrier:realplace}
Let $E$ be a CM field. The \emph{archimedean coordinate space} is $E_\infty=\prod_x\C$, the product over the infinite
places $x$ of $E$, with the Euclidean structure of $\C^{\{x\}}$. For integers $k\ge1$, $r\ge0$ let
$\mathrm{Mat}_{k,r}=M_{2r+k+1}(\C)$ and let $G_{k,r}=\Uni(H_{k,r})$ be the unitary group of a fixed Hermitian form
$H_{k,r}$ on $\C^{2r+k+1}$. A \emph{real place datum} consists of vectors $z,z'\in\C^{2r+k+1}$, unitary representations
$\pi,\pi'$ on finite-dimensional spaces $W_p,W_q$ of the stabilisers of $z,z'$, Haar measures, a section and a cut-off
$\chi$. It defines the \emph{Bessel transform} $F\mapsto J_F$ of functions $F\colon\mathrm{Mat}_{k,r}\to\Hom(W_q,W_p)$
(an integral of translates of $F$ over the stabilisers of $z,z'$ and two unipotent subgroups, against a character of the
latter), a real function $G_f(z,z';g)$ on $G_{k,r}$, and the norm $\|J_F\|^2\in[0,\infty]$.
\formal{Formalism.FourierDecayAlt.J.CV}
\formal{Formalism.QuantumDeterminants.Hpk}
\formal{Formalism.FourierDecayAlt.B.RealPlaceData}
\formal{Formalism.FourierDecayAlt.B.JF}
\formal{Formalism.QuantumDeterminants.Gf}
\formal{Formalism.FourierDecayAlt.B.normSq}
\end{definition}

\begin{definition}[Hida settings]\label{pf:thmS-carrier:hida}
A \emph{Hida setting} $S$ over $E$ with local spaces $W_p(x),W_q(x)$ consists of: the number $f$ of complex places with an
enumeration $v\mapsto x_v$ and embeddings $\sigma_v\colon E\to\C$ inducing $x_v$; \emph{global size data} $G$, which
contain integers $k,r,g_1,g_2$ and, for each place, a function $P_v\colon\R\to\R$ and cut parameters; numbers $C_A,T,T'$,
functions $d_e,d_e'\colon\N\to\R$ and an integer $l$; finite \emph{grids} $\mathrm{grid}_v(a),\mathrm{grid}^1_v\subset\R$;
and, for each $v$, \emph{place data} consisting of a real place datum $\mathrm{bs}_v$, a weight $G_v$ on local points and
\emph{local kernels} $RR_v(\eps,i,a,p,q)\in\C$, where $a$ is a \emph{local word} (a finite list of letters
$(b,Y)\in\{0,1\}\times\R$) and $p,q$ are local points (finite sequences in $\N\times\C$). Put $g=g_1g_2$,
$\thmScarrierNtop=20(r+1)$, and let $[\alpha_v,\beta_v]$ be the cut interval at $v$.

The \emph{local Gram kernels} are $\gamma^{v}_{\eps,i,a}(p,q)=\sqrt{n_v(p)n_v(q)}\,RR_v(\eps,i,a,p,q)$, with explicit
positive normalising factors $n_v$. A \emph{word} is a family $w=(w_v)_v$ of local words; its \emph{count} is
$\max_v\len(w_v)$ and its \emph{weight} is $\thmScarrierWt(w)=\prod_v\prod_{(b,Y)\in w_v}|P_v(Y)|$. A word is
\emph{admitted} if, for every $v$, each letter $(b,Y)$ of $w_v$ in position $j$ satisfies
$Y\in\mathrm{grid}_v(\len(w_v)-j)\cup\mathrm{grid}_v^1$ (if $b=1$), resp.\
$Y-1\in\mathrm{grid}_v(\len(w_v)-j)\cup\mathrm{grid}_v^1$ (if $b=0$), and distinct letters $(b,Y),(b',Y')$ of $w_v$
satisfy $|Y-Y'+\delta|>2$ for all integers $|\delta|\le2$. Let $W(S)$ be the set of admitted words of count
$\le\thmScarrierNtop$. For a point $p$ of $E$ let $\sigma_vp$ be the local point obtained by applying $\sigma_v$ to the
second coordinates.
\formal{Formalism.FourierDecayAlt.C.HidaSetting}
\formal{Formalism.FourierDecayAlt.Sz.Glob}
\formal{Formalism.FourierDecayAlt.C.nLoc}
\end{definition}

\begin{definition}[The energy datum of a Hida setting]\label{pf:thmS-carrier:mkdata}
Let $S$ be a Hida setting. For $\eps$, $i$, $w\in W(S)$ and points $p,q$ put
\[
\gamma_{\eps,i,w}(p,q)=\prod_{v}\gamma^{v}_{\eps,i,w_v}(\sigma_vp,\sigma_vq)\quad(i\le r+1),\qquad
\gamma_{\eps,i,w}=0\quad(i>r+1).
\]
For $c\in E_\infty$ let $N^1_{\eps,i,w}(Y,\kappa;c)=\thmScarrierWt(w)^{-1}\sum_{h,h'\in Y}\kappa(h,h')\prod_v
\gamma^{v}_{\eps,i,w_v}(\sigma_vh\cdot(1,c_v),\sigma_vh'\cdot(1,c_v))$ and
\begin{equation}\label{eq:pf:thmS-carrier:Nhat}
\thmScarrierNhat_{\eps,i,w}(Y,\kappa;l)=\covol(\Lambda_i)^{-1}\int_{E_\infty}N^1_{\eps,i,w}(Y,\kappa;c)\,
\exp\Bigl(\frac{2\pi i}{g}\sum_v2\Re\bigl(\sigma_v(l)c_v\bigr)\Bigr)\,dc,
\end{equation}
where $\Lambda_i\subset E_\infty$ is the image of the period lattice $L_i=\Lambda_{i,1}$ under the diagonal embedding.
Given kernels $\mathcal K_i$ with $\mathcal K_0(\emptyset,\emptyset)=1$, theta kernels $\Theta$ and numbers $e_i$, the
energy datum $\thmScarrierDat_S(\mathcal K,\Theta,e)$ has $r=G.r$, $f$, $K=5$, $|l|_v=$ the absolute value of $l$ at
$x_v$, word set $W(S)$ with the above weights and counts, $W_{\mathrm{all}}$ the words of $W(S)$ all of whose letters
$(b,Y)$ have $Y$ in the (finite) set of grid points and their translates by $1$, $\mathrm{emp}$ the empty word, the Gram
kernels and Fourier coefficients above, the given $\mathcal K,\Theta,e$, explicit lattices
$\Lambda_{b,s},\Lambda'_{b,s},\Lambda^\vee_b$ attached to $(g,r)$, and sizes $\rho^D,\rho^L,V$ given by the thresholds
of the global size data. The \emph{base datum} $S_{\mathrm{base}}$ is $\thmScarrierDat_S$ with trivial kernels.
\formal{Formalism.FourierDecayAlt.J.eLat}
\formal{Formalism.FourierDecayAlt.J.LaB}
\formal{Formalism.FourierDecayAlt.C.HidaSetting.mkData}
\end{definition}

\begin{definition}[Local-to-global conditions]\label{pf:thmS-carrier:conds}
Let $S$ be a Hida setting, $C_4\in\R$, $j=(j_v)_v$ real numbers and $V\colon\thmScarrierPt(E)\to H$. We consider the
following conditions.
\begin{description}
\item[(H1) local positivity] for all $v$, $\eps>0$, $i\le r+1$, $w\in W(S)$ and finite $Y$, the local Gram kernel
  $\gamma^v_{\eps,i,w_v}$ is PSD on $Y$;
\item[(H2) grid condition] the grids are well spaced inside the cut intervals, and $P_v(y)\ne0\ne P_v(y+1)$ for every
  grid point $y$;
\item[(H3) integral formula] for $\eps>0$, $1\le i\le r+1$, $w$, $Y$ of level $i-1$ and PSD $\kappa$, the function
  $c\mapsto N^1_{\eps,i,w}(Y,\kappa;c)$ on $\{c:c_v\ne0\ \forall v\}$ extends to a continuous function $\varphi$ on
  $E_\infty$ with $|\varphi(c)|\le C_\varphi\prod_v(1+|c_v|)^{-3}$, and
  $\thmScarrierNhat_{\eps,i,w}(Y,\kappa;0)=\thmScarrierNz_{\eps,i-1,w}(Y,\kappa)$;
\item[(H4) plane-wave bound] the bound of \ref{pf:thmS-carrier:R3} for $S_{\mathrm{base}}$, with $A$ the set
  $\mathrm{big}_i(l)$ of places $v$ with $\rho^D_{i,v}<|l|_v$;
\item[(H5) lift bound] condition \ref{pf:thmS-carrier:R5} for $S_{\mathrm{base}}$;
\item[(H6) block bound with constant $C_4$] condition \ref{pf:thmS-carrier:R6} for $S_{\mathrm{base}}$;
\item[(H7) energy limit with local energies $j$] for $w\in W(S)$ and every $v$,
  $RR_v(\eps,0,w_v,\emptyset,\emptyset)\to\mathrm{size}_v(w_v)\,j_v$ as $\eps\to0^+$, where
  $\mathrm{size}_v(a)=\prod_{(b,Y)\in a}|P_v(Y)|$;
\item[(H8) top Gram identity for $V$] for every finite set $Y$ of points of level $r+1$, for all small $\eps>0$:
  $\gamma_{\eps,r+1,\mathrm{emp}}(p,q)=\langle V(q),V(p)\rangle$ for $p,q\in Y$.
\end{description}
\formal{Formalism.FourierDecayAlt.C.ThmPPos}
\formal{Formalism.FourierDecayAlt.C.GridsOK}
\formal{Formalism.FourierDecayAlt.C.ThmPInt}
\formal{Formalism.FourierDecayAlt.C.PlaneWaveBoundAllPlaces}
\formal{Formalism.FourierDecayAlt.C.OffDiagonalLiftBound}
\formal{Formalism.FourierDecayAlt.C.CorollaryT}
\formal{Formalism.FourierDecayAlt.C.LemmaE}
\formal{Formalism.FourierDecayAlt.C.PropGram}
\end{definition}

\subsubsection{Built places, leading-form settings and model places}\label{sec:pf:thmS-carrier:built-defs}
\begin{definition}[Built places and leading-form settings]\label{pf:thmS-carrier:builtdef}
A Hida setting $S$ is \emph{built} from data $(T,f,\chi,\nu',m_0,\mathrm{top},A_0)$ if at each place $v$ the local kernels
are those of the \emph{component family} attached to a level tower $T_v$ (a sequence of real place data for the ranks
$r,r-1,\dots,0$), a smooth test function $f_v\colon\mathrm{Mat}_{k,r}\to\Hom(W_q,W_p)$ all of whose derivatives satisfy
the growth bounds with constant $A_0$, a cut-off $\chi$, a base frame $(\nu'_v,m_{0,v})$ and a top kernel, normalised by
$g$ and the root discriminant of $E$, and if moreover $f_v$ is an eigenfunction of the Casimir-type operators,
$\mathrm{CV}(Y)f_v=P_v(Y)f_v$ on $G_{k,r}$. A \emph{built place with constant $C_c$} is a built Hida setting with
$g_1,g_2>0$, $E$ unramified over its maximal totally real subfield $F$ at all finite primes, $1\le l\le k-2$, the grid
condition, the long-cut condition
$C_c(k\thmScarrierNtop(r+1)^2+\thmScarrierNtop^2(r+1)^2+\thmScarrierNtop^3)\le\beta_v-\alpha_v$,
$C_c\le C_A$, smooth $f_v$, a lower bound for $A_0$ in terms of $C_c$, and the condition that the whole cut interval lies
on the cut of the parameter. A \emph{leading-form setting with constants $(C_b,C_c)$} is the analogous condition in
which the local data come from \emph{place packages} $p_v$ (a tower, a base frame, a top lift $\Theta_v$ of class
$D^+$ with constant $C_b$, and a type projector), and the grids are moreover \emph{rich}; a \emph{leading-form place}
satisfies in addition the eigenfunction and cut-inside conditions.
\formal{Formalism.FourierDecayAlt.D.LevTower}
\formal{Formalism.FourierDecayAlt.D.IsBuilt}
\formal{Formalism.FourierDecayAlt.JCT.IsBuiltPlace}
\formal{Formalism.FourierDecayAlt.KB.Pl}
\formal{Formalism.FourierDecayAlt.KB.IsLeadingFormSetting}
\formal{Formalism.FourierDecayAlt.KB.IsLeadingFormPlace}
\end{definition}

\begin{definition}[Model places]\label{pf:thmS-carrier:modelplace-def}
Let $S$ be built from $(T,f,\chi,\dots)$ and $v$ a place. Model data at $v$ are a level datum $L$ of the model operator
algebra, a family $j=(j_\rho)$ of smooth functions on the domain $Q_{r,T_T}=\{(x,z)\in\C^r\times\C:x_l\ne0,\
|z|^2>T_T\}$ indexed by types $\rho$, real numbers $\lambda_1,\dots,\lambda_{k+2r+1}$, a cut $c$ and constants
$C_m,\sigma$. The place $v$ is a \emph{model place} for these data if $c$ has cut-off function $\chi$, $C_m,\sigma>0$, the
model Hida conditions hold on $[\alpha_v,\beta_v]$, $j$ satisfies the decay, adjointness and
eigenfunction conditions of the model (the first two to order $2\thmScarrierNtop+2$), and for all $\eps>0$ and all
admissible local words $a$
\begin{equation}\label{eq:pf:thmS-carrier:modelenergy}
\mathcal E_v(a;\mathrm{cut}_\eps f_v)=C_m\,\thmScarrierEn^{\mathrm{mod}}_v(a;\sigma\eps),
\end{equation}
where the left side is the principal Hermitian lattice energy of the cut test function and the right side is the cut-off
energy form of the model (\cref{pf:thmS-carrier:modelsys}); finally the model system admits every admissible $a$, with
$\mathrm{size}(a)=\mathrm{size}_v(a)$. The \emph{model energy} at $v$ is
\[
j_v^{\mathrm{mod}}=C_m\int_{Q_{r,T_T}}\sum_\rho|j_\rho(q)|^2\,d\mu_0(q)\Big/\Gamma_0,
\]
with $\mu_0$ the measure of density $\prod_l|x_l|^{-2}$ and $\Gamma_0$ an explicit normalising constant.
\formal{Formalism.FourierDecayAlt.C7.principalHermitianLattice}
\formal{Formalism.FourierDecayAlt.JC.IsModelPlace}
\formal{Formalism.FourierDecayAlt.JCT.jvModel}
\end{definition}

\subsubsection{The carrier statement}
We use the global data of Theorem~S as in \cref{sec:pf:thmS}: a datum $D$ (ranks, host, spaces), size data $L$, spectral
ties $T$ with Hida setting $X=T.X$ over the totally real subfield $F$ of the host field $E$, words $u,u'$, a bijection $e$
between the infinite places of $E$ and the index set $\mathrm{Tau}$, and the derived global size data
$G_0=G(D,e)$, words $d_e=d(D,\tau,u)$, $d_e'=d(D,\tau,u')$ and level
$l=l(d_e,d_e')$.

\begin{definition}[Carriers]\label{pf:thmS-carrier:carrierfull}
Let $V_{\mathrm{pt}}(p)\in\mathrm{HS}(X.V_\pi(u'),X.V_\pi(u))$ be the scaled operators at point representatives and
$j_X$ the energy constant of the mass data. The \emph{carrier bound} with constants $(C_4,C_A,G_0,d_e,d_e',l)$
holds if there are local spaces $W_p(x),W_q(x)$ ($W_q(x)$ finite-dimensional), a Hida setting $S$ over $E$, level towers
$T_v$, test functions $f_v$, base frames $(\nu'_v,m_{0,v})$, a bijection $\tau$ between the places of $S$ and the
infinite places of $F$, matrices $Q_y,Q^{\mathrm i}_y$ and isometries $W_p(x_{\tau^{-1}y})\cong X.V_{\mathrm{loc}}(u,y)$,
$W_q(x_{\tau^{-1}y})\cong X.V_{\mathrm{loc}}(u',y)$ such that: $S$ has constants $C_A,G_0,d_e,d_e',l$; the ranks, levels,
primes $T,T'$ and number of places of $S$ are those of $D$; the base frames are those of $X$; the top Gram identity (H8)
holds for $V_{\mathrm{pt}}$, together with the compatibility of the embeddings and of the top weights $G_v$ with $X$; the
operators at point representatives factor as tensor products of the local vectors attached to $(T_v,\nu'_v,m_{0,v},f_v)$;
and the energy datum $\thmScarrierDat^X_S=\thmScarrierDat_S(\mathcal K^{X},\Theta^X,e^X)$, formed with the kernels attached to $X$, satisfies the real inputs with
constants $(C_4,j_X)$ and top vectors $V_{\mathrm{pt}}$.
\formal{Formalism.L2.SG.SpectralTiesReal.gOf}
\formal{Formalism.L2.SG.RealObj.VtPt}
\formal{Formalism.L2.SG.SpectralTiesReal.massDataY}
\formal{Formalism.L2.SG.SpectralTiesReal.CarrierFull}
\end{definition}

\begin{theorem}[The carrier input]\label{pf:thmS-carrier:main}
There are $C_4\ge1$ and $C_{A,0}$ with the following property. For all $C_A\ge C_{A,0}$ with $C_A\ge1$ and all $C_0\ge1$,
$C_G\ge0$, $C_\Delta>4$, $c_T>0$, $C_b>0$ with $c_TC_b>1$ there is $k_0\in\N$ such that for all $C_g$, $c_b>0$, $C_w$,
$c_{Tf}>0$, $C_T$, $C_{0f}$ there are $C'_0,C_{l,0}$ such that for all $C'\ge C'_0$ and $C_l\ge C_{l,0}$: for every datum
$D$, size data $L$, spectral ties $T$, words $u\ne u'$, bijection $e$, index $\tau$, numbers $C_p,\Theta'$, index choices
$i_0=i_0'$ at the infinite places of $F$, measures $\nu_q$ and auxiliary mass data $(X_0,s_d,C,\Theta,F_0,v_R,Q)$, if
\begin{itemize}
\item $D$ is tied, $k_0\le k(G_0)$, and the count hypotheses hold for $D$, $L$ and the constants $C'$, $C_l$, $C_g$, $c_b$,
  $C_w$, $c_{Tf}$, $C_T$, $C_{0f}$;
\item the embeddings of $X$ induce the places of $E$ through $e$;
\item the size hypotheses hold for $G_0$ with the constants $C_A$, $C_0$, $C_G$, $C_\Delta$, $c_T$, $C_b$, the scale
  $b=L.\mathrm{sc}.b$, the numbers $C_p,\Theta'$, the words $d_e,d_e'$, the primes $T_u,T_{u'}$, the dimensions
  $d_u,d_{u'}$ and the level $l$;
\item point representatives exist, and the measures $\nu_q$ are pinned,
\end{itemize}
then the carrier bound with constants $(C_4,C_A,G_0,d_e,d_e',l)$ holds for these data. That is, $\sfname{CarrierP}$ holds.
\formal{Formalism.FourierDecayAlt.J.SzHyp}
\formal{Formalism.L2.SG.RealObj.PtRepExists}
\formal{Formalism.L2.SG.RealObj.NuQPin}
\formal{Formalism.L2.SG.PoincareAmplitudeLowerBounds.CarrierP}
\formallabel{GlobalTheory.CarrierGap.carrierP}
\end{theorem}

\begin{lemma}[The energy constant]\label{pf:thmS-carrier:intN}
The energy constant $j_X$ of the mass data equals an explicit expression in $X$, the words $u,u'$ and the additive
character $\psi$; in particular it does not depend
on the auxiliary data $X_0,s_d,C,\Theta,v_R,Q$.
\formallabel{Formalism.L2.SG.Pf.Carrier.EnergyPf.intN_massDataY}
\end{lemma}

\begin{proof}
Both sides unfold to the same expression; the formal proof is by definitional equality.
\end{proof}

\begin{remark}\label{pf:thmS-carrier:rem-statement}
(1) None of the eight real inputs involves the kernels $\mathcal K_i$, $\Theta_{i,s,t}$ or the numbers $e_i$ of the energy
datum; accordingly they are proved for $\thmScarrierDat_S(\mathcal K,\Theta,e)$ with arbitrary $(\mathcal K,\Theta,e)$, and
the particular kernels of $\thmScarrierDat^X_S$ play no role in this section. (2) By \cref{pf:thmS-carrier:intN} the
auxiliary data enter the conclusion only through a quantity that does not depend on them. (3) The index choices are
assumed equal, $i_0=i'_0$.
\end{remark}

% ---------------------------------------------------------------------------
\subsection{Proof of the carrier input}\label{sec:pf:thmS-carrier:assembly}

The proof combines five inputs: three archimedean statements, the model lemma, and a bound on the global side. We state
them and deduce \cref{pf:thmS-carrier:main}.

\begin{proposition}[Archimedean inputs]\label{pf:thmS-carrier:needs}
The following three statements hold.
\begin{enumerate}[label=(I\arabic*)]
\item\label{pf:thmS-carrier:I1} (Theorem~P for compactly supported test functions) For every parameter $P$, real place
  datum with vectors $z,z'$ and representations $\pi,\pi'$, and actual-operator datum $(X,\mu,L,\mathrm{st})$ at a level
  $1\le l_0\le k$, there are submodules $K_0,\dots,K_r$ of the radial space for which Theorem~P, the slice inclusions and
  the composite level-zero formula hold, and such that for every smooth compactly supported $F$ satisfying all growth
  bounds and supported in $\{G_f\ge\eps\}$ for some $\eps>0$, the restriction of the Bessel transform $J_F$ to the radial
  domain lies in $K_0$.
\item\label{pf:thmS-carrier:I2} (Model places) There is $C_c$ such that for every Hida setting $S$ which is a built place
  with constant $C_c$, whose place data are admissible and whose words $d_e,d_e'$ are admissible, every cut $c$ with
  cut-off function $\chi$, and every place $v$, there are model data at $v$ for which $v$ is a model place
  (\cref{pf:thmS-carrier:modelplace-def}).
\item\label{pf:thmS-carrier:I3} (Built bounds) Assuming \ref{pf:thmS-carrier:I1}: for every $C_b$ there are $C_c$ and
  $C_4\ge1$ such that every leading-form setting with constants $(C_b,C_c)$ satisfies local positivity (H1), the local
  integral formulas at all places, the place-wise wave bound together with the curvature barriers, the lift bounds below
  the top and on the top diagonal together with their thresholds, and the coset-sum bound with constant $C_4$ (see
  \cref{pf:thmS-carrier:planewave,pf:thmS-carrier:liftplaces,pf:thmS-carrier:corT,pf:thmS-carrier:intplaces}).
\end{enumerate}
\formal{Formalism.FourierDecayAlt.P5.ThmPActualCpt}
\formal{Formalism.FourierDecayAlt.JCT.ModelPlaceBuilt2}
\formal{Formalism.FourierDecayAlt.KB.BuiltBounds2Cpt}
\formallabel{GlobalTheory.carrierNeedsFromD}
\end{proposition}

\begin{proof}
\ref{pf:thmS-carrier:I1} is \lref{Formalism.FourierDecayAlt.P.CutSE.actual_compact_theorem}, which proves exactly this
statement by a compactness argument for the spectral operators. \ref{pf:thmS-carrier:I2} is \cref{pf:thmS-carrier:modelplace}
and \ref{pf:thmS-carrier:I3} is \cref{pf:thmS-carrier:builtbounds} below.
\end{proof}

\begin{lemma}[Top lifts]\label{pf:thmS-carrier:liftat3}
There is $C_b>0$ with the following property (TL$_{C_b}$): for every parameter $P$, real place datum $B_s$ and
actual-operator datum $(X,\mu,L,\mathrm{st})$ at a level $1\le l_0\le k$ there is a top lift $\Theta$ of the datum which
lies in the class $D^+$ with constant $C_b$.
\formal{Formalism.L2.SG.Pf.CarrierXAll.LiftAt3}
\formallabel{Formalism.L2.SG.Pf.CarrierXAll.liftAt3_exists}
\end{lemma}

\begin{proof}
This is \lref{Formalism.FourierDecayAlt.P.top_lift_exists}, which constructs such top lifts, with $C_b=11$, from the
lifts of functions, of the actual operators and of the Capelli operators by a recursion over words.
\end{proof}

\begin{definition}[The global-side bound]\label{pf:thmS-carrier:gsdef}
Let $C_6,C_c,C_b\in\R$. The \emph{carrier data bound} with constants $(C_6,C_c,C_b)$ holds for global data as in
\cref{pf:thmS-carrier:main} if there are finite-dimensional local spaces, a Hida setting $S$ with the constants
$C_A,G_0,d_e,d_e',l$ and the ranks, levels, primes and number of places of $D$, place packages $p_v$, test functions $f_v$, a
cut-off $\chi$, a number $A_0$, a cut $c$ with cut-off function $\chi$, and $\tau,Q,Q^{\mathrm i}$ and isometries as in
\cref{pf:thmS-carrier:carrierfull}, such that
\begin{enumerate}[label=(\roman*)]
\item $S$ is a leading-form place with constants $(C_b,C_c)$ for $(p,f,\chi,A_0)$, the slice budget for $A_0$ holds, and
  every cut interval has length at least $C_6\bigl(k\thmScarrierNtop(r+1)^2+\thmScarrierNtop^2(r+1)^2\bigr)$;
\item the top vectors factor: $\langle V_{\mathrm{pt}}(q),V_{\mathrm{pt}}(p)\rangle=\prod_v\sqrt{n_v(p)n_v(q)}\,
  \langle\mathrm{vec}_v(\sigma_vq),\mathrm{vec}_v(\sigma_vp)\rangle$ for points $p,q$ of level $r+1$, where
  $\mathrm{vec}_v$ are the local vectors of the place packages;
\item the towers, base frames and top weights of $S$ are built from those of $X$, and the operators at point
  representatives factor through the local vectors, as in \cref{pf:thmS-carrier:carrierfull};
\item for all model data $(L,j,\lambda,C_m,\sigma)$ for which every place is a model place, $\prod_vj^{\mathrm{mod}}_v
  =j_X$.
\end{enumerate}
The \emph{global-side bound} (GS$_{C_b}$) is the statement obtained from \cref{pf:thmS-carrier:main} by replacing ``there are
$C_4\ge1$ and $C_{A,0}$'' with ``for all $C_6>0$ and $C_c$ there is $C_{A,0}$'' and the carrier bound with the carrier data
bound with constants $(C_6,C_c,C_b)$.
\formal{Formalism.L2.SG.Pf.CarrierT.CarrierDataBound}
\formal{Formalism.L2.SG.SpectralTiesReal.CarrierSideBoundPositive}
\end{definition}

\begin{proposition}[The compact-model theorem]\label{pf:thmS-carrier:compactmodel}
Assume \ref{pf:thmS-carrier:I1}--\ref{pf:thmS-carrier:I3}, the model lemma (\cref{pf:thmS-carrier:lemmamodel}) and the
global-side bound (GS$_{C_b}$) for some $C_b$. Then $\sfname{CarrierP}$ holds.
\formallabel{Formalism.L2.SG.SpectralTiesReal.carrierP_of_compact_model}
\end{proposition}

The proof is given in \cref{sec:pf:thmS-carrier:compactproof}; the global-side bound is
\cref{pf:thmS-carrier:xside-proof} in \cref{sec:pf:thmS-carrier:X}.

\begin{proof}[Proof of \cref{pf:thmS-carrier:main}]
By \cref{pf:thmS-carrier:needs} the statements \ref{pf:thmS-carrier:I1}--\ref{pf:thmS-carrier:I3} hold. By
\cref{pf:thmS-carrier:liftat3} there is $C_b$ with (TL$_{C_b}$). The uniqueness of Bessel models with large gap holds by
\lref{CompanionTheory.Lime.BM.uniqGap_holds}: for a discrete series representation of $\Uni(H_{k,r})$ whose
Harish-Chandra parameter has gap at least $120(2r+k+1)^3$, any two smooth Bessel maps with values in a representation of
the stabiliser are linearly dependent on smooth vectors. The existence of $K$-finite vectors with non-zero Bessel
functional holds by \lref{Formalism.L2.SG.Pf.CarrierXAll.kFinNZ_of}, which deduces it from the density of $K$-finite
vectors \lref{Formalism.L2.SG.Pf.CarrierXAll.kFinDense_holds}. By \cref{pf:thmS-carrier:xpins,pf:thmS-carrier:xleaves}
the pinning and leaf statements follow, and \cref{pf:thmS-carrier:xside-proof} gives (GS$_{C_b}$). The model lemma holds by
\cref{pf:thmS-carrier:lemmamodel}. Now \cref{pf:thmS-carrier:compactmodel} gives $\sfname{CarrierP}$.
\end{proof}

% ---------------------------------------------------------------------------
\subsection{The real inputs from local statements}\label{sec:pf:thmS-carrier:parts}

\subsubsection{Reduction to the Hida setting}
\begin{proposition}[Real inputs of a Hida setting]\label{pf:thmS-carrier:realconcrete}
Let $S$ be a Hida setting with $g_1,g_2>0$, $j=(j_v)_v$ real numbers, $V\colon\thmScarrierPt(E)\to H$ and $C_4\in\R$.
Assume (H1)--(H8) of \cref{pf:thmS-carrier:conds} (with $C_4$, $j$ and $V$). Then for all $(\mathcal K,\Theta,e)$ the energy
datum $\thmScarrierDat_S(\mathcal K,\Theta,e)$ satisfies the real inputs with constants $(C_4,\prod_vj_v)$ and top vectors $V$.
\formallabel{Formalism.FourierDecayAlt.P.real_inputs_concrete_holds}
\formallabel{Formalism.FourierDecayAlt.P.real_inputs_by_parts}
\end{proposition}

\begin{proof}
Conditions \ref{pf:thmS-carrier:R1}--\ref{pf:thmS-carrier:R8} are
\cref{pf:thmS-carrier:gram,pf:thmS-carrier:words,pf:thmS-carrier:wave,pf:thmS-carrier:poisson,pf:thmS-carrier:liftblock,pf:thmS-carrier:levelzero,pf:thmS-carrier:cutoff}
below, in this order, \cref{pf:thmS-carrier:liftblock} giving \ref{pf:thmS-carrier:R5} and \ref{pf:thmS-carrier:R6}.
\end{proof}

\begin{lemma}[Gram positivity]\label{pf:thmS-carrier:gram}
If (H1) and (H2) hold, then \ref{pf:thmS-carrier:R1} holds for every $\thmScarrierDat_S(\mathcal K,\Theta,e)$.
\formallabel{Formalism.FourierDecayAlt.P.gram_positive_semidefinite}
\end{lemma}

\begin{proof}
Let $\eps>0$, $w\in W(S)$ and $Y$ finite. If $i>r+1$ then $\gamma_{\eps,i,w}=0$, which is PSD. Let $i\le r+1$. For each place
$v$ the kernel $(p,q)\mapsto\gamma^v_{\eps,i,w_v}(\sigma_vp,\sigma_vq)$ is PSD on $Y$: for $x\in\C^Y$,
\[
\sum_{p,q\in Y}x_p\overline{x_q}\,\gamma^v(\sigma_vp,\sigma_vq)=\sum_{a,b\in\sigma_vY}y_a\overline{y_b}\,\gamma^v(a,b)\ge0,
\qquad y_a=\sum_{p\in Y,\ \sigma_vp=a}x_p,
\]
by (H1) on the finite set $\sigma_vY$ of local points; Hermitian symmetry is pulled back in the same way. The kernel
$\gamma_{\eps,i,w}$ is the entrywise (Hadamard) product of these $f$ kernels, and the empty product is the all-ones matrix
$\mathbf 1\mathbf 1^*$, which is PSD. By the Schur product theorem \cite[Section 7.5]{HornJohnson2013} and induction on the
number of factors, $\gamma_{\eps,i,w}$ is PSD on $Y$.
\end{proof}

The grid condition (H2) is not used in this proof.

\begin{lemma}[Positivity of the weights]\label{pf:thmS-carrier:words}
If (H2) holds, then $\thmScarrierWt(w)>0$ for every $w\in W(S)$; hence \ref{pf:thmS-carrier:R2} holds.
\formallabel{Formalism.FourierDecayAlt.P.words_pos_of_parts}
\end{lemma}

\begin{proof}
$\thmScarrierWt(w)$ is a product of the numbers $|P_v(Y)|$ over the letters $(b,Y)$ of the $w_v$. By admissibility, $Y$ (if
$b=1$) or $Y-1$ (if $b=0$) is a grid point $y$ at $v$, and (H2) gives $P_v(y)\ne0$ and $P_v(y+1)\ne0$. Hence every factor is
positive.
\end{proof}

\begin{lemma}[Wave estimate]\label{pf:thmS-carrier:wave}
If (H2), (H3) and (H4) hold, then \ref{pf:thmS-carrier:R3} holds for every $\thmScarrierDat_S(\mathcal K,\Theta,e)$.
\formallabel{Formalism.FourierDecayAlt.P.wave_estimate_of_integral}
\end{lemma}

\begin{proof}
Take the probability vectors $\theta_A$ of (H4). The identity $\thmScarrierNhat(Y,\kappa;0)=\thmScarrierNz_{\eps,i-1,w}(Y,
\kappa)$ is the second part of (H3), and for $l\ne0$ the bound holds with $A=\mathrm{big}_i(l)$. Neither side depends on
$(\mathcal K,\Theta,e)$.
\end{proof}

\begin{lemma}[Poisson input]\label{pf:thmS-carrier:poisson}
Assume the embedding facts, the Euclidean frame property, trace duality and the lattice facts recalled in the proof. If
$g_1,g_2>0$ and (H1)--(H3) hold, then \ref{pf:thmS-carrier:R4} holds for every $\thmScarrierDat_S(\mathcal K,\Theta,e)$.
\formallabel{Formalism.FourierDecayAlt.P.poisson_input_of_parts}
\formallabel{Formalism.FourierDecayAlt.P.Cx2.first_norm_real_nonneg}
\end{lemma}

\begin{proof}
Take $V_\R=E_\infty$ with the coordinate maps $\pi_x$, the diagonal embedding $\iota\colon E\to E_\infty$ as
$\mathrm{emb}_i$, the lattices $\Lambda_i=\iota(\Lambda_{i,1})$, and $\mathrm{dm}_i(l)=\mathrm{dm}(-l)$, where
$\mathrm{dm}(l)_x=\frac2g\,\overline{x(l)}$ for the embedding $x$. The auxiliary statements are: the embedding facts
\lref{Formalism.FourierDecayAlt.P.embedding_facts}, which state that $\iota$ is injective, that the coordinates form a
bijection $E_\infty\to\C^{\{x\}}$, and that $|\iota(t)_x|=|t|_x$ and $|\mathrm{dm}(l)_x|=\frac2g|l|_x$; the Euclidean frame
property \lref{Formalism.FourierDecayAlt.P.euclidean_frame} and the lattice facts
\lref{Formalism.FourierDecayAlt.P.lattice_facts_general}, which together show that $\Lambda_i$ is discrete of full rank
(with covolume equal to that of the corresponding lattice in the mixed space); and trace duality
\lref{Formalism.FourierDecayAlt.P.trace_dual_alt}, which states that for a full $\Z$-lattice $L\subset E$ an element
$\mu\in E_\infty$ lies in $\iota(L)^*$ if and only if $\mu=\mathrm{dm}(l)$ for some $l$ in the trace-dual lattice of $L$.
For $L=\Lambda_{i,1}$ the trace dual is $\Lambda^\vee_i$, so every $\mu\in\Lambda_i^*$ is $\mathrm{dm}_i(-l)$ with
$-l\in\Lambda^\vee_i$; injectivity, $\iota(\Lambda_{i,1})\subset\Lambda_i$ and $\mathrm{dm}_i(0)=0$ are clear. This gives
the frames.

For the wave functions fix $\eps,i,w,Y,\kappa$ and let $\varphi$ be the continuous extension of $c\mapsto
N^1_{\eps,i,w}(Y,\kappa;c)$ given by (H3), with its decay bound. For $c$ with all $c_x\ne0$, the kernel
$M_c(h,h')=\prod_v\gamma^v_{\eps,i,w_v}(\sigma_vh\cdot(1,c_v),\sigma_vh'\cdot(1,c_v))$ is PSD (as in
\cref{pf:thmS-carrier:gram}), so $\sum_{h,h'}\kappa(h,h')M_c(h,h')$ is real (both kernels are Hermitian) and non-negative
(it is the sum of the entries of the Hadamard product of two PSD matrices); with $\thmScarrierWt(w)>0$ we get
$N^1(c)\in[0,\infty)$. By continuity and density, $\varphi$ is real and non-negative everywhere. For $t\ne0$ all
$\sigma_v(t)\ne0$ and $\varphi(\iota t)=N^1(\iota t)=\thmScarrierNst_{\eps,i,1,w}(Y,\kappa;t,t)$. Finally
$\langle c,\mathrm{dm}(-l)\rangle=-\frac2g\sum_v\Re(\sigma_v(l)c_v)$, so the character in \eqref{eq:pf:thmS-carrier:Nhat} is
$e^{-2\pi i\langle c,\mathrm{dm}_i(l)\rangle}$; since $N^1=\varphi$ outside the null set $\{\exists x:c_x=0\}$,
\eqref{eq:pf:thmS-carrier:Nhat} reads $\thmScarrierNhat_{\eps,i,w}(Y,\kappa;l)=\covol(\Lambda_i)^{-1}\widehat\varphi(
\mathrm{dm}_i(l))$.
\end{proof}

\begin{lemma}[Lift and block estimates]\label{pf:thmS-carrier:liftblock}
(H5) implies \ref{pf:thmS-carrier:R5}, and (H6) with constant $C_4$ implies \ref{pf:thmS-carrier:R6}, for every
$\thmScarrierDat_S(\mathcal K,\Theta,e)$.
\formallabel{Formalism.FourierDecayAlt.P.lift_estimate_of_base}
\formallabel{Formalism.FourierDecayAlt.P.block_estimate_of_base}
\end{lemma}

\begin{proof}
The hypotheses are \ref{pf:thmS-carrier:R5} and \ref{pf:thmS-carrier:R6} for $S_{\mathrm{base}}$, and these conditions involve
only the Gram kernels, weights, counts, word sets, lattices and sizes, which do not depend on $(\mathcal K,\Theta,e)$.
\end{proof}

\begin{lemma}[Energy at level zero]\label{pf:thmS-carrier:levelzero}
If (H2) and (H7) hold, then \ref{pf:thmS-carrier:R7} holds for every $\thmScarrierDat_S(\mathcal K,\Theta,e)$, with the
constant $\prod_vj_v$.
\formallabel{Formalism.FourierDecayAlt.P.level_zero_estimate_product}
\end{lemma}

\begin{proof}
Let $w\in W(S)$. The normalising factors of the empty local point are $n_v(\emptyset)=1$ (empty product), so
$\gamma_{\eps,0,w}(\emptyset,\emptyset)=\prod_vRR_v(\eps,0,w_v,\emptyset,\emptyset)$, and
$\thmScarrierWt(w)=\prod_v\mathrm{size}_v(w_v)>0$ by \cref{pf:thmS-carrier:words}. By (H7) and continuity of finite products,
$\gamma_{\eps,0,w}(\emptyset,\emptyset)\to\prod_v\mathrm{size}_v(w_v)j_v=\thmScarrierWt(w)\prod_vj_v$ as $\eps\to0^+$. Taking
real parts and dividing by $\thmScarrierWt(w)$ gives \ref{pf:thmS-carrier:R7}.
\end{proof}

\begin{lemma}[Cut-off]\label{pf:thmS-carrier:cutoff}
(H8) for $V$ implies \ref{pf:thmS-carrier:R8} for every $\thmScarrierDat_S(\mathcal K,\Theta,e)$.
\formallabel{Formalism.FourierDecayAlt.P.cutoff_of_parts}
\end{lemma}

\begin{proof}
(H8) gives the identity for all $p,q\in Y$, without the condition $\mathcal K_{r+1}(p,q)\ne0$.
\end{proof}

\subsubsection{From places to the Hida setting}
Conditions (H1)--(H8) are deduced from local statements, one place at a time. The decay estimates use the following
elementary lemma.

\begin{lemma}[Big places]\label{pf:thmS-carrier:combine}
Let $K\in\N$, and for $v$ in a finite set let $a_v>0$ and $0\le\lambda_v\le x_v$. Let $A=\{v:x_v<a_v\}$, $M\ge0$ and $N\in\R$
with $\bigl(\prod_{v\in A}a_v^K\bigr)N\le\bigl(\prod_{v\in A}\lambda_v^K\bigr)M$. Then
$N\le\prod_v\min\bigl(1,(x_v/a_v)^K\bigr)\,M$.
\formallabel{Formalism.FourierDecayAlt.P.JCx1.decay.combine}
\end{lemma}

\begin{proof}
$\min(1,(x_v/a_v)^K)$ equals $(x_v/a_v)^K$ for $v\in A$ and $1$ for $v\notin A$, so the right side is
$\prod_{v\in A}(x_v/a_v)^K\,M\ge\prod_{v\in A}(\lambda_v/a_v)^KM\ge N$.
\end{proof}

\begin{lemma}[Plane-wave bound from places]\label{pf:thmS-carrier:planewave}
Let $\lambda_{i,v}\ge0$ satisfy the \emph{curvature barriers} $\lambda_{i,v}\le\rho^D_{i,v}$ ($1\le i\le r+1$). Assume (H1),
(H2) and the \emph{place-wise wave bound with constants $\lambda$}: for $\eps>0$, $1\le i\le r+1$ and $w\in W_{\mathrm{all}}$
with $\thmScarrierCnt(w)+5\le\thmScarrierNtop$ there are probability vectors $\theta_A$ on
$\thmScarrierWb(\thmScarrierCnt(w)+5)$ such that for all $Y$ of level $i-1$, PSD $\kappa$, $l\in\Lambda^\vee_i\smallsetminus
\{0\}$ and all sets $A$ of places
\[
\Bigl(\prod_{v\in A}|l|_v^5\Bigr)\bigl|\thmScarrierNhat_{\eps,i,w}(Y,\kappa;l)\bigr|\le\Bigl(\prod_{v\in A}
\lambda_{i,v}^5\Bigr)\sum_{w'}\theta_A(w')\thmScarrierNz_{\eps,i-1,w'}(Y,\kappa).
\]
Then the plane-wave bound (H4) holds.
\formal{Formalism.FourierDecayAlt.JC.WaveByPlaces}
\formallabel{Formalism.FourierDecayAlt.P.plane_exponential_bound_of_places}
\end{lemma}

\begin{proof}
Apply the hypothesis with $A=\mathrm{big}_i(l)=\{v:\rho^D_{i,v}<|l|_v\}$. The sum on the right is $\ge0$, since
$\thmScarrierNz_{\eps,k,w'}(Y,\kappa)\ge0$ for PSD $\kappa$ (it is the sum of the entries of the Hadamard product of $\kappa$
and the PSD kernel $\gamma_{\eps,k,w'}$ of \cref{pf:thmS-carrier:gram}, divided by $\thmScarrierWt(w')>0$). Since $l\ne0$,
$|l|_v>0$ for all $v$. \Cref{pf:thmS-carrier:combine} with $x_v=\rho^D_{i,v}$, $a_v=|l|_v$ and $K=5$ gives the bound with
$\thmScarrierMaj(\rho^D_i,l)$.
\end{proof}

\begin{lemma}[Lift bound from places]\label{pf:thmS-carrier:liftplaces}
Let $\lambda_{b,s,v},\lambda'_{b,s,v}\ge0$ satisfy the \emph{lift thresholds} $\lambda_{b,s,v}\le\rho^L_{b,s,v}$ ($b+s\le
r+1$) and $2\lambda'_{r+2-s,s,v}\le\rho^L_{r+2-s,s,v}$ ($2\le s\le r+1$). Assume (H1), (H2), the \emph{lift bound below the
top with constants $\lambda$} (for $i\le r$ and every set $A$ of places, the bound
$(\prod_{v\in A}|t^{-1}-t'^{-1}|_v^5)|\thmScarrierNst(t,t')|\le(\prod_{v\in A}\lambda_{i-s+1,s,v}^5)M_A$, with $M_A$ an
average with probability weights $\theta_A$ of the diagonal terms as in \ref{pf:thmS-carrier:R5}), and the \emph{top diagonal
bound with constants $\lambda'$} (for $i=r+1$, $s\ge2$ and every set $B$ of places, the bound of each diagonal term
$\Re\thmScarrierNst(t,t)$ by $\prod_{v\in B}(|t|_v\lambda'_{r+2-s,s,v})^{10}$ times an average with probability weights
$\theta_B$). Then the lift bound (H5) holds.
\formal{Formalism.FourierDecayAlt.JC.LiftBelowTop}
\formal{Formalism.FourierDecayAlt.JC.LiftTopDiag}
\formallabel{Formalism.FourierDecayAlt.P.off_diagonal_lift_bound_of_places}
\end{lemma}

\begin{proof}
Let $\eps,i,s,w$ be as in \ref{pf:thmS-carrier:R5}. If $i\le r$, the averages $M_A$ are $\ge0$ by \cref{pf:thmS-carrier:gram}.
Put $\theta_\pi=\theta_{\{v:\pi(v)\ne0\}}$ and choose $\pi(v)=1$ if $\rho^L_{i-s+1,s,v}<|t^{-1}-t'^{-1}|_v$ and $\pi(v)=0$
otherwise; \cref{pf:thmS-carrier:combine} gives the bound.

If $i=r+1$ then $s\ge2$. By the Cauchy--Schwarz inequality for the PSD kernel $\kappa\cdot\gamma$,
$|\thmScarrierNst(t,t')|^2\le\Re\thmScarrierNst(t,t)\,\Re\thmScarrierNst(t',t')$. Let $o_v=|t^{-1}-t'^{-1}|_v$ and
$\lambda'_v=\lambda'_{r+2-s,s,v}$, and let $\pi(v)=0$ if $2\lambda'_v\ge o_v$, $\pi(v)=1$ if $2\lambda'_v<o_v$ and
$|t|_v\le|t'|_v$, and $\pi(v)=2$ otherwise; take $\theta_\pi(x_1,x_2)=\theta_{\pi^{-1}(1)}(x_1)\theta_{\pi^{-1}(2)}(x_2)$, a
probability vector on pairs. If $|t|_v\le|t'|_v$ then $o_v\le|t|_v^{-1}+|t'|_v^{-1}\le2|t|_v^{-1}$, so
$|t|_v\lambda'_v\le2\lambda'_v/o_v$, and symmetrically. The elementary several-places inequality
\lref{Formalism.FourierDecayAlt.P.third_polynomial_prime_combine} states: if $|N_{tt'}|^2\le N_tN_{t'}$,
$N_t\le\prod_{\pi(v)=1}(|t|_v\lambda'_v)^{10}S_1$ and $N_{t'}\le\prod_{\pi(v)=2}(|t'|_v\lambda'_v)^{10}S_2$ with
$S_1,S_2\ge0$, then $|N_{tt'}|\le\prod_v\min(1,(2\lambda'_v/o_v)^5)\,(S_1+S_2)/2$. With the top diagonal bound for
$B=\pi^{-1}(1)$ and $B=\pi^{-1}(2)$ and the threshold $2\lambda'_v\le\rho^L_v$ this gives the bound.
\end{proof}

\begin{lemma}[Block bound from the coset sums]\label{pf:thmS-carrier:corT}
Assume (H1), (H2) and the \emph{coset-sum bound with constant $C_4$}: in the situation of \ref{pf:thmS-carrier:R6} there is
$M\le C_4^{fs}\prod_{1\le s'<s}V_{i-s+1,s'}$ for which the coset-sum form of the energies of $S_{\mathrm{base}}$ holds with
constant $M$. Then the block bound (H6) holds with constant $C_4$.
\formal{Formalism.FourierDecayAlt.JC.SumFormConcrete}
\formallabel{Formalism.FourierDecayAlt.P.corollary_T_of_sum_form}
\end{lemma}

\begin{proof}
By \cref{pf:thmS-carrier:gram,pf:thmS-carrier:words}, $S_{\mathrm{base}}$ has PSD Gram kernels and positive weights. For such
energy data, \lref{Formalism.FourierDecayAlt.P.blockEstimateOfCosetSum} states that the coset-sum form with constant $M$
implies the block estimate with the same constant.
\end{proof}

\begin{lemma}[Integral formula from places]\label{pf:thmS-carrier:intplaces}
Let $d_{i,v}\in\R$ with $\prod_vd_{i,v}=\covol(\Lambda_i)\ne0$ for $1\le i\le r+1$. Assume (H1) and the \emph{local integral
formula with constants $d_{\cdot,v}$} at every $v$: for $\eps>0$, $1\le i\le r+1$, $w$ and local points $h,h'$ of size $i-1$,
the function $c\mapsto\gamma^v_{\eps,i,w_v}(h\cdot(1,c),h'\cdot(1,c))$ on $\C^\times$ extends to a continuous $\varphi$ on
$\C$ with $|\varphi(c)|\le C(1+|c|)^{-3}$ and $\int_\C\varphi=d_{i,v}\gamma^v_{\eps,i-1,w_v}(h,h')$. Then the integral formula
(H3) holds.
\formal{Formalism.FourierDecayAlt.JC.ThmPIntPlace}
\formallabel{Formalism.FourierDecayAlt.P.theorem_P_integral_of_places}
\formallabel{Formalism.FourierDecayAlt.P.JCx1.needs.thmPInt_of_covolume_prod}
\end{lemma}

\begin{proof}
$N^1_{\eps,i,w}(Y,\kappa;c)$ is a finite linear combination of products $\prod_v\varphi_{h,h',v}(c_v)$ of the local functions;
the product of the continuous extensions is continuous on $E_\infty$ and bounded by $\prod_vC_v(1+|c_v|)^{-3}$. By Fubini,
\[
\int_{E_\infty}N^1=\frac1{\thmScarrierWt(w)}\sum_{h,h'}\kappa(h,h')\prod_vd_{i,v}\gamma^v_{\eps,i-1,w_v}(\sigma_vh,\sigma_vh')
=\frac{\covol(\Lambda_i)}{\thmScarrierWt(w)}\sum_{h,h'}\kappa(h,h')\gamma_{\eps,i-1,w}(h,h'),
\]
and the last sum is real since $\kappa$ and $\gamma_{\eps,i-1,w}$ are Hermitian. Dividing by $\covol(\Lambda_i)$ gives
$\thmScarrierNhat(Y,\kappa;0)=\thmScarrierNz_{\eps,i-1,w}(Y,\kappa)$.
\end{proof}

\begin{lemma}[Positivity from modules]\label{pf:thmS-carrier:posmodule}
Let $S$ be built from $(T,f,\chi,\nu',m_0,\mathrm{top},A_0)$. Assume that for every $v$ there are submodules $K_j$
($0\le j\le r$) of coefficient functions on which the level forms $\mathrm{form}_{\mathrm{lev}}(T_v,\pm,j,a)$ are sesquilinear
and non-negative for all admissible local words $a$ (Theorem~P in module form), and which contain the slices of the cut test
functions, and that the top kernels are PSD. Then local positivity (H1) holds.
\formal{Formalism.FourierDecayAlt.D.ThmPModule}
\formallabel{Formalism.FourierDecayAlt.P.positivity_of_module}
\end{lemma}

\begin{proof}
At a built place the local kernels at level $i\le r$ are, up to positive normalising factors, the values
$\mathrm{form}_{\mathrm{lev}}(x_p,x_q)$ of a level form on the slices $x_p\in K_j$ of the cut test function at the local
points, and a matrix of values of a non-negative Hermitian form is PSD; at the top level the kernels are PSD by assumption.
This deduction is \lref{Formalism.FourierDecayAlt.P.Dx1.d_thmPPosOfModule}, which identifies the component-family kernels
with these Gram matrices.
\end{proof}

\begin{lemma}[Gram identity at the top]\label{pf:thmS-carrier:gramtop}
Let $\mathrm{vec}_v$ be maps from local points to Hilbert spaces with
$\langle V(q),V(p)\rangle=\prod_v\sqrt{n_v(\sigma_vp)n_v(\sigma_vq)}\langle\mathrm{vec}_v(\sigma_vq),\mathrm{vec}_v(\sigma_vp)
\rangle$ for points $p,q$ of level $r+1$, and assume that for every $v$ and every finite set $Y$ of local points of size $r+1$,
for small $\eps>0$, $RR_v(\eps,r+1,\emptyset,p,q)=\langle\mathrm{vec}_v(q),\mathrm{vec}_v(p)\rangle$ for $p,q\in Y$. Then the top
Gram identity (H8) holds for $V$. At a built place with $g_1,g_2>0$ the hypothesis holds for the local vectors $\mathrm{vec}_v$
attached to the towers, base frames and test functions.
\formal{Formalism.FourierDecayAlt.JC.vecPlace}
\formallabel{Formalism.FourierDecayAlt.P.gram_proposition_of_places}
\formallabel{Formalism.FourierDecayAlt.P.gram_at_built_place}
\end{lemma}

\begin{proof}
For the first part: the empty word has empty local words, $\gamma_{\eps,r+1,\mathrm{emp}}(p,q)=\prod_v\sqrt{n_vn_v}\,
RR_v(\eps,r+1,\emptyset,\sigma_vp,\sigma_vq)$, and finitely many eventual identities hold simultaneously. For the second
part, at a built place the top local kernel is an orbital integral of the cut test function $\mathrm{cut}_\eps f_v$ at the
relevant matrices $\gamma$. By \lref{Formalism.FourierDecayAlt.P.orbital_integral_cutoff_scalar}, the orbital integral of
$\mathrm{cut}_\eps f$ is $\chi(G_f(\gamma)/\eps)$ times that of $f$; since $\chi=1$ on $[1,\infty)$ it equals the uncut value
once $\eps\le G_f(\gamma)$, and the uncut value is the inner product of the local vectors. This identification uses the
specification of the permutations of points \lref{Formalism.FourierDecayAlt.P.ptPerm_spec} and of the weights $G$
\lref{Formalism.FourierDecayAlt.P.ptG_spec}, the fact that the points in question are relevant
\lref{Formalism.FourierDecayAlt.P.relevant_points_holds}, and the Haar property of the radial measure
\lref{Formalism.FourierDecayAlt.P.radial_measure_is_haar}.
\end{proof}

\begin{proposition}[Local parts at built places]\label{pf:thmS-carrier:parts-built}
Assume \ref{pf:thmS-carrier:I1}. Let $C_4,C_6\in\R$, and let $S$ be built from $(T,f,\chi,\nu',m_0,\mathrm{top},A_0)$ with
$g_1,g_2>0$ and $E$ unramified over $F$ at all finite primes. Assume local integral formulas with constants $d_{\cdot,v}$
with $\prod_vd_{i,v}=\covol(\Lambda_i)$, and the \emph{built inputs}: $1\le l\le k-2$; the actual levels and the slice
regularity of the towers; PSD top kernels at all places; (H2); the place-wise wave bound with constants $\lambda$ and the
curvature barriers; the lift bound below the top with constants $\lambda$, the top diagonal bound with constants
$\lambda'$ and the lift thresholds; the coset-sum bound with constant $C_4$; smoothness of the $f_v$; the cut-length bound
$C_6(k\thmScarrierNtop(r+1)^2+\thmScarrierNtop^2(r+1)^2)\le\beta_v-\alpha_v$; model places at all $v$; and the
factorisation of $\langle V(q),V(p)\rangle$ of \cref{pf:thmS-carrier:gramtop}. Then (H1)--(H6) with constant $C_4$ and (H8)
for $V$ hold.
\formal{Formalism.FourierDecayAlt.JCT.BuiltInputs}
\formallabel{Formalism.FourierDecayAlt.P.jct_parts_loc_cpt}
\end{proposition}

\begin{proof}
\emph{Positivity.} Fix $v$. Every admissible local word has length at most $\thmScarrierNtop$. Theorem~P in module form
for compactly supported test functions, \lref{Formalism.FourierDecayAlt.P.theorem_P_module_compact_support}, provides
submodules $K_j$ on which the level forms are non-negative and which contain every function satisfying the growth bounds,
smoothness and a support condition $G_f\ge\eps'$; its inputs are \ref{pf:thmS-carrier:I1}, the radial module form
\lref{Formalism.FourierDecayAlt.P.radial_module_form} (the level forms are positive multiples of the radial forms of the
Bessel transforms), positivity of the density on compactly supported functions
\lref{Formalism.FourierDecayAlt.P.positive_density_compact_support}, sesquilinearity of the level forms on functions of
moderate growth \lref{Formalism.FourierDecayAlt.P.level_form_sesquilinear_moderate}, and the actual levels at level $l+1$.
The slice regularity shows that the slices of the cut test functions satisfy these conditions. Hence
\cref{pf:thmS-carrier:posmodule}, together with the PSD top kernels, gives (H1).

\emph{Integral formula.} $\covol(\Lambda_i)\ne0$ by \lref{Formalism.FourierDecayAlt.P.unipotent_dual_lattice_covolume},
which gives the covolume of the dual unipotent lattice as an explicit positive expression in $g$ and the root discriminant of
$E$ (using that $E$ is unramified); \cref{pf:thmS-carrier:intplaces} gives (H3).

\emph{The rest.} (H2) is part of the built inputs; \cref{pf:thmS-carrier:planewave,pf:thmS-carrier:liftplaces,pf:thmS-carrier:corT}
give (H4), (H5) and (H6); \cref{pf:thmS-carrier:gramtop} gives (H8).
\end{proof}

\begin{lemma}[Energy from the model]\label{pf:thmS-carrier:lemE}
Assume the model lemma (\cref{pf:thmS-carrier:lemmamodel}). There is $C_6>0$ such that for every $S$ built from
$(T,f,\chi,\nu',m_0,\mathrm{top},A_0)$ for which the cut test functions $\mathrm{cut}_\eps f_v$ are smooth, every cut interval
has length at least $C_6\bigl(k\thmScarrierNtop(r+1)^2+\thmScarrierNtop^2(r+1)^2\bigr)$, and every place is a model place for
model data $(L_v,j_v,\lambda_v,c,C_{m,v},\sigma_v)$, the energy limit (H7) holds with the model energies
$j^{\mathrm{mod}}_v$ of \cref{pf:thmS-carrier:modelplace-def}.
\formallabel{Formalism.FourierDecayAlt.P.lemma_E_of_model}
\end{lemma}

\begin{proof}
Let $C_2,C_6>0$ be the constants of the model lemma. By the residue reduction
\lref{Formalism.FourierDecayAlt.P.lemma_E_of_residue}, which uses the level-zero formula
\lref{Formalism.FourierDecayAlt.P.level_zero_holds} and the goodness of admissible words
\lref{Formalism.FourierDecayAlt.P.good_of_admissible}, it suffices to show, for $w\in W(S)$ and every $v$, that
\[
\mathcal E_v(w_v;\mathrm{cut}_\eps f_v)\longrightarrow\mathrm{size}_v(w_v)\,C_{m,v}\int\sum_\rho|j_{v,\rho}|^2\,d\mu_0
\qquad(\eps\to0^+);
\]
indeed the reduction states that at a built place $RR_v(\eps,0,a,\emptyset,\emptyset)$ is the energy
$\mathcal E_v(a;\mathrm{cut}_\eps f_v)$ divided by $\Gamma_0$. Fix $v$ and $a=w_v$. The model place conditions give that the
model system with box length $L_b=C_2\thmScarrierNtop(r+1)^2$ admits $a$ and has $\mathrm{size}(a)=\mathrm{size}_v(a)$;
moreover $\len(a)\le\thmScarrierNtop$ and $r+1\le\thmScarrierNtop=20(r+1)$. The model lemma, applied with
$N_{\mathrm{tot}}=\thmScarrierNtop$, the interval $[\alpha_v,\beta_v]$ and the model data, gives
$\thmScarrierEn^{\mathrm{mod}}_v(a;\delta)\to\mathrm{size}(a)\int\sum_\rho|j_{v,\rho}|^2d\mu_0$ as $\delta\to0^+$. Since
$\sigma_v>0$, $\sigma_v\eps\to0^+$ as $\eps\to0^+$; by \eqref{eq:pf:thmS-carrier:modelenergy},
$\mathcal E_v(a;\mathrm{cut}_\eps f_v)=C_{m,v}\thmScarrierEn^{\mathrm{mod}}_v(a;\sigma_v\eps)$ for $\eps>0$, and the claim
follows.
\end{proof}

\begin{theorem}[Real inputs at leading-form places]\label{pf:thmS-carrier:compact}
Assume \ref{pf:thmS-carrier:I1}, \ref{pf:thmS-carrier:I2}, \ref{pf:thmS-carrier:I3} and the model lemma, and let $C_b\in\R$.
There are $C_6>0$, $C_c$ and $C_4\ge1$ such that for every CM field $E$, every Hida setting $S$ over $E$ with
finite-dimensional local spaces, place packages $p$, test functions $f$, cut-off $\chi$, number $A_0$, cut $c$ with cut-off
function $\chi$, and map $V$: if $S$ is a leading-form place with constants $(C_b,C_c)$, the slice budget for $A_0$ holds,
every cut interval has length at least $C_6\bigl(k\thmScarrierNtop(r+1)^2+\thmScarrierNtop^2(r+1)^2\bigr)$, and
$\langle V(q),V(p)\rangle$ factors as in \cref{pf:thmS-carrier:gramtop} through the local vectors of the place packages, then
there are model data $(L,j,\lambda,C_m,\sigma)$ for which every place is a model place, and for all $(\mathcal K,\Theta,e)$
the energy datum $\thmScarrierDat_S(\mathcal K,\Theta,e)$ satisfies the real inputs with constants
$(C_4,\prod_vj_v^{\mathrm{mod}})$ and top vectors $V$.
\formal{Formalism.FourierDecayAlt.KB.ArchimedeanUniformConstantsCompactSupport}
\formallabel{Formalism.FourierDecayAlt.P.RealInputsCompactSupport.real_inputs_closed_compact_support_of}
\end{theorem}

\begin{proof}
Let $C_6$ be the constant of \cref{pf:thmS-carrier:lemE}; let $c_1,c_2,c_7$ be the built-place constants of the actual levels
of built settings \lref{Formalism.FourierDecayAlt.P.levels_of_built_setting} (at a built place the towers carry actual
levels), of the slice regularity at built places \lref{Formalism.FourierDecayAlt.P.slices_finite_regularity_built} (under
the slice budget the slices of the cut test functions are regular), and of \ref{pf:thmS-carrier:I2}; and let $(c_3,C_4)$ be
given by \ref{pf:thmS-carrier:I3} for $C_b$. Put $C_c=\max(c_1,c_2,c_3,c_7)$. Let $S,p,f,\chi,A_0,c,V$ be as in the
statement.

Being a built place is monotone in the constant, so the leading-form place is a built place for each of the constants
$c_1,c_2,c_7$, with tower, base frames and top kernels taken from the place packages; it is also a leading-form setting with
constant $c_3$. By \ref{pf:thmS-carrier:I3} we obtain (H1), the local integral formulas at all places, $\lambda$ with the
place-wise wave bound and the curvature barriers, $\lambda,\lambda'$ with the lift bounds and thresholds, and the coset-sum
bound with constant $C_4$. The top kernels are PSD at every place by
\lref{Formalism.FourierDecayAlt.P.top_positive_semidefinite}. By \ref{pf:thmS-carrier:I2} (the admissibility of the place
data and of the words is part of the leading-form condition) every place is a model place for some model data
$(L_v,j_v,\lambda_v,C_{m,v},\sigma_v)$ and the cut $c$. Thus the built inputs of \cref{pf:thmS-carrier:parts-built} hold. The
local integral formulas hold with $d_{i,v}=\mathrm{rd}_E^{2\nu_i}/\delta_g$ (with an explicit exponent $\nu_i$ and an explicit normalising factor $\delta_g>0$), and the product of these
numbers over the places is $\covol(\Lambda_i)$ by \lref{Formalism.FourierDecayAlt.P.PBJoint.prod_dlB}, which computes this
product when $E$ is unramified. \Cref{pf:thmS-carrier:parts-built} gives (H1)--(H6) and (H8); \cref{pf:thmS-carrier:lemE},
whose smoothness hypothesis holds by \lref{Formalism.FourierDecayAlt.P.cutoff_smooth} (the cut test functions are smooth),
gives (H7) with the model energies. Now \cref{pf:thmS-carrier:realconcrete} applies.
\end{proof}

\subsubsection{Proof of the compact-model theorem}\label{sec:pf:thmS-carrier:compactproof}
\begin{proof}[Proof of \cref{pf:thmS-carrier:compactmodel}]
Let $(C_6,C_c,C_4)$ be the constants of \cref{pf:thmS-carrier:compact} for $C_b$. Apply (GS$_{C_b}$) with $(C_6,C_c)$; it gives
$C_{A,0}$. We claim that $(C_4,C_{A,0})$ are constants as in \cref{pf:thmS-carrier:main}. Given the further constants,
(GS$_{C_b}$) provides $k_0$ and then $C'_0,C_{l,0}$; given the data and the hypotheses of \cref{pf:thmS-carrier:main}, it
provides the carrier data bound: a Hida setting $S$ with place packages $p$ and $(f,\chi,A_0,c,\tau,Q,Q^{\mathrm i})$ and
isometries with (i)--(iv) of \cref{pf:thmS-carrier:gsdef}. By (i) and (ii), \cref{pf:thmS-carrier:compact} applies and gives
model data with model places everywhere such that $\thmScarrierDat_S(\mathcal K,\Theta,e)$ satisfies the real inputs with
constants $(C_4,\prod_vj^{\mathrm{mod}}_v)$ and top vectors $V_{\mathrm{pt}}$ for all $(\mathcal K,\Theta,e)$, in particular
for those of $\thmScarrierDat^X_S$. By (iv), $\prod_vj_v^{\mathrm{mod}}=j_X$. Taking for $T_v,\nu'_v,m_{0,v}$ the
data of the place packages, the remaining clauses of the carrier bound are among (iii) and the constants of $S$.
\end{proof}

% ---------------------------------------------------------------------------
\subsection{The global side}\label{sec:pf:thmS-carrier:X}

The Hida setting of the carrier is constructed from the spectral data of the global datum. At every infinite place of $F$ the
spectral ties provide a discrete series representation $\sigma_B$ of the local unitary group and the stabiliser
representations $X.V_{\mathrm{loc}}(u,\cdot)$, $X.V_{\mathrm{loc}}(u',\cdot)$; a \emph{choice} $\mathcal C$ consists of
matrices $Q,Q^{\mathrm i}$ putting the local Hermitian form into the standard form $H_{k,r}$. The test function at the place
$j$ is built from the matrix coefficients of two $K$-finite vectors $a_j,b_j$ of $\sigma_B$, two vectors
$\xi_j\in X.V_{\mathrm{loc}}(u)$, $\xi'_j\in X.V_{\mathrm{loc}}(u')$ and a scalar $c_j$.

\begin{definition}[Pins and leaves]\label{pf:thmS-carrier:pinsleaves}
Consider constants $C_A$, $C_0$, $C_G$, $C_\Delta$, $c_T$, $C_b'$, $C_g$, $c_b$, $C_w$, $c_{Tf}$, $C_T$, $C_{0f}$, $C'$, $C_l$
and data $D$, $L$, $T$, $u$, $u'$, $e$, $\tau$, $C_p$, $\Theta'$, $i_0$, $i_0'$, $\nu_q$ as in \cref{pf:thmS-carrier:main}, with
$u\ne u'$, $c_b,c_{Tf}>0$, $D$ tied, $C'\ge29$, $C_l\ge1000+|C_{\mathrm{orb}}|$, $k\ge1801$, the count hypotheses, the tie of
places, the size hypotheses, $i_0=i'_0$, the existence of point representatives and the pinning of $\nu_q$. Here
$C_{\mathrm{orb}}$ is the constant provided by \lref{Formalism.FourierDecayAlt.P.orbital_integral_converges}, the
convergence of the orbital integrals of the Bessel kernels.
\begin{itemize}
\item The \emph{pinning statement} asserts: for all such constants and data there are $\theta$, a choice $\mathcal C$ with
  Wu--Bessel data $B$, and $K$-finite vectors $a_j,b_j$ and vectors $\xi_j,\xi'_j$ at every place, such that the trace
  hypothesis for $(u,u')$ holds, the traces $\mathrm{Tr}(D,u),\mathrm{Tr}(D,u')$ are positive, the local goodness
  conditions on the weights and on the Harish-Chandra parameters hold at every place, the coefficients are pinned at scale
  $1$ by $(a_j,b_j,\xi_j,\xi'_j)$, and the scalars $\mathrm{scl}(B,u,u',\cdot)$ and $c(B,u,\cdot)$ are non-zero at every
  place.
\item The \emph{leaf statement} asserts: for all such constants and data and all auxiliary mass data there are $\theta$, a
  choice $\mathcal C$, $K$-finite $a_j,b_j$, vectors $\xi_j,\xi'_j$ and scalars $c_j$ such that the \emph{leaf condition}
  holds: for all top lifts, type projectors and $A_0$, the carrier built from the test functions attached to
  $(a,b,\xi,\xi',c)$ satisfies the pinning covariance.
\end{itemize}
\formal{Formalism.L2.SG.LocWSOf.Choice}
\formal{Formalism.L2.SG.LocWSOf.WuBessel}
\formal{Formalism.L2.SG.Pf.CarrierT.fcT}
\formal{Formalism.L2.SG.Pf.CarrierXAll.XPinsP}
\formal{Formalism.L2.SG.Pf.CarrierXAll.XLeafAt}
\formal{Formalism.L2.SG.Pf.CarrierXAll.XLeavesP}
\end{definition}

\begin{proposition}[Pins]\label{pf:thmS-carrier:xpins}
The existence of $K$-finite vectors with non-zero Bessel functional implies the pinning statement. (The formal statement also
assumes the uniqueness of Bessel models with large gap, which is not used in its proof.)
\formal{Formalism.L2.SG.Pf.CarrierXAll.KFinNZ}
\formallabel{GlobalTheory.CarrierGap.xPinsP_ofG}
\end{proposition}

\begin{proof}
Fix the constants and data. The gap hypothesis for $D$ holds by
\lref{Formalism.L2.SG.ThmSOwn.gap_hypothesis_of_hypotheses}, which derives it from the count hypotheses when $C_l\ge1000$
and $C'\ge29$. Choose a choice $\mathcal C$ with $\theta=1$, which exists by \lref{Formalism.L2.SG.Pf.choiceExists}, and
Wu--Bessel data $B$ for it, which \lref{Formalism.L2.SG.Pf.wuBesselOfLoc_all} assembles from the Bessel functional of the
Whittaker model, the local factor, and the transfers of discrete series, Haar measures and column types. For every word $w$
and index $t$ the local Bessel function $\varphi(B,w,t)$ has non-zero finite $L^2$-norm by
\lref{Formalism.L2.SG.Pf.GoodS.l2SG_phiOf_ne_zero_and_lt_infty}; hence $c(B,w,t)>0$, and
$\mathrm{scl}(B,u,u',t)=\bigl(c(B,u,t)\dim X.V_{\mathrm{loc}}(u',t)/\|\varphi(B,u',t)\|^2\bigr)^{1/2}\ne0$, because the dimension
is positive: on a zero space $\varphi(B,u',t)$ would vanish. By hypothesis (which applies under the gap hypothesis) there are
$K$-finite $a_j,b_j$ with non-zero Bessel functionals $B_u(a_j)$, $B_{u'}(b_j)$; put $\xi_j=B_u(a_j)/\|B_u(a_j)\|^2$ and
$\xi'_j=B_{u'}(b_j)/\|B_{u'}(b_j)\|^2$, so that the pinning condition (a product of the inner products
$\langle\xi_j,B_u(a_j)\rangle$ and $\langle B_{u'}(b_j),\xi'_j\rangle$, both equal to $1$) holds. The trace hypothesis holds
for all pairs of words by \lref{Formalism.MainNotions.SchurRelationsDiscreteSeries.hypTT_all}; $\mathrm{Tr}(D,u)=T_u>0$
since $T_u$ is prime; the goodness of the weights holds at all places by \lref{Formalism.L2.SG.Pf.AmpWinRX.goodW_all};
and the goodness of the Harish-Chandra parameters follows from the discrete-series tie of $D$ by
\lref{Formalism.L2.SG.Pf.goodSRest_all}.
\end{proof}

\begin{proposition}[Leaves]\label{pf:thmS-carrier:xleaves}
The uniqueness of Bessel models with large gap and the pinning statement imply the leaf statement.
\formallabel{GlobalTheory.CarrierGap.xLeavesP_of_pinsPG}
\end{proposition}

\begin{proof}
Fix the constants and data, and take the pins $(\theta,\mathcal C,B,a,b,\xi,\xi')$. Let $n=2r_S+k+1$ be the rank and
$h=C_ln^{19}f^5L$ with $L=L.\mathrm{sc}.L$. We verify the numerical hypotheses of the leaf construction: $C_l\ge1000$; $k\ge5$;
the scale parameter satisfies $L\ge\frac12$ under the hypotheses; $f\ge2$, so $f^5\ge32$ and $16C_ln^{19}\le h$; and
$n\ge1801$. The gap condition of Wu decay holds at $h$: by \lref{Formalism.L2.SG.Pf.WuDec.gapOK_of_big} it holds at gap
$n^{19}$ when $k\ge5$ and $n\ge1801$, and it is monotone in the gap. The elementary inequality
$\max\bigl(C_{\mathrm{orb}}(r_S+1)^2(r_Sn+1),0\bigr)\le h/2$ holds by \lref{Formalism.L2.SG.Pf.CarrierXAll.ac_le} (valid for
$C_l\ge1000+|C_{\mathrm{orb}}|$, $f\ge2$, $L\ge\frac12$). The Harish-Chandra parameters at the places have gap at least $h$ by
\lref{Formalism.MainNotions.BridgeLemma.szh_gap}, which extracts this gap from the standing hypotheses. Finally, with
explicit normalising scalars $c_j$ whose definition uses the uniqueness of Bessel models,
\lref{GlobalTheory.CarrierGap.xLeafAt_of_pinsG} establishes the leaf condition from the pins, these scalars and the gap
conditions.
\end{proof}

\begin{theorem}[Carrier data at built carriers]\label{pf:thmS-carrier:cdbuilt}
Let $C_b$ satisfy (TL$_{C_b}$) and assume the capacity bound from Harish-Chandra characters. Let the data of
\cref{pf:thmS-carrier:pinsleaves} satisfy the standing hypotheses with $C_l\ge1000$ and $k\ge1801$; let $C_c,C_6\in\R$, words
$u,u'$, a bijection $e$ of places tied to $X$, constants $C_A,d_e,d'_e,l$, base indices $i_0=i'_0$, a scale $s_c>0$, measures
$\nu_q$ pinned for $(u,u')$ such that point representatives exist, a number $j_n$, and leaf data $(a,b,\xi,\xi',c)$ satisfying
the leaf condition. Write $\thmScarrierNtop=20(r_S+1)$ and $\Delta'=\Delta/2-5/2$, and let $C_d$, $C_e$ and $C_{\mathrm{fd}}$ be
the three fixed constants of the existence of carrier data, of the energy of built settings and of \cref{pf:thmS-carrier:lemE}
(they reappear in the proof of \cref{pf:thmS-carrier:xside-proof}). Assume:
\begin{enumerate}[label=(\alph*)]
\item the Harish-Chandra coordinates at the non-top places are those of the standing data, $d_e,d'_e$ are the words of
  $u,u'$, all places are admissible, $d_e,d'_e$ are admissible, $1\le l\le k-2$ and $d_e(l)\ne d'_e(l)$;
\item $6\thmScarrierNtop^3+6\thmScarrierNtop\le\Delta'$; $c\,(k\thmScarrierNtop(r_S+1)^2+\thmScarrierNtop^2(r_S+1)^2+
  \thmScarrierNtop^3)\le\Delta'$ for $c\in\{C_c,C_d,C_e\}$ and the same without the term $\thmScarrierNtop^3$ for
  $c\in\{C_6,C_{\mathrm{fd}}\}$; and $C_c,C_d,C_e\le C_A$;
\item there is $h$ such that the Harish-Chandra parameters at all places have gap at least $h$, the gap condition of Wu decay
  holds at $h$, and $A(|C_c|+|C_d|+|C_e|,k,r_S)\le h/2$ for the decay threshold $A$;
\item (energy) for all top lifts $\Theta$, type projectors and $A_0$ for which the Hida setting of the test functions of the
  leaf data is built, $\prod_v(\|J_{f_v}\|^2/\Gamma_0)=j_n$.
\end{enumerate}
Then the carrier data bound with constants $(C_6,C_c,C_b)$ holds, with energy constant $j_n$.
\formallabel{Formalism.L2.SG.Pf.CarrierXAll.carrierDataBuilt_of_hyp}
\end{theorem}

\begin{proof}
\emph{Preliminaries.} The test functions $f_j$ attached to the leaf data are smooth, because the vectors $b_j$ are $K$-finite
(\lref{Formalism.L2.SG.Pf.CarrierXAll.fcT_smooth}). By (a), the parameters and types of the places are those of the
standing data (\lref{Formalism.L2.SG.Pf.CarrierXAll.parNumsE_of} and \lref{Formalism.L2.SG.Pf.CarrierXAll.typesE_of}), and type
projectors $TP$ exist (\lref{Formalism.L2.SG.Pf.CarrierXAll.tpOf}, from
\lref{Formalism.FourierDecayAlt.P.type_projectors_exist}). Let $\Theta_0$ be the identity top lift and
$A_0=A(|C_c|+|C_d|+|C_e|,k,r_S)\ge0$. For $|c|\le|C_c|+|C_d|+|C_e|$,
\[
c\,(r_S+1)^4\bigl(\thmScarrierNtop(2r_S+k+1)^2+(2r_S+k+1)^3\bigr)\le A_0,
\]
since the left side increases with $|c|$ and \lref{Formalism.L2.SG.Pf.CarrierT.poly_le_decayThreshold} gives the bound for
$|C_c|+|C_d|+|C_e|$.

\emph{A built place at the identity lift.} By \lref{Formalism.L2.SG.Pf.CoefCCT.coefCC_carrierInTc}, the gap hypotheses (c)
and the $K$-finiteness of $a_j,b_j$ imply that the coefficients satisfy the growth bounds with constant $A_0$. By construction the
cut interval has length $\Delta'$, so (b) is the long-cut condition for $C_d$; the
standard grids satisfy the grid condition (\lref{Formalism.L2.SG.Pf.CarrierT.gridsOK_ITc}, which uses
$6\thmScarrierNtop^3+6\thmScarrierNtop\le\Delta'$). With the frames of the towers of $X$
(\lref{Formalism.L2.SG.Pf.CarrierT.frames_carrierInT}) and the capacity bound
(\lref{Formalism.L2.SG.Pf.CarrierXAll.coefCap_ITc}), \lref{Formalism.L2.SG.Pf.CarrierT.isBuiltPlace_carrierInT}
shows that the Hida setting of the test functions with $(\Theta_0,TP,A_0)$ is a built place with constant $C_d$. Then
\lref{Formalism.L2.SG.Pf.CarrierT.exists_Cc_data} provides the carrier data of this built place.

\emph{Top lifts.} By (TL$_{C_b}$), \lref{Formalism.L2.SG.Pf.CarrierXAll.lift_compact_test_towers} replaces $\Theta_0$ by top
lifts $\Theta$ of class $D^+$ with constant $C_b$, keeping the carrier data. For $(\Theta,TP,A_0)$: the leaf condition gives the
pinning covariance of the carrier; the capacity bound holds as before; the weights are tied to those of $X$
(\lref{Formalism.L2.SG.Pf.Carrier.CarPinVt.wtTie_carrierInT}); the joint tensor identity holds at built settings
(\lref{Formalism.L2.SG.Pf.CarrierXAll.jointTensor_compact_test}, which uses $i_0=i'_0$, the pinning of $\nu_q$ and $s_c>0$);
and the base conditions at the empty point hold
(\lref{Formalism.L2.SG.Pf.Carrier.CarPinVt.lZero_of_choice}, \lref{Formalism.L2.SG.Pf.Carrier.CarPinVt.ptSome_of}, and \lref{Formalism.L2.SG.Pf.Carrier.CarPinVt.rZero_carrierInT}).

\emph{Assembly.} \Lref{Formalism.L2.SG.Pf.CarrierXAll.carrierDataBuilt_of_towers} verifies the four conditions of the
carrier data bound (\cref{pf:thmS-carrier:gsdef}) from these facts. Conditions (i)--(iii) follow from the built-place property
with constant $C_c$, the slice budget of $A_0$, the cut-length conditions for $C_6$, and the factorisation of the top vectors
through the joint tensor. Condition (iv) is the energy identity: if every place is a model place for model data
$(L,j,\lambda,C_m,\sigma)$, then by \cref{pf:thmS-carrier:lemE} the energy limit (H7) holds with the model energies
$j^{\mathrm{mod}}_v$, and \lref{Formalism.FourierDecayAlt.P.energy_of_built_setting} (which, given the limit of the
regularised cut-offs \lref{Formalism.FourierDecayAlt.P.beta_cutoff_limit}, states that local energies satisfying (H7) at a
built place are $\|J_{f_v}\|^2/\Gamma_0$) gives $j^{\mathrm{mod}}_v=\|J_{f_v}\|^2/\Gamma_0$ for every $v$; so
$\prod_vj^{\mathrm{mod}}_v=\prod_v(\|J_{f_v}\|^2/\Gamma_0)=j_n$ by (d).
\end{proof}

\begin{theorem}[Energy of the compactly supported test functions]\label{pf:thmS-carrier:energycpt}
Let $T$ be spectral ties with Hida setting $X$, $\mathcal C$ a choice, $u,u'$ words and $e$ a bijection of places tied to $X$,
with $v_{\mathrm{host}}=\mathrm{rd}(E)/2$; let $(a,b,\xi,\xi',c)$ be leaf data, $\Theta$ top lifts, $TP$ type projectors and
$A_0\in\R$, and let the mass data of $T$ be given (with their auxiliary parameters). If the carrier of the Hida setting of the
test functions satisfies the pinning covariance and the setting is built, then
\[
\prod_v\Bigl(\|J_{f_v}\|^2\Big/\Gamma_0\bigl(k,\mathrm{de}(g,\mathrm{rd}(E))\bigr)\Bigr)=j_X,
\]
the energy constant of the mass data.
\formallabel{Formalism.L2.SG.Pf.CarrierXAll.energy_compact_test_bound}
\end{theorem}

\begin{proof}
The carrier is built from $X$: its base is the base constructed from $X$, the choice, the test functions and the top weights
(this holds by construction). Moreover $g=g_0g_1>0$, $\mathrm{rd}(E)$ agrees with the root discriminant of the host, and $k\ge5$.
The transfer of the matrices $\hat Z$ to the carrier (\lref{Formalism.L2.SG.Pf.Carrier.zhatTr_carrierOf}) yields the norm
bridge (\lref{Formalism.L2.SG.Pf.Carrier.JointPinNormSt.normBridge_of_zhat}) relating the local norms of the carrier to
the local operators of $X$. By \lref{Formalism.L2.SG.Pf.Carrier.JointPin.placeNormG_of_pin}, the pinning covariance, the
norm bridge and the invariance of the Hilbert--Schmidt norm under conjugation by isometries give the place-norm identity: at
every place the local norms of the carrier are the Hilbert--Schmidt norms of the pinned local operators of $X$. Finally
\lref{Formalism.L2.SG.Pf.Carrier.EnergyNF.energy_norm_bound_third_form} deduces the claimed identity from the place-norm
identity by a change of variables on the local unitary groups, using the product--convolution identity
\lref{Formalism.L2.SG.Pf.ProdConvCut2.prodConv_holds} and the standard mass lemma
\lref{Formalism.L2.S.Pv.MSTypP.lemmaMSStd_holds}.
\end{proof}

\begin{theorem}[Carrier data from leaves]\label{pf:thmS-carrier:xside-proof}
Assume (TL$_{C_b}$) and the leaf statement. Then the global-side bound (GS$_{C_b}$) of \cref{pf:thmS-carrier:gsdef} holds.
\formallabel{Formalism.L2.SG.Pf.CarrierXAll.carrierXSideP_of_leaves}
\end{theorem}

\begin{proof}
Let $C_6>0$ and $C_c$ be given. Let $C_d$ be the constant of the existence of carrier data, $C_e$ the constant of the energy
of built settings, and $C_{\mathrm{fd}}$ the constant $C_6$ of \cref{pf:thmS-carrier:lemE} (three fixed numbers). Put
$C_{A,0}=\max(C_c,C_d,C_e,1)$ and
\[
C_q=8400\bigl(|C_c|+|C_d|+|C_6|+|C_e|+|C_{\mathrm{fd}}|\bigr)+48120 .
\]
Given $C_A\ge C_{A,0}$ and the constants $C_0,C_G,C_\Delta,c_T,C'_b$, \lref{Formalism.L2.SG.Pf.CarrierXAll.lenI_big}
provides $k_1$ such that for $k\ge k_1$ the size hypotheses force the length of the cut interval to satisfy
\begin{equation}\label{eq:pf:thmS-carrier:lenI}
C_q\bigl(k(r_S+1)^3+(r_S+1)^4\bigr)\le\Delta/2-5/2 .
\end{equation}
Put $k_0=\max(k_1,1801)$, $C'_0=29$ and $C_{l,0}=\max\bigl(1000,25000+21(|C_c|+|C_d|+|C_e|),1000+|C_{\mathrm{orb}}|\bigr)$. Let
the data satisfy the hypotheses; the leaf statement gives leaf data $(\theta,\mathcal C,a,b,\xi,\xi',c)$. From
\eqref{eq:pf:thmS-carrier:lenI} and elementary polynomial inequalities in $k$ and $r_S$
(\lref{Formalism.L2.SG.Pf.CarrierXAll.polys}) we obtain, with $\thmScarrierNtop=20(r_S+1)$, the cut-length conditions
$c\,\bigl(k\thmScarrierNtop(r_S+1)^2+\thmScarrierNtop^2(r_S+1)^2+\thmScarrierNtop^3\bigr)\le\Delta/2-5/2$ for
$c\in\{C_c,C_d,C_e\}$, the same without the term $\thmScarrierNtop^3$ for $c\in\{C_6,C_{\mathrm{fd}}\}$, and the grid condition
$6\thmScarrierNtop^3+6\thmScarrierNtop\le\Delta/2-5/2$; moreover $C_c,C_d,C_e\le C_A$. As in \cref{pf:thmS-carrier:xleaves}, the
gap $h=C_ln^{19}f^5L$ satisfies the Wu decay gap condition, the Harish-Chandra parameters have gap at least $h$, and the decay
threshold satisfies $A(|C_c|+|C_d|+|C_e|,k,r_S)\le h/2$, because $(25000+21(|C_c|+|C_d|+|C_e|))n^7\le C_ln^{19}$. Now
\cref{pf:thmS-carrier:cdbuilt}, with the capacity bound from Harish-Chandra characters
\lref{Formalism.L2.SG.Pf.CarrierT.capacity_from_hermitian_char} and $j_n=j_X$, produces the carrier data bound with constants
$(C_6,C_c,C_b)$: its hypotheses (a)--(c) were just verified, and its energy hypothesis (d), that
$\prod_v(\|J_{f_v}\|^2/\Gamma_0)=j_X$ for every built carrier, is \cref{pf:thmS-carrier:energycpt}, which applies because the
leaf condition gives the pinning covariance.
\end{proof}

\begin{remark}\label{pf:thmS-carrier:rem-facts}
The statements of this subsection are formalised under four auxiliary hypotheses: decay of low $K$-type matrix coefficients
of discrete series with gap $h\ge120n^3$, of order $\|x\|^{-h/(8\sqrt{2(r+1)})}$; non-vanishing of a Bessel integral of matrix
coefficients; capacity bounds for the carrier test functions; and finiteness of a weighted $L^1$-norm of low $K$-type matrix
coefficients. They are discharged by typeclass instances proved in the development \cite{LeanSource2026}, and are not
treated in this section.
\end{remark}

% ---------------------------------------------------------------------------
\subsection{Bounds at leading-form settings}\label{sec:pf:thmS-carrier:built}

\begin{theorem}[Built bounds]\label{pf:thmS-carrier:builtbounds}
Statement \ref{pf:thmS-carrier:I3} holds: assuming \ref{pf:thmS-carrier:I1}, for every $C_b$ there are $C_c$ and $C_4\ge1$ such
that every leading-form setting $(S,p,f,\chi,A_0)$ with constants $(C_b,C_c)$ satisfies (H1), the local integral formulas at
all places, the place-wise wave bound and the curvature barriers for some $\lambda$, the lift bound below the top and the
top diagonal bound with the lift thresholds for some $\lambda,\lambda'$, and the coset-sum bound with constant $C_4$.
\formallabel{Formalism.FourierDecayAlt.Final.pb_builtBounds2Cpt}
\end{theorem}

The local statements used in the proof are of \emph{block-PSD} type.

\begin{definition}[Block positivity; local conditions]\label{pf:thmS-carrier:blockdef}
For finite sets $Y_1,Y_2$ and kernels $a$ on $Y_1$, $b$ on $Y_2$ and $w$ on $Y_1\times Y_2$, the triple $(a,w,b)$ is
\emph{block-PSD} if the kernel $\left(\begin{smallmatrix}a&w\\ w^*&b\end{smallmatrix}\right)$ on $Y_1\sqcup Y_2$ is PSD. For a
Hida setting $S$ and a place $v$ let $\tilde\gamma^v_{\eps,i,a}=\gamma^v_{\eps,i,a}/\mathrm{size}_v(a)$, let
$\widehat{\tilde\gamma}_l$ be its Fourier coefficient at frequency $l\in\C$ in the variable $c$ of a last block $(1,c)$, and let
$\mathrm{mix}^\pm$ denote averages of $\tilde\gamma^v$ over probability vectors on local words with five more letters.
\begin{itemize}
\item The \emph{local wave condition} with constants $X_i$: for each level $i$ and local word $a$ there are such probability
  vectors with $(\tilde\gamma_{i-1},\widehat{\tilde\gamma}_l,\tilde\gamma_{i-1})$ and $(X_i^5\mathrm{mix}^+_{i-1},u\,l^5
  \widehat{\tilde\gamma}_l,X_i^5\mathrm{mix}^-_{i-1})$ block-PSD for some unit $u$, for all $l$ and all $Y$ of level $i-1$; so
  the Fourier coefficient at frequency $l$ decays like $(X_i/|l|)^5$ relative to averages of the kernel.
\item The \emph{local lift condition} with constants $\Lambda_{b,s}$: the analogous block condition for the off-diagonal
  kernels $(p,q)\mapsto\tilde\gamma(p\cdot(s,t),q\cdot(s,t'))$ with weight $(t^{-1}-t'^{-1})^5$ and diagonal entries
  $\Lambda_{b,s}^5\mathrm{mix}^\pm$.
\item The \emph{local top condition} with constants $\Lambda'$: at the top level, $(|t|\Lambda')^{10}\mathrm{mix}-\tilde\gamma$
  is PSD on the points $p\cdot(s,t)$.
\item The \emph{germ condition}: continuity and integrability of the one-variable kernels, and a block-PSD representation of
  the Fourier coefficients $\widehat{\tilde\gamma}$ as limits at infinity of kernels in an inverted variable.
\end{itemize}
A coefficient $F\colon\mathrm{Mat}_{k,r}\to\Hom(W_q,W_p)$ is \emph{compactly supported} if it is smooth on $G_{k,r}$, its
support has bounded Frobenius norm, and $G_f(z,z';g)\ge\eps'$ on its support for some $\eps'>0$.
\formal{Formalism.FourierDecayAlt.KB.BlockPSD}
\formal{Formalism.FourierDecayAlt.KB.WaveLoc}
\formal{Formalism.FourierDecayAlt.KB.LiftLoc}
\formal{Formalism.FourierDecayAlt.KB.TopLoc}
\formal{Formalism.FourierDecayAlt.KB.GermLoc}
\formal{Formalism.FourierDecayAlt.KB.CptCoef}
\end{definition}

\begin{lemma}[Block Schur inequality]\label{pf:thmS-carrier:blockschur}
If $(\kappa_a,\kappa_w,\kappa_b)$ and $(a_v,w_v,b_v)$, for $v$ in a finite set, are block-PSD on $(Y_1,Y_2)$, then
\begin{equation}\label{eq:pf:thmS-carrier:blockschur}
\Bigl|\sum_{p\in Y_1,q\in Y_2}\kappa_w(p,q)\prod_vw_v(p,q)\Bigr|\le\frac12\Bigl(\Re\sum_{p,p'\in Y_1}\kappa_a\prod_va_v+
\Re\sum_{q,q'\in Y_2}\kappa_b\prod_vb_v\Bigr).
\end{equation}
\formal{Formalism.FourierDecayAlt.KB.BlockSchur}
\end{lemma}

\begin{proof}
By the Schur product theorem the entrywise product of the block kernels, which is the block kernel of
$(\kappa_a\prod a_v,\kappa_w\prod w_v,\kappa_b\prod b_v)$, is PSD. Evaluating its quadratic form at the vector equal to $1$ on
$Y_1$ and to $-u$ on $Y_2$, where $|u|=1$ is chosen with $\bar u\sum\kappa_w\prod w_v=|\sum\kappa_w\prod w_v|$, gives
$\Re\sum\kappa_a\prod a_v+\Re\sum\kappa_b\prod b_v-2|\sum\kappa_w\prod w_v|\ge0$.
\end{proof}

This is how local decay is turned into decay of products.

\begin{theorem}[Compact approximation]\label{pf:thmS-carrier:core}
Assume the following eleven statements, each of the form ``for every $C_b$ there is $C_c$ such that for every leading-form
setting $(S,p,f,\chi,A_0)$ with constants $(C_b,C_c)$'':
(a) there are constants $X\ge0$ satisfying the curvature barriers such that the local wave condition holds for the Hida setting
$S_F$ with the local kernels of $F$, for every family $F$ of compactly supported coefficients; (b) the same for the local lift
condition, with some $0\le\Lambda\le\rho^L$; (c) the same for the local top condition, with the top constants and the top
thresholds, given three statements on the degree at the top, the types of the local vectors and the top ratio; (d) the germ
condition for the approximants $S_n$ below; (e) local positivity (H1) for $S_F$, $F$ compactly supported; (f) the local
integral formulas for $S_n$; (g) bounds $|\gamma^v(h\cdot(1,c),h'\cdot(1,c))|\le C(1+|c|)^{-3}$ and $|\widehat{\gamma^v}(l)|\le
C(1+|l|)^{-3}$ for the approximants, uniform in $n$; (h) positivity of the thresholds. Then the conclusion of
\cref{pf:thmS-carrier:builtbounds} holds.
\formallabel{Formalism.FourierDecayAlt.P.PBASM.core}
\formallabel{Formalism.FourierDecayAlt.P.PBASM.psd_of_tendsto}
\end{theorem}

\begin{proof}
\emph{Constants.} Let $c_1$ be the constant of \lref{Formalism.FourierDecayAlt.P.limit_of_module_forms} for $C_b$ (a
statement on the convergence of module forms under truncation, recalled below), let $C_4\ge1$ be the constant of
\cref{pf:thmS-carrier:sumform} (the coset-sum bound with constant $C_4$ from the local statements), let $c_2,\dots,c_8$ be the
constants of the hypotheses for $C_b$, and put
$C_c=\max(c_1,\dots,c_8,28e)$. Let $(S,p,f,\chi,A_0)$ be a leading-form setting with constants $(C_b,C_c)$; it is one with each
constant $c_i$.

\emph{Approximation.} Let $\beta\colon\R\to[0,1]$ be the smooth transition function ($0$ on $(-\infty,0]$, $1$ on $[1,\infty)$)
and $\beta_n(g)=\beta\bigl((4-\|g\|^2/(n+1)^2)/3\bigr)$ for the Frobenius norm, so $\beta_n=1$ for $\|g\|\le n+1$ and
$\beta_n=0$ for $\|g\|\ge2(n+1)$. Put $F_{n,v,\eps}=\beta_n\cdot\mathrm{cut}_\eps f_v$. By
\lref{Formalism.FourierDecayAlt.P.limit_of_module_forms}, for $\eps>0$ the $F_{n,v,\eps}$ are compactly supported
coefficients and, for every admissible local word of length $\le\thmScarrierNtop$, the local kernels of $F_{n,v,\eps}$ converge
to those of $\mathrm{cut}_\eps f_v$ as $n\to\infty$; after the normalisation of the component family, the local kernels
$RR^{(n)}_v$ of the approximating Hida setting $S_n$ converge to $RR_v$ for all admissible words and local points.

\emph{Local statements for $S_n$.} Hypotheses (e) and (f) give (H1) and the local integral formulas for $S_n$; (a) gives $X$
(extended by $0$ outside $1\le i\le r+1$) with the curvature barriers and the local wave condition; (b) gives $\Lambda$ with
$0\le\Lambda\le\rho^L$ and the local lift condition; (c) gives the top thresholds and the local top condition; (d) gives the
germ condition; (h) gives positivity of the thresholds.

\emph{From places to $S_n$.} With \cref{pf:thmS-carrier:blockschur}, the place-wise wave bound for $S_n$ follows by
\lref{Formalism.FourierDecayAlt.P.plane_wave_bound_of_local_bounds}, the lift bound below the top by
\lref{Formalism.FourierDecayAlt.P.PBA.liftOfLoc_statement}, the top diagonal bound by
\lref{Formalism.FourierDecayAlt.P.PBA.top_of_local}, and the coset-sum bound with constant $C_4$ by
\cref{pf:thmS-carrier:sumform}; each of the cited results states that the global bound follows from the corresponding local
conditions at all places (for the coset-sum bound, \cref{pf:thmS-carrier:sumform} proves this below).

\emph{Passage to the limit.} The lift bound below the top, the top diagonal bound and the coset-sum bound are non-strict
inequalities between finite sums of kernel values, hence closed under pointwise convergence of the local kernels
(\lref{Formalism.FourierDecayAlt.P.PBA.limit_closed}). The local integral formulas and the place-wise wave bound pass to
the limit by dominated convergence, using the uniform bounds (g)
(\lref{Formalism.FourierDecayAlt.P.integral_formula_passes_to_limits}). Finally, (H1) for $S$: each local Gram matrix of
$S$ is the entrywise limit of the PSD local Gram matrices of $S_n$, and a limit of PSD matrices is PSD, since Hermitian symmetry
and non-negativity of the quadratic forms pass to the limit.
\end{proof}

\subsubsection{The local statements for compactly supported coefficients}\label{sec:pf:thmS-carrier:leaves}

We now verify the hypotheses of \cref{pf:thmS-carrier:core}. For a place package $p$ we write $\mathrm{form}_j(w)$ for its
module form at level $j$ on vectors of local words $w$, $\mathrm{slv}_j(F,w,h)$ for the slice vectors of a coefficient $F$
at a local point $h$ of size $j$, $V_p(F)(q)$ for the local vector of $F$ at a local point $q$ of size $r+1$, and $\Theta$ for
its top lift. For a list $y=(y_1,\dots,y_K)$ of reals, $w^+_y$ and $w^-_y$ are the words obtained from $w$ by adding the letters
$(1,y_i)$, resp.\ $(0,y_i)$. A \emph{grid} of $K$ blocks of $d+1$ nodes is a family $(y_a(i))_{a<K,\,i\le d}$ of reals with
$i\mapsto y_a(i)$ injective; for $n=(n_a)_a$ put $y_n=(y_a(n_a))_a$ and
\[
\lambda_n=\prod_a\ell_a(n_a),\qquad\ell_a(i)=\prod_{i'\ne i}\bigl(y_a(i)-y_a(i')\bigr)^{-1}.
\]
For a polynomial $Q(y)=\sum_{m\le d}c_my^m$ and distinct nodes $y(0),\dots,y(d)$, Lagrange interpolation gives
$c_d=\sum_i\ell(i)Q(y(i))$; applied in each of $K$ variables in turn this extracts the coefficient of the top monomial.

\begin{theorem}[Built bounds from the leaves]\label{pf:thmS-carrier:builtleaves}
Assume the following statements.
\begin{enumerate}[label=(L\arabic*)]
\item (Radial head) For place packages of class $D^+$, compactly supported $F$, cut words $w$, local points $h,h'$ of size $j$
  and radial representatives $\Phi,\Phi'$ of the slice vectors of $F$ at $h,h'$, the radial pairing of $\Phi,\Phi'$ in the
  first radial variable at $t\ne0$ equals $\mathrm{lev}_{j+1}(t)\,RR_p(F)(j+1,w,h\cdot(1,t),h'\cdot(1,t))$.
\item (Head-top) The analogous identity at the top level $j=r$: the top density of radial representatives at a point $x$ of the
  radial domain equals $RR_p(F)(r+1,w,h\cdot(1,\zeta(x)),h'\cdot(1,\zeta(x)))$.
\item (Degree of the lift pairing) For compactly supported $F$ and grids of $K$ blocks of $2(j-s)+s+1$ nodes,
  \begin{multline*}
  \bigl((-1)^{s-1}(t^{-1}-t'^{-1})\bigr)^KRR_p(F)(j,w,h\cdot(s,t),h'\cdot(s,t'))\\
  =\sum_n\lambda_n\,\mathrm{pair}_{y_n}\bigl(\mathrm{slv}_j(F,w^+_{y_n},h),\mathrm{slv}_j(F,w^-_{y_n},h')\bigr),
  \end{multline*}
  with the pairing $\mathrm{pair}_y$ of $\mathrm{form}_j$ attached to $y$.
\item (Local top statement) The conclusion of \cref{pf:thmS-carrier:toploc}, given its three hypotheses.
\item (Degree at the top) The statement of \cref{pf:thmS-carrier:degtop}.
\item (Top ratio) For leading-form settings with $C_c\ge28e$, cut words $(1,y)\cdot w$ with $y$ a grid point,
  $\len(w)<\thmScarrierNtop$ and $y$ at distance at least $\thmScarrierNtop^2+1$ from the letters of $w$ (and from their
  translates by $-1$ for letters $(0,\cdot)$), and all indices $e,\rho$: the top ratio $R(w,e,\rho,y)$ is real and positive,
  and $|P_v(y)|\le e\,A_v(TT')^{-1/2}R(w,e,\rho,y)$, where $A_v$ is the sup constant of the place.
\item (Germ parts) For every $C_b$ there is $C_c$ such that the approximants $S_n$ of leading-form settings with constants
  $(C_b,C_c)$ satisfy the germ-fibre and germ-continuity conditions at every place.
\item (Uniform decay) \ref{pf:thmS-carrier:I1} implies hypothesis (g) of \cref{pf:thmS-carrier:core}.
\end{enumerate}
Then \ref{pf:thmS-carrier:I1} implies the conclusion of \cref{pf:thmS-carrier:builtbounds}.
\formallabel{Formalism.FourierDecayAlt.P.built_bounds_compact_of_leaves}
\formallabel{Formalism.FourierDecayAlt.P.PBASM.builtBounds2CptOfLeaves}
\formallabel{Formalism.FourierDecayAlt.P.PBASM.builtBounds2CptOfParts}
\formallabel{Formalism.FourierDecayAlt.P.PBASM.waveLocB_of_cpt}
\formallabel{Formalism.FourierDecayAlt.P.PBASM.liftLocB_of_cpt}
\end{theorem}

\begin{proof}
Assume \ref{pf:thmS-carrier:I1}. We verify hypotheses (a)--(h) of \cref{pf:thmS-carrier:core}.

\emph{Standard pairing decay.} By \lref{Formalism.FourierDecayAlt.P.StandardPairingDecay.pStdC_of}, \ref{pf:thmS-carrier:I1},
(L1) and (L2) imply the \emph{standard pairing decay}: for place packages of class $D^+$, compactly supported $F$, cut words $w$
and points $h,h'$ of size $j$, there is a continuous integrable $\varphi$ on $\C$ with $\varphi(c)=\mathrm{lev}_{j+1}(c)\,
RR_p(F)(j+1,w,h\cdot(1,c),h'\cdot(1,c))$ for $c\ne0$, whose Fourier transform is the plane-wave pairing of the slice vectors.

\emph{(f) and (g).} (g) is (L8). By \lref{Formalism.FourierDecayAlt.P.integral_bound_of_parts}, the standard pairing decay
and (g) imply the local integral formulas (f) for the approximants.

\emph{(a) Wave.} By \lref{Formalism.FourierDecayAlt.P.wave_local_positive_compact_support}, \ref{pf:thmS-carrier:I1} and
the standard pairing decay imply: for every leading-form setting with constants $(C_b,C_c)$, $C_c\ge28e$, there are constants
$X_{i,v}$ satisfying the curvature barriers such that the local wave condition with constants $X_{\cdot,v}$ holds for $S_F$ at
every place, for every family $F$ of compactly supported coefficients. So (a) holds with the constant $28e$.

\emph{(b) Lift.} By \lref{Formalism.FourierDecayAlt.P.lift_local_positive_compact_support}, \ref{pf:thmS-carrier:I1} and
(L3) give, for such settings, constants $\Lambda^0_{b,s,v}$ with $0\le\Lambda^0_{b,s,v}\le\rho^L_{b,s,v}$ for $1\le b$, $1\le s$,
$b+s\le r+1$, for which the local lift condition holds for $S_F$. Extending $\Lambda^0$ by $0$ outside this range gives
constants $\Lambda$ with $0\le\Lambda\le\rho^L$ and the local lift condition, so (b) holds with the constant $28e$.

\emph{(c) Top.} This is (L4), with the three statements it needs supplied by (L5), the vector-valued types statement
\lref{Formalism.FourierDecayAlt.P.vector_valued_types} (the type components of a local top vector sum to the vector), and
(L6).

\emph{(d) Germ.} By \lref{Formalism.FourierDecayAlt.P.PBASM.germLocB_of}, the local integral formulas (f) and (L7) imply
the germ condition for the approximants.

\emph{(e) Positivity.} By \lref{Formalism.FourierDecayAlt.P.positivity_bound}, \ref{pf:thmS-carrier:I1} implies (H1) for
$S_F$, $F$ compactly supported.

\emph{(h) Thresholds.} \Lref{Formalism.FourierDecayAlt.P.PBASM.thrPos} states that the thresholds of leading-form
settings are positive.

\Cref{pf:thmS-carrier:core} now gives the conclusion of \cref{pf:thmS-carrier:builtbounds}.
\end{proof}

\paragraph{The degree at the top.}
For a real place datum and $y\in\R$, the \emph{raising operator at the top} is $A_yG=D_{\mathrm{last}}^{-1}\,
\mathrm{AA}\bigl(y+\tfrac{2r+k}{2}\bigr)G$, where $\mathrm{AA}(u)$ is the Capelli-type determinant operator with parameter
$u$ and $D_{\mathrm{last}}$ the last leading minor attached to $(z,z')$.

\begin{lemma}[Polynomiality at the top]\label{pf:thmS-carrier:polytop}
Let $p$ be a place package, $G$ a compactly supported coefficient, $s\ge2$, and $q=h\cdot(s,t)$ a local point of size $r+1$;
let $d$ be the size of $h$. There are $c_0,\dots,c_d\in\Hom(W_q,W_p)$ with $c_d=((-1)^{s-1}/t)\,V_p(G)(q)$ and
$V_p(A_yG)(q)=\sum_{m\le d}y^mc_m$ for all $y\in\R$.
\formallabel{Formalism.FourierDecayAlt.P.polynomial_top}
\formallabel{Formalism.FourierDecayAlt.P.VandermondePolynomialIntegral.polyTop_of}
\formallabel{Formalism.FourierDecayAlt.P.VandermondePolynomialIntegral.poly_integral}
\end{lemma}

\begin{proof}
If $q$ is not a local point of size $r+1$ with positive weight $G_{\mathrm{pt}}(q)>0$, or if there is no frame element
$\gamma\in G_{k,r}$ with $\hat Z(\gamma)$ equal to the normal-form matrix of $q$, then $V_p(\cdot)(q)=0$ by definition and
$c=0$ works. Otherwise fix such a $\gamma$; then $G_f(z,z';\gamma)>0$. The operator $A_y$ preserves compact support
(\lref{Formalism.FourierDecayAlt.P.raising_compact_support}). By
\lref{Formalism.FourierDecayAlt.P.VandermondePolynomialIntegral.vecV_split}, for compactly supported $F$,
\[
V_p(F)(q)=\int a(x)\,\pi(x_1)\,K(F)\bigl(x_1^{-1}x_2^{-1}\gamma\bigr)\,dx,
\]
an integral over a product of the unipotent group $U_W(z)$ and a unipotent group $N$, with an integrable scalar weight $a(x)$
and the radial kernel $K(F)$ of $F$; the integrand is integrable. By
\lref{Formalism.FourierDecayAlt.P.VandermondePolynomialIntegral.pt_poly} (which uses that $K$ commutes with the operators
$\mathrm{AA}(u)$, \lref{Formalism.FourierDecayAlt.P.VandermondePolynomialIntegral.raising_lowering_relations}), for almost
every $x$ the integrand for $A_yG$ is a polynomial in $y$ of degree at most $d$ whose coefficient of $y^d$ is $(-1)^{s-1}/t$
times the integrand for $G$: at the point $x_1^{-1}x_2^{-1}\gamma$, which is in the top normal form, the Capelli-type
determinant divided by the last minor is a polynomial in its parameter with this leading coefficient.

It remains to integrate. Let $f(y,x)=\sum_{m\le d}y^mc_m(x)$ (almost every $x$), with $f(y,\cdot)$ integrable for each $y$ and
$c_d=\kappa g$ almost everywhere, $g$ integrable, $\kappa=(-1)^{s-1}/t$. Choose distinct $y_0,\dots,y_d$ and let $(v_{mi})$ be
the inverse of the Vandermonde matrix $(y_i^m)$; then $c_m(x)=\sum_iv_{mi}f(y_i,x)$ almost everywhere, so each $c_m$ is
integrable, and $\int f(y,x)\,dx=\sum_my^m\int c_m$, with $\int c_d=\kappa\int g=\kappa V_p(G)(q)$. Put $c_m:=\int c_m(x)\,dx$.
\end{proof}

\begin{theorem}[Degree at the top]\label{pf:thmS-carrier:degtop}
Let $p$ be a place package of class $D^+$ with constant $C_b$, $F$ a compactly supported coefficient, $s\ge2$, $w$ a local
word with an index $e$, and $q=h\cdot(s,t)$ a local point of size $r+1$; let $d$ be the size of $h$. Let $(y_a(i))$ be a grid
of $K$ blocks of $d+1$ nodes such that all words $w^+_{y_n}$ are cut words. Then
\begin{equation}\label{eq:pf:thmS-carrier:degtop}
\Bigl(\frac{(-1)^{s-1}}{t}\Bigr)^KV_p(\Theta_{w,e}F)(q)=\sum_n\lambda_n\,V_p\bigl(\Theta_{w^+_{y_n},e^+}F\bigr)(q),
\end{equation}
where $e^+$ is the index induced by $e$ on the longer word.
\formallabel{Formalism.FourierDecayAlt.P.degree_top}
\formallabel{Formalism.FourierDecayAlt.P.pb_degTop_of_polyTop}
\formallabel{Formalism.FourierDecayAlt.P.pb_degTop_of}
\formallabel{Formalism.FourierDecayAlt.P.degree_top_of_parts}
\end{theorem}

\begin{proof}
\emph{Compact representatives.} The word $w$ is a cut word, since the words $w^+_{y_n}$ are. By
\lref{Formalism.FourierDecayAlt.P.theta_compact_support} there is a compactly supported $F_0$ which has the same radial
representatives and the same local vectors as $\Theta_{w,e}F$. Let $\Phi$ be a radial representative of $F$; the top lift
maps it to a radial representative $\Phi_0$ of $\Theta_{w,e}F$, hence of $F_0$. For a list $y=(y_1,\dots,y_m)$ put $F_y=A_{y_m}\cdots A_{y_1}F_0$.
By \lref{Formalism.FourierDecayAlt.P.radial_raising_top} (the raising operator at the top preserves compact support and
maps a radial representative $\Phi$ of a coefficient to the radial representative $a(y)\Phi$ of its raising, where $a(y)$ is
the raising operator of the level datum), $F_y$ is compactly supported with radial representative $a(y_m)\cdots a(y_1)\Phi_0$.
By the word recursion for radial representatives (\lref{Formalism.FourierDecayAlt.P.PBN.rad_plusW}, which uses
\lref{Formalism.FourierDecayAlt.P.word_recursion}), the same radial function represents $\Theta_{w^+_y,e^+}F$ when $w^+_y$
is a cut word. Two compactly supported coefficients with the same radial representative have the same local vectors
(\lref{Formalism.FourierDecayAlt.P.density_constant}). Hence
\begin{equation}\label{eq:pf:thmS-carrier:vit}
V_p(F_y)(q)=V_p\bigl(\Theta_{w^+_y,e^+}F\bigr)(q)\qquad\text{whenever }w^+_y\text{ is a cut word}.
\end{equation}

\emph{Extraction.} For $0\le m\le K$ and $y\in\R^K$ let $W_m(y)=V_p(F_{(y_1,\dots,y_m)})(q)$. By \cref{pf:thmS-carrier:polytop}
applied to the compactly supported $G=F_{(y_1,\dots,y_m)}$, for each $m<K$ the function $t'\mapsto W_{m+1}(y_1,\dots,y_m,t')$ is
a polynomial of degree at most $d$ whose coefficient of $t'^d$ is $((-1)^{s-1}/t)\,W_m(y)$. By the $K$-fold Lagrange extraction
\lref{Formalism.FourierDecayAlt.P.extract_K_bound} (Lagrange interpolation in each variable, applied $K$ times; the blocks
are used in reverse order, which only permutes the sum),
\[
\Bigl(\frac{(-1)^{s-1}}{t}\Bigr)^KW_0=\sum_n\lambda_n\,W_K(y_n).
\]
Here $W_0=V_p(F_0)(q)=V_p(\Theta_{w,e}F)(q)$, and $W_K(y_n)=V_p(\Theta_{w^+_{y_n},e^+}F)(q)$ by \eqref{eq:pf:thmS-carrier:vit}.
\end{proof}

\paragraph{The local top statement.}
For kernels $\kappa,\kappa'$ on a finite set $Y$ write $\kappa\preceq\kappa'$ if $\kappa'-\kappa$ is PSD; for a family $u$ of
vectors of a Hilbert space indexed by $Y$, let $G[u](P,Q)=\langle u(P),u(Q)\rangle$ be its Gram kernel; it is PSD, since
$\sum_{P,Q}\bar x_Px_QG[u](P,Q)=\|\sum_P\bar x_Pu(P)\|^2$.

\begin{lemma}[Top kernel inequality]\label{pf:thmS-carrier:topkernel}
Let $Y$ be finite, let $E_a,R,N$ and $E_n$ ($n\in N$) be finite sets with injections $g_n\colon E_a\to E_n$, let
$V_a(e,\rho)$ and $V_n(e',\rho)$ be families of vectors indexed by $Y$, $d_a(e,\rho),d_n(e',\rho)\ge0$ weights, $\lambda_n\in
\R$, $P_n>0$, $c\in\C$, $C_0\ge0$ and $K\in\N$. Assume $V_a(e,\rho)=c\sum_n\lambda_nV_n(g_ne,\rho)$ on $Y$ and
$d_a(e,\rho)P_n\le C_0^Kd_n(g_ne,\rho)$. Then, with $\Lambda=\sum_n|\lambda_n|$,
\[
\sum_{e,\rho}d_a(e,\rho)\,G[V_a(e,\rho)]\preceq|c|^2\Lambda C_0^K\sum_n\frac{|\lambda_n|}{P_n}\sum_{e',\rho}d_n(e',\rho)\,
G[V_n(e',\rho)] .
\]
\formallabel{Formalism.FourierDecayAlt.P.PBL.top_kernel}
\formallabel{Formalism.FourierDecayAlt.P.PBL.loe_convex}
\end{lemma}

\begin{proof}
\emph{Convexity.} For vectors $u_n(P)$ and $x\in\C^Y$, the quadratic form of $G[\sum_n\lambda_nu_n]$ at $x$ is
$\|\sum_n\lambda_nU_n\|^2$ with $U_n=\sum_P\bar x_Pu_n(P)$ (the $\lambda_n$ are real), and by the Cauchy--Schwarz
inequality $\|\sum_n\lambda_nU_n\|^2\le(\sum_n|\lambda_n|)\sum_n|\lambda_n|\|U_n\|^2$. Hence
$G[\sum_n\lambda_nu_n]\preceq\Lambda\sum_n|\lambda_n|G[u_n]$. Applied to $u_n=V_n(g_ne,\rho)$ and multiplied by
$d_a(e,\rho)|c|^2\ge0$:
$d_a(e,\rho)G[V_a(e,\rho)]\preceq|c|^2\Lambda\sum_n|\lambda_n|\,d_a(e,\rho)G[V_n(g_ne,\rho)]$.

\emph{Weights.} Since $G[V_n(g_ne,\rho)]$ is PSD and $d_a(e,\rho)\le C_0^Kd_n(g_ne,\rho)/P_n$,
$d_a(e,\rho)G[V_n(g_ne,\rho)]\preceq(C_0^K/P_n)\,d_n(g_ne,\rho)G[V_n(g_ne,\rho)]$.

\emph{Summation.} Summing over $(e,\rho)$ and using that $g_n$ is injective, $\sum_{e,\rho}d_n(g_ne,\rho)G[V_n(g_ne,\rho)]$ is
a sub-sum of $\sum_{e',\rho}d_n(e',\rho)G[V_n(e',\rho)]$, the omitted terms being PSD. Combining the three steps gives the
claim.
\end{proof}

\begin{theorem}[Local top statement at built places]\label{pf:thmS-carrier:toploc}
Assume the degree at the top (the statement of \cref{pf:thmS-carrier:degtop}), the vector-valued types statement and the
top ratio (L6). Let $C_b,C_c$ with $C_c\ge28e$ and let $(S,p,f,\chi,A_0)$ be a leading-form setting with constants
$(C_b,C_c)$. For $1\le b\le r+1$ define the \emph{top constants} $\lambda^{\mathrm{top}}_{b,v}\ge0$ by
\[
(\lambda^{\mathrm{top}}_{b,v})^2=\Bigl(\frac g{2\pi}\Bigr)^{2\nu}\Bigl(\frac{28\thmScarrierNtop}{\beta_v-\alpha_v}
\Bigr)^{2(b-1)}\,\frac{e\,A_v}{\sqrt{TT'}},
\]
with $\nu=\nu(b,r+2-b)$ the block exponent and $A_v$ the sup constant of the place. Then:
\begin{enumerate}[label=(\roman*)]
\item $0\le\lambda^{\mathrm{top}}_{r+2-s,v}$ and $2\lambda^{\mathrm{top}}_{r+2-s,v}\le\rho^L_{r+2-s,s,v}$ for $2\le s\le r+1$;
\item for every family $F$ of compactly supported coefficients and every place $v$, the local top condition with constants
  $\lambda^{\mathrm{top}}$ holds for $S_F$: for $2\le s\le r+1$ and every admissible local word $a$ with
  $\len(a)+5\le\thmScarrierNtop$ there is a probability vector $\theta$ on local words such that for all $\eps>0$, $t\ne0$ and
  finite sets $Y$ of points of size $r+1-s$,
  \[
  \tilde\gamma_{\eps,r+1,a}\bigl(P\cdot(s,t),Q\cdot(s,t)\bigr)\preceq(|t|\lambda^{\mathrm{top}}_{r+2-s,v})^{10}\sum_{a'}
  \theta(a')\,\tilde\gamma_{\eps,r+1,a'}\bigl(P\cdot(s,t),Q\cdot(s,t)\bigr)\quad\text{on }Y.
  \]
\end{enumerate}
\formallabel{Formalism.FourierDecayAlt.P.top_local_built_places}
\formallabel{Formalism.FourierDecayAlt.P.PBL.topLocB}
\end{theorem}

\begin{proof}
(i) is \lref{Formalism.FourierDecayAlt.P.PBL.lamTop_thr}, a comparison of explicit constants which uses $C_c\ge28e$.

(ii) Fix $v$, $s$ and $a$; let $d=r+1-s$ and $K=5$. By \lref{Formalism.FourierDecayAlt.P.new_letters_admissible} there is a
grid $(y_b(i))$ of $K$ blocks of $d+1$ nodes, taken from the grids of $S$ at $v$, such that every word $a^+_{y_n}$ is an
admissible cut word, $P_v$ does not vanish at the nodes, $\sum_i|\ell_b(i)|\le B:=(28\thmScarrierNtop/(\beta_v-\alpha_v))^d$
for each block, and $\theta(a^+_{y_n})=|\lambda_n|/\Lambda$ (summed over the $n$ giving the same word) is a probability vector,
where $\Lambda=\sum_n|\lambda_n|>0$. Expanding the product over the blocks, $\Lambda=\prod_b\sum_i|\ell_b(i)|\le B^K$.

\emph{Factorisation at the top.} By \lref{Formalism.FourierDecayAlt.P.PBL.gN_top}, which uses the vector-valued types
statement, for every local word $a'$ and points $P,Q$ of size $d$,
\[
\tilde\gamma_{\eps,r+1,a'}\bigl(P\cdot(s,t),Q\cdot(s,t)\bigr)=\frac{\omega(P)\omega(Q)}{\Gamma\,\mathrm{size}_v(a')}\sum_{e,\rho}
\delta(a',e,\rho)\,G[u_{a',e,\rho}](P,Q),
\]
where $u_{a',e,\rho}(P)$ is the $\rho$-component of the local top vector of $\Theta_{a',e}F_{v,\eps}$ at the standardised point
of $P\cdot(s,t)$, the \emph{second defects} $\delta(a',e,\rho)$ are $\ge0$, and $\omega(P)>0$, $\Gamma>0$ are explicit.

\emph{Degree.} With $\tau=(g/2\pi)^{\nu}i^{s-1}t$ (the last coordinate of the standardised point) and
$c=((-1)^{s-1}/\tau)^{-K}$, \lref{Formalism.FourierDecayAlt.P.PBL.vecT_deg} (the degree at the top at the standardised
point) gives $u_{a,e,\rho}=c\sum_n\lambda_n\,u_{a^+_{y_n},e^+,\rho}$ on points of size $d$.

\emph{Weights.} By \lref{Formalism.FourierDecayAlt.P.PBL.second_defect_plus_word_le}, whose hypothesis is supplied by the
top ratio through \lref{Formalism.FourierDecayAlt.P.PBL.rhyp_nodes},
$\delta(a,e,\rho)P_n\le C_0^K\delta(a^+_{y_n},e^+,\rho)$ with $P_n=\prod_b|P_v(y_b(n_b))|>0$ and $C_0=eA_v(TT')^{-1/2}$.

\emph{Loewner bound.} \Cref{pf:thmS-carrier:topkernel} (with $g_n(e)=e^+$, injective) gives
$\sum\delta(a)G[u_a]\preceq|c|^2\Lambda C_0^K\sum_n(|\lambda_n|/P_n)\sum\delta(a^+_{y_n})G[u_{a^+_{y_n}}]$. Multiply both sides
by $\omega(P)\omega(Q)/(\Gamma\,\mathrm{size}_v(a))$ (a congruence by a positive diagonal matrix and a positive factor, which
preserves $\preceq$) and use $\mathrm{size}_v(a^+_{y_n})=P_n\,\mathrm{size}_v(a)$:
\[
\tilde\gamma_{\eps,r+1,a}\preceq|c|^2\Lambda^2C_0^K\,\mathrm{mix},\qquad\mathrm{mix}=\sum_n\frac{|\lambda_n|}{\Lambda}\,
\tilde\gamma_{\eps,r+1,a^+_{y_n}},
\]
at the points $P\cdot(s,t),Q\cdot(s,t)$; $\mathrm{mix}$ is the $\theta$-average and is PSD, by the factorisation above.

\emph{Constants.} $|c|=|\tau|^K=((g/2\pi)^\nu|t|)^K$, $\Lambda\le B^K$, and $(\lambda^{\mathrm{top}}_{r+2-s,v})^2=
(g/2\pi)^{2\nu}B^2C_0$. Hence
$|c|^2\Lambda^2C_0^K\le\bigl(|t|^2(g/2\pi)^{2\nu}B^2C_0\bigr)^K=(|t|\lambda^{\mathrm{top}}_{r+2-s,v})^{2K}$, and since
$\mathrm{mix}$ is PSD, $|c|^2\Lambda^2C_0^K\,\mathrm{mix}\preceq(|t|\lambda^{\mathrm{top}}_{r+2-s,v})^{10}\,\mathrm{mix}$.
Transitivity of $\preceq$ gives (ii).
\end{proof}

\paragraph{Uniform decay.}
Let $\psi(s,t)=\exp(st-\bar s\bar t)$ for $s,t\in\C$, a unitary character in $t$; then $\psi(s,t)\psi(\sigma,t)=\psi(s+\sigma,t)$.
For an integrable $\varphi$ on $\C$ put $\hat\varphi(s)=\int_\C\psi(s,t)\varphi(t)\,dt$, and let
$I_3=\int_\C(1+|s|)^{-3}\,ds<\infty$. For a place package $p$, a level $j$ and a word $w$, let $\mathrm{wave}(s)$ be the
plane-wave operator on vectors of $w$, let $\mathrm{wave}'(\sigma)$ be the plane-wave operator with parameter $\sigma$ acting on
the letters of $w$, and $\mathrm{wave}^{(3)}_\sigma=(\mathrm{wave}'(\sigma)-1)^3$.

\begin{theorem}[Uniform decay from the heads]\label{pf:thmS-carrier:unifdecay}
Assume the radial head and head-top statements (L1), (L2). Then \ref{pf:thmS-carrier:I1} implies: for every $C_b$ there is $C_c$
such that for every parameter, place package $p$ of class $D^+$ with constant $C_b$, cut-off $\chi$, $\eps>0$, smooth $f$ with
all growth bounds with constant $A\ge C_c(r+1)^4\bigl((N+3)(2r+k+1)^2+(2r+k+1)^3\bigr)$, cut interval $[\mathrm{lo},\mathrm{hi}]$
on the cut with $\mathrm{hi}-\mathrm{lo}\ge10(N+3(2r+1)+1)$, level $j\le r$, cut word $w$ with $\len(w)\le N$ and local points
$h,h'$ of size $j$, there is $C$ such that for all $n$, with $F_n=\beta_n\cdot\mathrm{cut}_\eps f$,
\begin{enumerate}[label=(\roman*)]
\item $|\mathrm{lev}_{j+1}(c)\,RR_p(F_n)(j+1,w,h\cdot(1,c),h'\cdot(1,c))|\le C(1+|c|)^{-3}$ for $c\ne0$, and
\item $|\mathrm{form}_j(w)\bigl(\mathrm{wave}(s)\,\mathrm{slv}_j(F_n,w,h),\mathrm{slv}_j(F_n,w,h')\bigr)|\le C(1+|s|)^{-3}$ for
  $s\in\C$.
\end{enumerate}
\formallabel{Formalism.FourierDecayAlt.P.uniform_decay_of_heads}
\formallabel{Formalism.FourierDecayAlt.P.PBO.decayUnifOfHeads}
\formallabel{Formalism.FourierDecayAlt.P.PBO.decayUnif_of}
\formallabel{Formalism.FourierDecayAlt.P.PBO.transform_diff3}
\formallabel{Formalism.FourierDecayAlt.P.PBO.fsv_sigma}
\end{theorem}

\begin{proof}
\emph{Constants.} Let $C_F$ be the constant of \lref{Formalism.FourierDecayAlt.P.limit_of_module_forms} for $C_b$ (for
$A\ge C_F(r+1)^4((N+3)(2r+k+1)^2+(2r+k+1)^3)$ the $F_n$ are compactly supported, and for cut words of length $\le N+3$ the
module forms of the $F_n$ converge as $n\to\infty$), and $C_W$ the constant of the wave-difference statement
\lref{Formalism.FourierDecayAlt.P.wave_difference_form}: for cut words $w^+_y$ with $\len(y)\le3$ there is $C_y$ such that
for all $n$ and $|\sigma|\le2$
\begin{equation}\label{eq:pf:thmS-carrier:wavediff}
\bigl|\mathrm{form}_j(w^+_y)\bigl(\mathrm{wave}^{(3)}_\sigma x_n,\mathrm{wave}^{(3)}_\sigma x_n\bigr)\bigr|\le C_y|\sigma|^6,
\end{equation}
where $x_n=\mathrm{slv}_j(F_n,w^+_y,h)$ and $\mathrm{wave}^{(3)}_\sigma$ acts on the letters of $w$. Put $C_c=\max(C_F,C_W)$.

\emph{Uniform bounds.} By \lref{Formalism.FourierDecayAlt.P.PBO.grid_exists} there is a grid of $3$ blocks of $2j+1$ nodes in
$[\mathrm{lo},\mathrm{hi}]$ such that all words $w^\pm_{y_n}$ are cut words (of length $\le N+3$) whose letters are at distance
$>2$. A convergent sequence is bounded, so there is $B\ge0$ bounding, for all $n$, the diagonal values
$|RR_p(F_n)(j,w',h_0,h_0)|=|\mathrm{form}_j(w')(\mathrm{slv}_j(F_n,w',h_0),\mathrm{slv}_j(F_n,w',h_0))|$ for
$w'\in\{w\}\cup\{w^\pm_{y_n}\}$ and $h_0\in\{h,h'\}$. Let
\[
K_0=\sum_n|\lambda_n|\,\Bigl(\prod_aq_w(y_a(n_a))\Bigr)^{-1/2},\qquad T_1=8(1+K_0)B,
\]
with the positive weights $q_w$ of the cut word $w$, and let $C_0'$, $C_n'$ be the constants of \eqref{eq:pf:thmS-carrier:wavediff}
for $y=\emptyset$ and $y=y_n$; put $C_W'=|C_0'|+\sum_n|C_n'|$.

\emph{(ii).} \Lref{Formalism.FourierDecayAlt.P.PBO.xi_bound} states: for a compactly supported coefficient, a linear map
$M$ on vectors of $w$ preserving compact support and niceness, and $a,b\ge0$ bounding the diagonal module forms of $Mx$ and of
the vectors of the plus words (for $M$ applied to the old letters), resp.\ of $y$ and of the vectors of the minus words,
\[
\bigl|\mathrm{form}_j(w)(\mathrm{wave}(s)Mx,y)\bigr|\le8(1+K_0)\sqrt{ab}\,(1+|s|)^{-3}\qquad(s\in\C).
\]
With $M=1$ and $a=b=B$ this gives (ii) with the constant $T_1$.

\emph{Fourier inversion.} By \lref{Formalism.FourierDecayAlt.P.PBO.pStdCpt_of_heads}, \ref{pf:thmS-carrier:I1}, (L1) and
(L2) give a continuous integrable $\varphi$ on $\C$ with $\varphi(c)=\mathrm{lev}_{j+1}(c)RR_p(F_n)(j+1,w,h\cdot(1,c),
h'\cdot(1,c))$ for $c\ne0$ and $\hat\varphi(s)=\mathrm{form}_j(w)(\mathrm{wave}(s)\,\mathrm{slv}_j(F_n,w,h),
\mathrm{slv}_j(F_n,w,h'))$. By (ii), $\hat\varphi$ is continuous and bounded by $T_1(1+|s|)^{-3}$, which is integrable on $\C=\R^2$; so $\hat\varphi$ is
integrable with $\int|\hat\varphi|\le T_1I_3$, and by Fourier inversion
(\lref{Formalism.FourierDecayAlt.P.PBO.norm_le_integral_fsv}: $|\varphi(t)|\le\int|\hat\varphi|$ for continuous integrable
$\varphi$ with integrable $\hat\varphi$), $|\varphi(c)|\le T_1I_3$. For $|c|\le1$ this gives $(1+|c|)^3|\varphi(c)|\le8T_1I_3$.

\emph{Large $|c|$.} Let $|c|>1$ and $\sigma=i\pi/(2c)$. Then $\sigma c-\bar\sigma\bar c=i\pi$, so $\psi(\sigma,c)=-1$, and
$|\sigma|=\pi/(2|c|)\le2$. Let $\phi(t)=(\psi(\sigma,t)-1)^3\varphi(t)$, continuous and integrable since $|\psi(\sigma,t)-1|\le2$.
By the multiplicativity of $\psi$,
\[
\hat\phi(s)=\hat\varphi(s+3\sigma)-3\hat\varphi(s+2\sigma)+3\hat\varphi(s+\sigma)-\hat\varphi(s),
\]
and by \lref{Formalism.FourierDecayAlt.P.PBO.form_wave_mDiff} (for compactly supported vectors the third difference in
the wave parameter is produced by $\mathrm{wave}^{(3)}_\sigma$) the right side equals
\[
\mathrm{form}_j(w)\bigl(\mathrm{wave}(s)\,\mathrm{wave}^{(3)}_\sigma\,\mathrm{slv}_j(F_n,w,h),\mathrm{slv}_j(F_n,w,h')\bigr).
\]
Apply the bound of \lref{Formalism.FourierDecayAlt.P.PBO.xi_bound} with $M=\mathrm{wave}^{(3)}_\sigma$ (which preserves niceness and
compact support), $a=C_W'|\sigma|^6$ by \eqref{eq:pf:thmS-carrier:wavediff}, and $b=B$:
\[
|\hat\phi(s)|\le8(1+K_0)|\sigma|^3\sqrt{C_W'B}\,(1+|s|)^{-3}.
\]
As before, $|\phi(c)|\le\int|\hat\phi|\le8(1+K_0)|\sigma|^3\sqrt{C_W'B}\,I_3$. Since $|\phi(c)|=|(-2)^3\varphi(c)|=8|\varphi(c)|$,
$|\varphi(c)|\le(1+K_0)\sqrt{C_W'B}\,I_3|\sigma|^3$. Finally $(1+|c|)^3\le8|c|^3$ and $8|c|^3|\sigma|^3=\pi^3$, so
$(1+|c|)^3|\varphi(c)|\le(1+K_0)\sqrt{C_W'B}\,I_3\pi^3$.

Hence (i) and (ii) hold for all $n$ with $C=T_1+8T_1I_3+(1+K_0)\sqrt{C_W'B}\,I_3\pi^3$.
\end{proof}

\paragraph{The germ rate.}
For a Hida setting $S$ let $\beta=4\pi/g$, let $\psi_\beta$ be the character of $\C$ attached to $\beta$, and let
$\mathrm{dl}_{b,s}=\mathrm{rd}(E)^{2\nu(b,s)}/\mathrm{de}(g)>0$ be the local covolume constants. For $D\ne0$ let $R_s(D)$ be the
set of $s$-th roots of $D$. We say that the \emph{germ rate} holds at $v$ if for all $\eps>0$, $b\ge1$, $s\ge2$ with
$b+s\le r+2$, admissible $a$, every Hilbert space $H$ and every map $\mathrm{Vec}$ from local points to $H$ with
$\tilde\gamma_{\eps,b+s-1,a}(P,Q)=\langle\mathrm{Vec}(Q),\mathrm{Vec}(P)\rangle$ for all $P,Q$, there are $\eps'\in\C$ with
$|\eps'|=1$ and $\phi\colon\C\to\C$ with $|\phi|=1$, continuous on $\C\smallsetminus\{0\}$, such that for every local point $h$
of size $b-1$ and every compact $K\subset\C\smallsetminus\{0\}$ there is $C$ with
\begin{equation}\label{eq:pf:thmS-carrier:germrate}
\bigl\|U(h,x,D)-M(h,x,D)\bigr\|\le C|x|^{-1}\qquad(x\ne0,\ D\in K),
\end{equation}
where
\begin{gather*}
U(h,x,D)=\frac{|x|^{s-1}}{\sqrt{\mathrm{dl}_{b+s-1,1}}}\,\mathrm{Vec}\bigl(h\cdot(s-1,x^{-1})\cdot(1,x^{s-1}D)\bigr),\\
M(h,x,D)=\sum_{z\in R_s(\eps'D)}\psi_\beta(szx)\,\phi(z)\,\frac{\mathrm{Vec}(h\cdot(s,z))}{s|z|^{s-1}\sqrt{\mathrm{dl}_{b,s}}} .
\end{gather*}

\begin{theorem}[Germ rate]\label{pf:thmS-carrier:germrate}
\ref{pf:thmS-carrier:I1} implies: for every $C_b$ there is $C_c$ such that for every leading-form setting with constants
$(C_b,C_c)$, all $n$ and all places $v$, the germ rate holds at $v$ for the approximant $S_n$.
\formallabel{Formalism.FourierDecayAlt.P.germ_rate_bound}
\formallabel{Formalism.FourierDecayAlt.P.germ_rate_bound_of_split}
\formallabel{Formalism.FourierDecayAlt.P.GermRateSplit.germRateB_of}
\formallabel{Formalism.FourierDecayAlt.P.GermRateSplit.germRate_of}
\formallabel{Formalism.FourierDecayAlt.P.germ_rate_top_bound}
\formallabel{Formalism.FourierDecayAlt.P.GermRateSplit.germ_sum}
\end{theorem}

\begin{proof}
Assume \ref{pf:thmS-carrier:I1} and let $C_b$ be given. Let $c_1$ be the constant of
\lref{Formalism.FourierDecayAlt.P.germ_rate_lower_bound}, which proves the germ rate for the levels $b+s\le r+1$, and $c_2$
the constant of \lref{Formalism.FourierDecayAlt.P.GermFibreKernels.cutSeq_cpt}, which states that for leading-form
settings with constant $\ge c_2$ the approximants have compactly supported coefficients. Put $C_c=\max(c_1,c_2)$; a
leading-form setting with constant $C_c$ is one with each smaller constant, since the conditions involving the constant
persist when it is decreased. If $b+s\le r+1$ we are done; let $b+s=r+2$, so
$b+s-1=r+1$. Take $\eps'=1$ and $\phi=1$.

\emph{An identity.} For $D\ne0$ the roots $z\in R_s(D)$ are non-zero. Writing out the definitions,
\[
U(h,x,D)-M(h,x,D)=\sum_{j\in\{\ast\}\sqcup R_s(D)}\mathrm{co}_j\,\mathrm{Vec}(P_j),
\]
with $P_\ast=h\cdot(s-1,x^{-1})\cdot(1,x^{s-1}D)$, $P_z=h\cdot(s,z)$, $\mathrm{co}_\ast=|x|^{s-1}\mathrm{dl}_{r+1,1}^{-1/2}$ and
$\mathrm{co}_z=-\psi_\beta(szx)(s|z|^{s-1})^{-1}\mathrm{dl}_{b,s}^{-1/2}$.

\emph{Gram structure at the top.} By \lref{Formalism.FourierDecayAlt.P.GermRateSplit.gN_top_gram} there are finitely many
$c_i\ge0$, compactly supported coefficients $F_i$ and bounded linear maps $L_i$ into a fixed Hilbert space such that
\[
\tilde\gamma_{\eps,r+1,a}(P,Q)=\sum_ic_i\bigl\langle\omega(P)L_iV(F_i)(\mathrm{std}\,P),\ \omega(Q)L_iV(F_i)(\mathrm{std}\,Q)
\bigr\rangle
\]
for all $P,Q$, with the explicit weights $\omega$ and the standardisation of points $\mathrm{std}$.

\emph{Comparison of norms.} Let $Y_i(P)=\omega(P)L_iV(F_i)(\mathrm{std}\,P)$ and $a_j=\mathrm{co}_j$. Then
\begin{align*}
\Bigl\|\sum_ja_j\mathrm{Vec}(P_j)\Bigr\|^2&=\sum_{j,l}a_j\overline{a_l}\,\tilde\gamma_{\eps,r+1,a}(P_l,P_j)
=\sum_ic_i\Bigl\langle\sum_l\overline{a_l}Y_i(P_l),\sum_j\overline{a_j}Y_i(P_j)\Bigr\rangle\\
&=\sum_ic_i\Bigl\|\sum_j\overline{a_j}\,Y_i(P_j)\Bigr\|^2,
\end{align*}
and since $\sqrt{\sum_ic_iX_i^2}\le
\sum_i\sqrt{c_i}X_i$ for $X_i\ge0$,
\[
\|U-M\|\le\sum_i\sqrt{c_i}\,\|L_i\|\,\Bigl\|\sum_j\overline{\mathrm{co}_j}\,\omega(P_j)V(F_i)(\mathrm{std}\,P_j)\Bigr\| .
\]

\emph{The single-coefficient estimate.} By \lref{Formalism.FourierDecayAlt.P.GermRateSplit.top_core}, for every compactly
supported coefficient $F'$, every $h$ of size $b-1$ and compact $K\not\ni0$ there is $C'$ with
$\|\sum_j\overline{\mathrm{co}_j}\,\omega(P_j)V(F')(\mathrm{std}\,P_j)\|\le C'|x|^{-1}$ for $x\ne0$, $D\in K$. Applying this to
the $F_i$ with constants $C_i$ gives \eqref{eq:pf:thmS-carrier:germrate} with $C=\sum_i\sqrt{c_i}\|L_i\|C_i$.
\end{proof}

\paragraph{The coset-sum bound.}
Recall the coset-sum form (\cref{pf:thmS-carrier:corT}): for $\eps>0$, $2\le s\le i\le r+1$ and $w\in W_{\mathrm{all}}$ with
$\thmScarrierCnt(w)+(s+1)\cdot5\le\thmScarrierNtop$ and $V_{i-s+1,s}\le1$, the coset-sum form of $S_{\mathrm{base}}$ with a
constant $M$ asserts a probability vector $\theta$ on $\thmScarrierWb(\thmScarrierCnt(w)+5s)$ such that for all $Y$ of level
$i-s$, PSD $\kappa$, $t_0$ and finite sets $\Sigma$ of good points in $t_0+\Lambda_{i-s+1,s}$,
\[
\sum_{t\in\Sigma}\Re\thmScarrierNst_{\eps,i,s,w}(Y,\kappa;t,t)\le M\sum_u\theta(u)\,\thmScarrierNz_{\eps,i-s,u}(Y,\kappa).
\]
For $\Lambda\in\R^f_{\ge0}$ let $m_\Lambda\ge0$ be the explicit majorant on $E_\infty$ with exponent $5$ attached to
$\Lambda$, and for a lattice $\mathcal L$ let $\mathcal L^\vee$ be its dual.

\begin{theorem}[Coset-sum bound from local statements]\label{pf:thmS-carrier:sumform}
There is $C_4\ge1$ such that every Hida setting $S$ with $g_1,g_2>0$ and $E$ unramified over $F$ at all finite primes which
satisfies (H1), (H2), the local integral formulas with constants $\mathrm{dl}_{i,1}$ at every place, the local wave condition at
every place with constants $X$ satisfying the curvature barriers, the local lift condition at every place with constants
$\Lambda$ satisfying $0\le\Lambda_{b,s,v}\le\rho^L_{b,s,v}$ for $b+s\le r+1$, positivity of the thresholds, and the germ
condition at every place, satisfies the coset-sum bound with constant $C_4$.
\formallabel{Formalism.FourierDecayAlt.P.sum_form_of_local_statements}
\formallabel{Formalism.FourierDecayAlt.P.PBC.sumFormOfLoc}
\formallabel{Formalism.FourierDecayAlt.P.PBC.sumFormConcrete_of}
\end{theorem}

\begin{proof}
\emph{Three uniform constants.} (1) By \lref{Formalism.FourierDecayAlt.P.PBC.pbc_lstdSum} there is $C_s\ge1$ such that for
the standard lattices $\mathcal L(R)$ ($R>0$), $g>0$ and $\Lambda\ge0$ the dual sum $\sum_{\sigma\in\mathcal L(R)^\vee}
m_\Lambda(g\sigma)$ converges and is at most $\prod_vC_s(1+(R\Lambda_v/g)^2)$. (2) By
\lref{Formalism.FourierDecayAlt.P.PBC.pbc_sizesSup} there is $B\ge1$ such that for all $b,n\ge1$ there is $R>0$ with
$\prod_v\mathrm{dl}_{b,n}\prod_vC_s(1+(R\rho_{b,n,v}/g)^2)/R^{2f}\le B^fV_{b,n}$, where $\rho_{b,n,v}$ are the thresholds.
(3) By \lref{Formalism.FourierDecayAlt.P.PBC.pbc_blockLatV} there is $A\ge1$ such that the block lattices
$\Lambda_{b,s}\subset E_\infty$ are lattices of covolume $\prod_v\mathrm{dl}_{b,s}$ whose duals, scaled by $g$, lie in a lattice,
and such that $\sum_{\sigma\in\Lambda_{b,s}^\vee}m_\Lambda(g\sigma)$ converges and is at most $A^f$ whenever $0\le\Lambda_v\le
\rho_{b,s,v}$ and $V_{b,s}\le1$. Put $C_4=\max(A,B)$.

\emph{The germ supremum.} Let $\eps$, $i$, $s$, $w$ be as in the coset-sum form and $b=i-s+1$. By
\lref{Formalism.FourierDecayAlt.P.PBC.gsup_all}, proved by induction on the level from the local statements and (1),
(2), there is $P'$ with $0\le P'\le B^{f(s-1)}\prod_{1\le m<s}V_{b,m}$ such that the diagonal terms at
level $(b,s-1)$ are bounded by $P'$ times $\theta$-averages of $\thmScarrierNz$ (the germ-supremum bound).

\emph{The germ step.} Let $M=\sum_{\sigma\in\Lambda_{b,s}^\vee}m_{\Lambda_{b,s-1}}(g\sigma)$; by (3), applied with
$\Lambda_{b,s-1,v}\le\rho_{b,s,v}$, $M$ converges and $M\le A^f$. By
\lref{Formalism.FourierDecayAlt.P.PBC.germ_step}, which combines the lift bound below the top at all places, the germ
condition (through a positivity argument of Herglotz type for the germs) and summation over the dual lattice of
$\Lambda_{b,s}$, there is a probability vector $\theta$ on
$\thmScarrierWb(\thmScarrierCnt(w)+5s)$ such that for all $Y,\kappa,t_0,\Sigma$ as above
\[
\covol(\Lambda_{b,s})\sum_{t\in\Sigma}\Re\thmScarrierNst(t,t)\le\prod_v\mathrm{dl}_{b,s}\cdot MP'\sum_u\theta(u)
\thmScarrierNz_{\eps,b-1,u}(Y,\kappa).
\]
Since $\covol(\Lambda_{b,s})=\prod_v\mathrm{dl}_{b,s}>0$, the coset-sum form holds with the constant $MP'$, and
$MP'\le A^fB^{f(s-1)}\prod_{1\le m<s}V_{b,m}\le C_4^{fs}\prod_{1\le m<s}V_{b,m}$.
\end{proof}

\begin{proof}[Proof of \cref{pf:thmS-carrier:builtbounds}]
The formal proof reduces the theorem in two steps to \cref{pf:thmS-carrier:builtleaves}: the top truncation of averages
\lref{Formalism.FourierDecayAlt.P.top_truncation_of_averages} implies the head-top statement (L2)
(\lref{Formalism.FourierDecayAlt.P.pb_headTop_of_mkTop}), and then the hypotheses of \cref{pf:thmS-carrier:builtleaves}
are supplied as follows: (L1) is the radial head \lref{Formalism.FourierDecayAlt.P.radial_head}; (L3) is the degree of the
lift pairing \lref{Formalism.FourierDecayAlt.P.degree_of_lift_pairing}; (L4) is \cref{pf:thmS-carrier:toploc}; (L5) is
\cref{pf:thmS-carrier:degtop}; (L6) is the top ratio \lref{Formalism.FourierDecayAlt.P.top_ratio}; (L7) is
\lref{Formalism.FourierDecayAlt.P.germ_parts_compact_support}, which derives the germ parts from the continuity of the
kernels at built places \lref{Formalism.FourierDecayAlt.P.kernel_continuity_built}, the germ rate
(\cref{pf:thmS-carrier:germrate}) and \ref{pf:thmS-carrier:I1}; and (L8) is
\lref{Formalism.FourierDecayAlt.P.uniform_decay_of_place}, which derives the uniform decay at all places of the
approximants from \ref{pf:thmS-carrier:I1}, (L1), (L2) and the single-place uniform decay of \cref{pf:thmS-carrier:unifdecay}.
\Cref{pf:thmS-carrier:builtleaves} then gives the theorem.
\end{proof}

% ---------------------------------------------------------------------------
\subsection{Model places}\label{sec:pf:thmS-carrier:model}

\begin{theorem}[Model places]\label{pf:thmS-carrier:modelplace}
Statement \ref{pf:thmS-carrier:I2} holds.
\formallabel{Formalism.FourierDecayAlt.P.CutMP.SmoothTransferRegularity.model_place_built_holds}
\end{theorem}

At a built place $v$ the energy $\mathcal E_v(a;\mathrm{cut}_\eps f_v)$ of a local word $a$ is a quadratic expression in the
Bessel transform $J_{\mathrm{cut}_\eps f_v}$. The theorem identifies it with the cut-off energy form of a model function.
The link is a \emph{gauge transfer} $\mathrm{tr}_e$: a frame $e$ of the fibres of the radial domain turns functions on
$\mathrm{Mat}_{k,r}$ into families of functions on $Q_{r,T_T}$, and the transfer intertwines the Casimir-type operators
$\mathrm{CV}(Y)$ with the model operators $C_0(Y)$. We set up the transfer, deduce the theorem from six \emph{transfer
statements} (\cref{pf:thmS-carrier:openQQ}), and then prove three of the transfer statements from deeper local identities
(\cref{pf:thmS-carrier:hatreal,pf:thmS-carrier:startransfer,pf:thmS-carrier:polyradial}). The proof of
\cref{pf:thmS-carrier:modelplace} is given at the end of this subsection.

\subsubsection{Frames, the gauge transfer and the model conditions}

\begin{definition}[Frames and the gauge transfer]\label{pf:thmS-carrier:transferdef}
Let $S$ be a Hida setting built from $(T,f,\chi,\nu',m_0,\mathrm{top},A_0)$ and let $v$ be a place, with parameter
$P_v=(k,r,T,T',u,u')$, threshold $T_T=TT'>0$ and real place datum $B_v$ (the base of the tower $T_v$), with vectors $z,z'$
and representations $\pi,\pi'$. We call $S$ an \emph{admissible built place with constant $C$} if it is a built place with
constant $C$, all its place data are admissible and the words $d_e,d'_e$ are admissible.
\begin{enumerate}[label=(\roman*)]
\item The \emph{open domain} at $v$ is $\Omega_v=\{g:G_f(z,z';g)>0,\ D_a(g)\ne0\ (1\le a\le r)\}\subset\mathrm{Mat}_{k,r}$,
  where the $D_a$ are the leading minors attached to $(z,z')$. The class $\mathcal F_v$ of \emph{smooth forms} consists of
  the functions $J\colon\mathrm{Mat}_{k,r}\to\Hom(W_q,W_p)$ which are smooth on $\Omega_v$, satisfy there the left and right
  differential equations of level $0$, and satisfy $J(hgh')=\pi(h)J(g)\pi'(h')$ for $g\in\Omega_v\cap G_{k,r}$ and $h,h'$ in
  the unipotent groups $U_W(z),U_W(z')$.
\item Let $A_v=\{\zeta\in\C:|\zeta|^2>T_T\}$ and, for $\zeta\in A_v$, let $\mathrm{Fib}(\zeta)\subset\Hom(W_q,W_p)$ be the
  fibre of the radial domain at $\zeta$ (defined through the section of $B_v$). A \emph{frame} at $v$ is a family
  $e=(e_\rho)$, indexed by the types $\rho$ of $P_v$, of smooth maps $e_\rho\colon A_v\to\Hom(W_q,W_p)$ such that for every
  $\zeta\in A_v$ the $e_\rho(\zeta)$ lie in $\mathrm{Fib}(\zeta)$, are orthonormal for the Hilbert--Schmidt inner product
  $\langle\cdot,\cdot\rangle_{\mathrm{HS}}$, and span $\mathrm{Fib}(\zeta)$.
\item The \emph{transfer} of $J\colon\mathrm{Mat}_{k,r}\to\Hom(W_q,W_p)$ by a frame $e$ is the family
  $\mathrm{tr}_e(J)=(\mathrm{tr}_e(J)_\rho)_\rho$ of functions on $Q=Q_{r,T_T}$ given by
  \[
  \mathrm{tr}_e(J)_\rho(x,z)=w_k(x,z)^{1/2}\bigl\langle J(g_q(x,z)),e_\rho(z)\bigr\rangle_{\mathrm{HS}},\qquad
  w_k(x,z)=\prod_l|x_l|^{2(m_l-1)}(|z|^2-T_T)^{k-1},
  \]
  if this function is smooth on $Q$, and by $\mathrm{tr}_e(J)_\rho=0$ otherwise; here $g_q(x,z)\in\mathrm{Mat}_{k,r}$ is the
  big-cell matrix attached to $(x,z)$ by $B_v$, and the $m_l$ are explicit integers.
\item A frame $e$ satisfies the \emph{gauge condition} if $\mathrm{tr}_e(\mathrm{CV}(Y)J)=C_0(Y)\,\mathrm{tr}_e(J)$ for all
  $J\in\mathcal F_v$ and $Y\in\R$, where $C_0(Y)$ is the operator of the model frame in the model operator algebra (acting on
  families); and the \emph{radial gauge condition} if for every radial family $\Phi$, $Y\in\R$ and $a,b\in\{0,1\}$ the
  transfer of $Q_L(Y,a,b)\Phi$ (the radial operator in the last variable) is $QQ_r(Y,a,b)$ applied to the transfer of
  $\Phi$.
\end{enumerate}
\formal{Formalism.FourierDecayAlt.P.CutMP.BuiltOW}
\formal{Formalism.FourierDecayAlt.P.CutMP.OmP}
\formal{Formalism.FourierDecayAlt.Q.FFl}
\formal{Formalism.FourierDecayAlt.P.CutMP.IsFrame}
\formal{Formalism.FourierDecayAlt.P.CutMP.trJ}
\formal{Formalism.FourierDecayAlt.P.CutMP.IsGauge}
\formal{Formalism.FourierDecayAlt.P.CutMP.IsGaugeQQ}
\end{definition}

\begin{definition}[Hat operators and the model conditions]\label{pf:thmS-carrier:hatdef}
Let $P$ be a parameter and $L$ a level datum of the model operator algebra which is a \emph{model level} (its structure
operators are those of the model). The \emph{hat set} $\widehat{\mathcal H}(L)$ consists of the operators $\widehat A(y)$,
$R_0(y)$, $\widehat A^\vee(y)$, $R_0^\vee(y)$ of $L$ ($y\in\C$). The \emph{bounded set} $\mathcal B(L)$ consists of
$\zeta^{-1}$, $\bar\zeta^{-1}$, $D\bar D^{-1}$, $\bar DD^{-1}$, the multiplications $\kappa(c)$ by functions $c$ of the type,
and the operators $\bar D_rD_r^{-1}S_i$ and their adjoints ($1\le i\le k-1$); $\langle\mathcal B(L)\rangle$ denotes the
monoid it generates. Let $\langle u,w\rangle_Q=\sum_\rho\langle u_\rho,w_\rho\rangle_{L^2(\mu_0)}$ be the \emph{model inner
product}. A family $j=(j_\rho)$ of smooth functions on $Q_{r,T_T}$ satisfies
\begin{itemize}
\item the \emph{model decay condition to order $m$} if $|(X_1\cdots X_pj)_\rho|^2$ is $\mu_0$-integrable for all $p\le m$,
  all $X_i\in\widehat{\mathcal H}(L)$ and all types $\rho$;
\item the \emph{model adjointness condition to order $m$} if for every log-smooth $\varphi$ which vanishes on
  $(-\infty,\eps]$ for some $\eps>0$ and has bounded logarithmic derivatives, all $b\in\langle\mathcal B(L)\rangle$,
  $X\in\widehat{\mathcal H}(L)$ and all products $V,W$ of $p$ and $q$ elements of $\widehat{\mathcal H}(L)$ with
  $p+q+1\le m$,
  \[
  \bigl\langle X(\varphi(G_q)\,bVj),Wj\bigr\rangle_Q=\bigl\langle\varphi(G_q)\,bVj,X^*Wj\bigr\rangle_Q .
  \]
\end{itemize}
\formal{Formalism.FourierDecayAlt.MA.IsModelLevel}
\formal{Formalism.FourierDecayAlt.MA.hatSet}
\formal{Formalism.FourierDecayAlt.MA.bddSet}
\formal{Formalism.FourierDecayAlt.MA.JDecay}
\formal{Formalism.FourierDecayAlt.MA.JAdj}
\end{definition}

These are the decay and adjointness conditions of the model used in \cref{pf:thmS-carrier:modelplace-def} and in the model
lemma (\cref{pf:thmS-carrier:lemmamodel}).

\subsubsection{Model places from the transfer statements}

\begin{lemma}[Local facts at built places]\label{pf:thmS-carrier:localfacts}
Let $S$ be an admissible built place with constant $C$ and let $v$ be a place.
\begin{enumerate}[label=(\roman*)]
\item Every letter $(b,Y)$ of a local word admissible at $v$ satisfies $Y\in[\alpha_v,\beta_v]$.
\item For $s\in\C$ and $Y\in\R$ the map $J\mapsto s\,\mathrm{CV}(Y)J$ preserves $\mathcal F_v$.
\item If $C\ge100$, then $(\mathrm{CV}(Y)J_{f_v})(g_q(x,z))=P_v(Y)\,J_{f_v}(g_q(x,z))$ for all $Y\in\R$ and
  $(x,z)\in Q_{r,T_T}$.
\end{enumerate}
\formallabel{Formalism.FourierDecayAlt.P.CutMP.admissible_in_truncation_holds}
\formallabel{Formalism.FourierDecayAlt.P.CutMP.coefficient_vector_preserves_form}
\formallabel{Formalism.FourierDecayAlt.P.CutMP.bessel_eigenvalue_equation_holds}
\end{lemma}

\begin{proof}
(i) By the definition of admitted words, a letter $(1,Y)$ has $Y$ in a grid $\mathrm{grid}_v(m)$ or in $\mathrm{grid}^1_v$,
and a letter $(0,Y)$ has $Y-1$ there. The grid condition, which is part of the definition of built places, states that
$y$ and $y+1$ lie in the cut interval of the place for every grid point $y$; by the definitions, this cut interval is
contained in $[\alpha_v,\beta_v]$.

(ii) By \lref{Formalism.FourierDecayAlt.P.operator_preserves_fusion_forms}, the Casimir-type operator $\mathrm{CV}(Y)$
maps $\mathcal F_v$ into itself (for the domain $\Omega_v$ of admissible vectors $z,z'$), and $\mathcal F_v$ is closed under
multiplication by scalars.

(iii) Let $g=g_q(x,z)$. By \lref{Formalism.FourierDecayAlt.P.big_cell_domain}, $g$ lies in the open domain of the
Bessel transform and in $G_{k,r}$. The Bessel transform is an integral $J_F(g)=\int\xi(p)\,(\rho(p)F)(g)\,d\mu(p)$ of
twisted two-sided translates $\rho(p)F$ against a character $\xi$ (\cref{pf:thmS-carrier:realplace}). The function $f_v$ is
smooth with the growth bounds of a built place, and the slice budget of a built place gives the threshold condition on
$A_0$ required by \lref{Formalism.FourierDecayAlt.P.casimir_under_jacquet_integral}, which states that then
$\mathrm{CV}(Y)J_{f_v}(g)=\int\xi(p)\,\mathrm{CV}(Y)(\rho(p)f_v)(g)\,d\mu(p)$. For almost every $p$,
\lref{Formalism.FourierDecayAlt.P.capelli_rho_commute} gives $\mathrm{CV}(Y)(\rho(p)f_v)(g)=\rho(p)(\mathrm{CV}(Y)f_v)(g)$,
and the translated points at which $\mathrm{CV}(Y)f_v$ is evaluated lie in $G_{k,r}$, where $\mathrm{CV}(Y)f_v=P_v(Y)f_v$ by
the eigenfunction condition of built places. Hence the integrand equals $P_v(Y)$ times the integrand of $J_{f_v}(g)$, and
integration gives the claim.
\end{proof}

\begin{lemma}[Cut-offs and the transfer]\label{pf:thmS-carrier:cutofflemma}
Let $S$ be an admissible built place such that $J_{f_v}\in\mathcal F_v$ for all $v$, and let $\eps>0$.
\begin{enumerate}[label=(\roman*)]
\item $J_{\mathrm{cut}_\eps f_v}(g)=\chi(G_f(z,z';g)/\eps)\,J_{f_v}(g)$ for all $g$, and $J_{\mathrm{cut}_\eps f_v}\in
  \mathcal F_v$.
\item For every frame $e$ at $v$ and every cut $c$ with cut-off function $\chi$,
  $\mathrm{tr}_e(J_{\mathrm{cut}_\eps f_v})=\chi_\eps(G_q)\,\mathrm{tr}_e(J_{f_v})$, where $G_q(x,z)=\prod_l|x_l|^2(|z|^2-T_T)$.
\end{enumerate}
\formallabel{Formalism.FourierDecayAlt.P.CutMP.regularity_smooth_forms_holds}
\formallabel{Formalism.FourierDecayAlt.P.CutMP.cutoff_transport_to_model}
\end{lemma}

\begin{proof}
(i) The identity is \lref{Formalism.FourierDecayAlt.P.JCTx2.JF_cutF}, which states that the Bessel transform of
$\mathrm{cut}_\eps F$ is $\chi(G_f/\eps)J_F$. The cut-off $\chi$ is smooth (part of the definition of built places), and
\lref{Formalism.FourierDecayAlt.P.CutMP.cut_mem_FFl} states that multiplication by $\chi(G_f/\eps)$ preserves
$\mathcal F_v$ for smooth $\chi$.

(ii) For $(x,z)\in Q_{r,T_T}$ one has $G_f(z,z';g_q(x,z))=G_q(x,z)$ by
\lref{Formalism.FourierDecayAlt.P.CutMP.cutTr_Gf_gq}. Hence, by (i), the function defining
$\mathrm{tr}_e(J_{\mathrm{cut}_\eps f_v})_\rho$ is $\chi(G_q/\eps)$ times the function
$u_\rho=w_k^{1/2}\langle J_{f_v}\circ g_q,e_\rho\rangle_{\mathrm{HS}}$ defining $\mathrm{tr}_e(J_{f_v})_\rho$. The function
$u_\rho$ is smooth on $Q_{r,T_T}$: $J_{f_v}\circ g_q$ is smooth there because $J_{f_v}$ is smooth on $\Omega_v$ and $g_q$ maps
$Q_{r,T_T}$ smoothly into $\Omega_v$ (\lref{Formalism.FourierDecayAlt.P.CutMP.cutTr_jgq_smooth}), $e_\rho$ is smooth, and
$w_k$ is smooth and positive on $Q_{r,T_T}$. So $\mathrm{tr}_e(J_{f_v})_\rho=u_\rho$, the product $\chi(G_q/\eps)u_\rho$ is
smooth, and $\mathrm{tr}_e(J_{\mathrm{cut}_\eps f_v})_\rho=\chi(G_q/\eps)u_\rho$.
\end{proof}

\begin{theorem}[Model places from the transfer statements]\label{pf:thmS-carrier:openQQ}
Assume the following six statements. Then \ref{pf:thmS-carrier:I2} holds.
\begin{enumerate}[label=(T\arabic*)]
\item (Smoothness) There is $C_1$ such that $J_{f_v}\in\mathcal F_v$ for every built place with constant $C_1$ and every $v$.
\item (Gauge core) For every parameter with $k\ge2$ and every real place datum which is a place for it, there is a frame
  of the fibres of that datum satisfying the radial gauge condition.
\item (Energy trace identity) There is $C_3$ such that for every admissible built place $S$ with constant $C_3$, every place
  $v$ and every frame $e$ at $v$: if $\mathcal F_v$ is preserved by the maps $J\mapsto s\,\mathrm{CV}(Y)J$, $\eps>0$ and
  $J_\eps=J_{\mathrm{cut}_\eps f_v}\in\mathcal F_v$, then for every admissible local word
  $a=((b_1,Y_1),\dots,(b_m,Y_m))$
  \begin{equation}\label{eq:pf:thmS-carrier:energytrace}
  \mathcal E_v(a;\mathrm{cut}_\eps f_v)=\bigl\langle\mathrm{tr}_e\bigl(\sigma\mathrm{CV}(Y_m)\cdots\sigma\mathrm{CV}(Y_1)
  J_\eps\bigr),\mathrm{tr}_e(J_\eps)\bigr\rangle_Q,
  \end{equation}
  where $\sigma=\pm1$ is the sign attached to $P_v$ and the level $l+1$.
\item (Reality) There is $C_4$ such that for every admissible built place with constant $C_4$ the energies
  $\mathcal E_v(a;\mathrm{cut}_\eps f_v)$, $\eps>0$, $a$ admissible, are real.
\item (Decay transfer) There is $C_5$ such that for every admissible built place with constant $C_5$, every place $v$, every
  frame $e$ at $v$ satisfying the gauge and radial gauge conditions and every model level $L$: if $J_{f_v}\in\mathcal F_v$,
  then $\mathrm{tr}_e(J_{f_v})$ satisfies the model decay condition to order $2\thmScarrierNtop+2$.
\item (Adjointness transfer) The same as (T5) with the model adjointness condition, with a constant $C_6$.
\end{enumerate}
\formallabel{Formalism.FourierDecayAlt.P.CutMP.modelPlaceBuilt2_of_openQQ}
\formallabel{Formalism.FourierDecayAlt.P.CutMP.modelPlaceBuilt3_of_openGC}
\formallabel{Formalism.FourierDecayAlt.P.CutMP.modelPlaceBuilt3_of_openQQ}
\formallabel{Formalism.FourierDecayAlt.P.CutMP.modelPlaceBuilt3_of_leavesQQ}
\formallabel{Formalism.FourierDecayAlt.P.CutMP.modelPlaceBuilt2_of_3}
\end{theorem}

\begin{proof}
\emph{Constants.} Put $C=\max(1,C_3,C_4,C_5,C_6,100)$ and $C_c=\max(C_1,C)$. Being a built place with a constant $C'$ implies
being one with every smaller constant, because the constant enters the definition only through conditions that persist
when it is decreased (the long-cut condition, $C'\le C_A$, and the lower bound for $A_0$ in terms of $C'$). Hence an admissible built place $S$ with constant $C_c$ satisfies the
hypotheses of (T1), (T3)--(T6) and \cref{pf:thmS-carrier:localfacts}(iii). Fix such an $S$, a cut $c$ with cut-off function
$\chi$ and a place $v$; by (T1), $J:=J_{f_v}\in\mathcal F_v$ for every place.

\emph{Step 1: a gauge frame.} By \lref{Formalism.FourierDecayAlt.P.CutMP.gaugeExistsQQ_of_gCore}, which derives this from
(T2) applied to the real place datum $B_v$, there is a frame $e$ at $v$ satisfying the gauge condition and the radial gauge
condition.

\emph{Step 2: the model data.} The parameter $P_v$ is admissible and $1\le l+1\le k$, since $1\le l\le k-2$ for a built place.
By \lref{Formalism.FourierDecayAlt.P.model_frame_analytic} (the analysis of the model frame: for every admissible
parameter and $1\le l'\le k$ there is a model level datum starting at the model frame of level $l'$) we choose a model level
$L$ starting at level $l+1$. Put
\[
j=\mathrm{tr}_e(J),\qquad\lambda_s=\lambda_s(v)\ (1\le s\le k+2r+1),\qquad C_m=\sigma=1,
\]
where $\lambda_s(v)$ are the numbers of the configuration of the place $v$ in the global size data.

\emph{Step 3: the model Hida conditions.} The parameter is admissible, the levels are in range and $L$ is a model level.
By \lref{Formalism.FourierDecayAlt.P.CutMP.eigenvalue_configuration_interlaces}, which holds for $k\ge2$ and admissible
place data and words, the configuration $(\lambda_s)$ interlaces with the parameters $u$ and $u'$ of $P_v$. Since $C_c\ge1$
and the long-cut condition gives $\beta_v-\alpha_v\ge C_c(k\thmScarrierNtop(r+1)^2+\thmScarrierNtop^2(r+1)^2+
\thmScarrierNtop^3)>0$, we have $\alpha_v<\beta_v$; and the cut interval lies on the cut of the parameter, by the definition of
built places.

\emph{Step 4: decay and adjointness.} These are (T5) and (T6) applied to $e$ and $L$, since $J\in\mathcal F_v$.

\emph{Step 5: the eigenfunction condition.} Let $Y\in\R$. By the gauge condition, $C_0(Y)j=\mathrm{tr}_e(\mathrm{CV}(Y)J)$. The
transfer only involves the values of its argument at the points $g_q(x,z)$, $(x,z)\in Q_{r,T_T}$, and there
$\mathrm{CV}(Y)J=P_v(Y)J$ by \cref{pf:thmS-carrier:localfacts}(iii). Since the transfer is linear on such functions,
$C_0(Y)j=P_v(Y)\,j$, and $P_v(Y)=\prod_s(Y+\lambda_s(v))$ by the definition of $P_v$ as the polynomial of the configuration.

\emph{Step 6: the energy.} Let $\eps>0$ and let $a=((b_1,Y_1),\dots,(b_m,Y_m))$ be admissible. By
\cref{pf:thmS-carrier:localfacts}(ii) and \cref{pf:thmS-carrier:cutofflemma}(i) the hypotheses of (T3) hold, so
\eqref{eq:pf:thmS-carrier:energytrace} holds. By \lref{Formalism.FourierDecayAlt.P.CutMP.trJ_foldr}, which iterates the
gauge condition (the intermediate functions $\sigma\mathrm{CV}(Y_i)\cdots\sigma\mathrm{CV}(Y_1)J_\eps$ lie in $\mathcal F_v$
by \cref{pf:thmS-carrier:localfacts}(ii)),
\[
\mathrm{tr}_e\bigl(\sigma\mathrm{CV}(Y_m)\cdots\sigma\mathrm{CV}(Y_1)J_\eps\bigr)=C^{(l+1)}_0(Y_m)\cdots C^{(l+1)}_0(Y_1)\,
\mathrm{tr}_e(J_\eps),
\]
where $C^{(l+1)}_0(Y)$ is the operator $C_0(Y)$ of the model start frame at level $l+1$. By
\cref{pf:thmS-carrier:cutofflemma}(ii), $\mathrm{tr}_e(J_\eps)=\chi_\eps(G_q)j$. Therefore
\[
\mathcal E_v(a;\mathrm{cut}_\eps f_v)=\Bigl\langle C^{(l+1)}_0(Y_m)\cdots C^{(l+1)}_0(Y_1)\chi_\eps(G_q)j,\ \chi_\eps(G_q)j
\Bigr\rangle_Q .
\]
By (T4) the left side is real, so it equals the real part of the right side, which is the cut-off energy form
$\thmScarrierEn^{\mathrm{mod}}_v(a;\eps)$ (\cref{pf:thmS-carrier:modelsys}). This is \eqref{eq:pf:thmS-carrier:modelenergy}
with $C_m=\sigma=1$.

\emph{Step 7: admission and size.} By \cref{pf:thmS-carrier:localfacts}(i) the letters of $a$ lie in $[\alpha_v,\beta_v]$.
Then \lref{Formalism.FourierDecayAlt.P.CutMP.level_system_admissible_size} states that the model system with interval
$[\alpha_v,\beta_v]$ and any box length admits $a$ and gives it the size $\mathrm{size}_v(a)$.

Steps 2--7 show that $v$ is a model place for $(L,j,\lambda,c,1,1)$, which proves \ref{pf:thmS-carrier:I2} with the
constant $C_c$.
\end{proof}

\subsubsection{Realisation of the hat operators}

\begin{theorem}[Realisation of the hat operators]\label{pf:thmS-carrier:hatreal}
There is $C_c\ge0$ with the following property. Let $S$ be an admissible built place with constant $C_c$, $v$ a place, $e$ a
frame at $v$ satisfying the gauge and radial gauge conditions, and $L$ a model level; put $n=2r+k+1$. Then every
$X\in\widehat{\mathcal H}(L)$ is \emph{realised}: there are finitely many pairs $(b_t,\Theta_t)$, where
$b_t\in\langle\mathcal B(L)\rangle$ and $\Theta_t$ is a differential operator on $\mathrm{Mat}_{k,r}$ of order at most
$C_cn$, with coefficients smooth on $\{G_f>0\}$ and all of whose derivatives grow at most like the
$(C_c(r+1)n)$-th power of the norm on $G_{k,r}$, commuting on $\{G_f>0\}$ with the translations $\rho(p)$ of the Bessel
integral, such that
\begin{equation}\label{eq:pf:thmS-carrier:realised}
X\,\mathrm{tr}_e(J)=\sum_tb_t\,\mathrm{tr}_e(\Theta_tJ)\qquad(J\in\mathcal F_v).
\end{equation}
\formallabel{Formalism.FourierDecayAlt.P.CutMP.hat_quadratic_operator_holds}
\formallabel{Formalism.FourierDecayAlt.P.CutMP.jDHatOpQQ_of_opA_opR_residue_splitting}
\end{theorem}

\begin{proof}
\emph{Closure properties.} Fix $C_c$ (chosen below). Call an element $X$ of the model operator algebra realised if it has a
representation \eqref{eq:pf:thmS-carrier:realised} with operators $\Theta_t$ of the stated kind. An element with
$X\mathrm{tr}_e(J)=\mathrm{tr}_e(\Theta J)$ for a single such $\Theta$ is realised (take $b=1$). Realised elements are closed
under left multiplication by $b\in\langle\mathcal B(L)\rangle$: if $X\mathrm{tr}_e(J)=\sum_tb_t\mathrm{tr}_e(\Theta_tJ)$, then
$bX\mathrm{tr}_e(J)=\sum_t(bb_t)\mathrm{tr}_e(\Theta_tJ)$, because the action of the model operator algebra on families is
multiplicative, and $bb_t\in\langle\mathcal B(L)\rangle$. They are closed under sums (concatenate the lists), and under
multiplication by a scalar $c$, since $cX=\kappa(c)X$ for the constant function $c$ of the type and $\kappa(c)\in\mathcal
B(L)$; hence also under differences and finite linear combinations.

\emph{The generators.} Combining the decay of admissible polynomial operators
\lref{Formalism.FourierDecayAlt.P.CutMP.polynomial_decay_admissible} with the Capelli identities for the minors
\lref{Formalism.FourierDecayAlt.P.CutMP.capelli_minors_holds}, we obtain $K_A\ge0$ such that for $C_c\ge K_A$
every $\widehat A(y)$ and $\widehat A^\vee(y)$ ($y\in\C$) is realised by a single operator:
$\mathrm{tr}_e(\Theta J)=\widehat A(y)\mathrm{tr}_e(J)$ and $\mathrm{tr}_e(\Theta^\vee J)=\widehat A^\vee(y)\mathrm{tr}_e(J)$ for
$J\in\mathcal F_v$. Likewise, by \lref{Formalism.FourierDecayAlt.P.CutMP.residue_operator_holds} there is $K_R\ge0$ such
that for $C_c\ge K_R$ and $y\in\R$ the \emph{homogeneous residue operators}
$R_h(y)=D_{\bar r}\bigl(t_{21}(y)A_r(y)+t_{22}(y)B_r(y)\bigr)$ and their mirrors $R_h^\vee(y)$ are realised by single
operators. Put $C_c=\max(K_A,K_R)$.

\emph{Nodes and residue splitting.} For $y\in\C$ let $x^{-1}(y)=\kappa\bigl((x_C(\pi,y)^{-1})_\pi\bigr)$, defined when $y$ is
\emph{regular}, i.e.\ $x_C(\pi,y)\ne0$ for all types $\pi$; it lies in $\mathcal B(L)$. By
\lref{Formalism.FourierDecayAlt.P.CutMP.hat_nodes_decay_holds} there are pairwise distinct regular real numbers
$\nu_0,\dots,\nu_{r+1}$. By the residue splitting \lref{Formalism.FourierDecayAlt.P.CutMP.residue_splitting_holds}, for
regular $y$
\[
R_0(y)=x^{-1}(y)R_h(y)-\sum_{l=1}^{k-1}x^{-1}(y)\,S^\natural_l\,\kappa\bigl((x_{\mathrm{Ex}}(\{l\},\pi,y))_\pi\bigr)\,
\widehat A(y),
\]
where $S^\natural_l=\bar D_rD_r^{-1}S_l\in\mathcal B(L)$; and the mirror splitting
\lref{Formalism.FourierDecayAlt.P.CutMP.mirror_residue_splitting_holds} expresses $R_0^\vee(y)$ in the same way through
$R_h^\vee(y)$, $\widehat A^\vee(y)$ and the adjoints of the $S^\natural_l$. By the closure properties, $R_0(\nu_j)$ and
$R_0^\vee(\nu_j)$ are realised for every $j$.

\emph{Interpolation.} The operators $R_0(y)$ and $R_0^\vee(y)$ are determined by their values at $r+2$ nodes:
\lref{Formalism.FourierDecayAlt.P.CutMP.hat_lagrange_decay_holds} states that for pairwise distinct real nodes
$\nu_0,\dots,\nu_{r+1}$ and all $y\in\C$
\[
R_0(y)=\sum_jc_j(y)R_0(\nu_j),\qquad R_0^\vee(y)=\sum_jc_j(y)R_0^\vee(\nu_j),
\]
with the Lagrange coefficients $c_j(y)=\prod_{i\ne j}(y-\nu_i)/(\nu_j-\nu_i)$. By the closure properties every $R_0(y)$ and
$R_0^\vee(y)$ is realised. Together with the realisation of $\widehat A(y)$ and $\widehat A^\vee(y)$ this covers all of
$\widehat{\mathcal H}(L)$.
\end{proof}

\subsubsection{The transfer of the star operation}

Let $\mathcal P_v\subset\mathcal F_v$ be the \emph{positive cone of smooth forms} at $v$ and let $\Theta=(\Theta_1,\dots,
\Theta_p)$ be a list of differential operators of the class used for the adjointness transfer (operators with an explicit
formal adjoint); write $\Theta J=\Theta_1\cdots\Theta_pJ$ and $\Theta^*J$ for the application of the list of formal adjoints
in the opposite order. A \emph{test family} is a family $U=(U_\rho)$ of smooth functions with compact support in
$Q_{r,T_T}$.

\begin{theorem}[Transfer of the star operation]\label{pf:thmS-carrier:startransfer}
There is $C_c$ with the following property. Let $S$ be an admissible built place with constant $C_c$, $v$ a place, $e$ a
frame at $v$ satisfying the gauge condition, $L$ a model level and $a$ an element of the monoid generated by
$\widehat{\mathcal H}(L)\cup\mathcal B(L)$. Let $(c_t,\Theta_t)_t$ be finitely many pairs of scalars and lists of operators
as above such that for all $J\in\mathcal P_v$ the forms $\Theta_tJ$ and $\Theta_t^*J$ lie in $\mathcal P_v$ and
$a\,\mathrm{tr}_e(J)=\sum_tc_t\,\mathrm{tr}_e(\Theta_tJ)$. Then for all $J\in\mathcal P_v$
\[
a^*\,\mathrm{tr}_e(J)=\sum_t\bar c_t\,\mathrm{tr}_e(\Theta_t^*J).
\]
\formallabel{Formalism.FourierDecayAlt.P.CutMP.JA.star_operator_transfer_holds}
\formallabel{Formalism.FourierDecayAlt.P.CutMP.JA.jAStarOp_of_pairL}
\formallabel{Formalism.FourierDecayAlt.P.CutMP.JA.SA.jAStarOp_of_parts}
\end{theorem}

\begin{proof}
We use four statements.
\begin{enumerate}[label=(\alph*)]
\item (Fundamental identity) \lref{Formalism.FourierDecayAlt.P.CutMP.JA.SA.fundamental_quadratic_identity}: two families
  $u,w$ of smooth functions on $Q_{r,T_T}$ with $\langle u,U\rangle_Q=\langle w,U\rangle_Q$ for all test families $U$ are
  equal.
\item (Model adjoint identity) \lref{Formalism.FourierDecayAlt.P.CutMP.JA.SA.model_adjoint_identity_holds}: for $T_T>0$,
  a model level $L$, $a$ in the monoid generated by $\widehat{\mathcal H}(L)\cup\mathcal B(L)$, a family $u$ and a test family
  $U$, $\langle a^*u,U\rangle_Q=\langle u,aU\rangle_Q$.
\item (Sections) \lref{Formalism.FourierDecayAlt.P.CutMP.JA.self_adjoint_section_holds}: there is a constant such that
  at admissible built places with that constant, for every frame $e$, every test family is $\mathrm{tr}_e(K)$ for some
  $K\in\mathcal P_v$.
\item (Pairing) \lref{Formalism.FourierDecayAlt.P.CutMP.JA.self_adjoint_pairing_list_holds}: there is a constant such that
  at admissible built places with that constant, for every frame $e$, every list $\Theta$ as above and all
  $K,J\in\mathcal P_v$ with $\mathrm{tr}_e(K)$ a test family, $\mathrm{tr}_e(\Theta K)$ is a test family and
  $\langle\mathrm{tr}_e(\Theta K),\mathrm{tr}_e(J)\rangle_Q=\langle\mathrm{tr}_e(K),\mathrm{tr}_e(\Theta^*J)\rangle_Q$.
\end{enumerate}
Let $C_c$ be the maximum of the constants of (c) and (d). Since $T_T=TT'>0$, (b) applies. By (a) it suffices to show
$\langle a^*\mathrm{tr}_e(J),U\rangle_Q=\langle\sum_t\bar c_t\mathrm{tr}_e(\Theta_t^*J),U\rangle_Q$ for every test family $U$.
By (c), $U=\mathrm{tr}_e(K)$ with $K\in\mathcal P_v$. By (b) and the hypothesis applied to $K$,
\[
\langle a^*\mathrm{tr}_e(J),\mathrm{tr}_e(K)\rangle_Q=\langle\mathrm{tr}_e(J),a\,\mathrm{tr}_e(K)\rangle_Q
=\Bigl\langle\mathrm{tr}_e(J),\sum_tc_t\mathrm{tr}_e(\Theta_tK)\Bigr\rangle_Q .
\]
Since $\mathrm{tr}_e(K)=U$ is a test family, (d) and the Hermitian symmetry of $\langle\cdot,\cdot\rangle_Q$ give
$\langle\mathrm{tr}_e(J),\mathrm{tr}_e(\Theta_tK)\rangle_Q=\langle\mathrm{tr}_e(\Theta_t^*J),\mathrm{tr}_e(K)\rangle_Q$ for every
$t$. For a sesquilinear Hermitian form one has $\langle w,cu\rangle_Q=\langle\bar cw,u\rangle_Q$ for every scalar $c$
(whichever variable is the linear one). Hence the right side of the last display equals
$\sum_t\langle\bar c_t\,\mathrm{tr}_e(\Theta_t^*J),\mathrm{tr}_e(K)\rangle_Q=\langle\sum_t\bar c_t\,\mathrm{tr}_e(\Theta_t^*J),
U\rangle_Q$, as required.
\end{proof}

\subsubsection{The polynomial radial statement}

Let $B_s$ be a real place datum of a parameter $P$ with $k\ge2$, with section $\mu$, and let $A=\{\zeta:|\zeta|^2>T_T\}$.
For $\zeta\in A$ the section provides vectors $\nu_0(\zeta),\nu'_0(\zeta)$ and a matrix $\iota_0(\zeta)$. Let $X$ be an
actual-operator datum and $d$ a normal-form datum on the radial space satisfying the \emph{action hypotheses}: $B_s$ is a
place for $P$, the level datum of $d$ is an actual level and a star level, the normal-form relations hold, and
$X.QQ(Y)=d.QQ(Y)$ for all $Y$. For $1\le l\le k-1$, a \emph{pointwise shift} at level $l$ is a family
$S(\zeta)$ of linear maps of $\Hom(W_q,W_p)$ such that $(S_l\Phi)(c,\mu(\zeta))=S(\zeta)\bigl(\Phi(c,\mu(\zeta))\bigr)$ for
every radial function $\Phi$, $c\in(\C^\times)^r$ and $\zeta\in A$, where $S_l$ is the shift operator of $d$. A \emph{type
projection family} is a smooth family $\mathrm{Pr}_\rho(a)$ of orthogonal projections on $W_p$, indexed by the types and
homogeneous of degree $0$ and equivariant in $a$, such that $\mathrm{Pr}_\rho(a)x=x$ exactly when $x$ is an eigenvector of
the Capelli representatives with the eigenvalues prescribed by $\rho$. For frames $B_w,B'_w$ of $z,z'$, let
$\theta_\zeta=\theta(\nu_0(\zeta),\nu'_0(\zeta),\iota_0(\zeta))$ be the explicit sum of sandwich operators
$X\mapsto\sum_{p,q}H(\iota_0B'_{w,q},B_{w,p})E_{p}XE'_{q}$ built from the frames, where $E_p,E'_q$ are the images under
$\pi,\pi'$ of elementary matrices attached to $B_{w,p}$, $\nu_0(\zeta)$ and $\nu'_0(\zeta)$, $B'_{w,q}$.

\begin{lemma}[Phase extraction]\label{pf:thmS-carrier:phaseextract}
Let $T_T>0$, $g(\zeta)=|\zeta|^2-T_T$ and $\varsigma(\zeta)=\sqrt{|\zeta|^2-T_T}$ on $A$. Let $F\colon\C\to\C$ be such that
$g^mF=p(\zeta,\bar\zeta,\varsigma)$ on $A$ for some $m\in\N$ and some polynomial $p$ in three variables, and such that
$F(e^{i\theta}\zeta)=F(\zeta)$ for $\theta\in\R$, $\zeta\in A$. Then there is a polynomial $q$ in one variable with
$g(\zeta)^mF(\zeta)=q(\varsigma(\zeta))$ on $A$.
\formallabel{Formalism.FourierDecayAlt.P.CutMP.JA.SSn.Rep.Core.TC.PR.phase_extraction_holds}
\formallabel{Formalism.FourierDecayAlt.P.CutMP.JA.SSn.Rep.Core.TC.PR.evenPart_poly}
\end{lemma}

\begin{proof}
Let $\zeta\in A$ and $x=|\zeta|$. Rotation invariance with $\theta=\arg\zeta$ and $\theta=\pi$ gives $F(\zeta)=F(x)=F(-x)$,
and $g,\varsigma$ take the same values at $\zeta$, $x$ and $-x$. Hence
$2g(\zeta)^mF(\zeta)=p(x,x,\varsigma)+p(-x,-x,\varsigma)$. Write $p(t,t,s)=\sum_{i,j}a_{ij}t^is^j$; the even part
$\sum_{i\text{ even}}2a_{ij}t^is^j$ is a polynomial in $t^2$ and $s$, and $x^2=\varsigma^2+T_T$. Substituting gives a polynomial
$q_0$ in one variable with $p(x,x,\varsigma)+p(-x,-x,\varsigma)=q_0(\varsigma)$; put $q=q_0/2$.
\end{proof}

\begin{theorem}[Polynomial radial statement]\label{pf:thmS-carrier:polyradial}
Let $B_s,X,d$ satisfy the action hypotheses, $1\le l\le k-1$, let $S$ be a pointwise shift at level $l$, $\mathrm{Pr}$ a type
projection family, $B_w,B'_w$ frames of $z,z'$, and let $\pi,\rho$ be types such that $\rho$ is obtained from $\pi$ by
raising at level $l$. Then there are $n\in\N$ and a polynomial $q$ such that for all $\zeta\in A$ and all $X_f$ in the fibre
$\mathrm{Fib}(\zeta)$ with $\mathrm{Pr}_\pi(\nu_0(\zeta))X_f=X_f$,
\begin{equation}\label{eq:pf:thmS-carrier:polyradial}
(|\zeta|^2-T_T)^n\,\mathrm{Pr}_\rho(\nu_0(\zeta))\,S(\zeta)X_f=q\bigl(\sqrt{|\zeta|^2-T_T}\bigr)\,
\mathrm{Pr}_\rho(\nu_0(\zeta))\,\theta_\zeta X_f .
\end{equation}
\formallabel{Formalism.FourierDecayAlt.P.CutMP.JA.SSn.Rep.Core.TC.PR.polynomial_radial_holds}
\formallabel{Formalism.FourierDecayAlt.P.CutMP.JA.SSn.Rep.Core.TC.PR.polyRadialS_of}
\formallabel{Formalism.FourierDecayAlt.P.CutMP.JA.SSn.Rep.Core.TC.PR.polynomialRadial_of_base}
\formallabel{Formalism.FourierDecayAlt.P.CutMP.JA.SSn.Rep.Core.TC.PR.polynomial_radial_base_holds}
\end{theorem}

\begin{proof}
\emph{Reduction to the canonical section.} \Lref{Formalism.FourierDecayAlt.P.CutMP.JA.SSn.Rep.Core.TC.PR.transfer_holds}
states that, given the identities for the theta operators
\lref{Formalism.FourierDecayAlt.P.CutMP.JA.SSn.Rep.theta_facts_hold}, the theorem for an arbitrary section follows from
the theorem for the canonical section $\mu_0$ of $(z,z')$: the two sections differ by $\mu(\zeta)=h(\zeta)\mu_0(\zeta)h'(\zeta)$
with $h,h'$ in the unipotent groups, and conjugating $S(\zeta)$ by $\pi(h),\pi'(h')$ gives a pointwise shift for $\mu_0$. So
let $\mu=\mu_0$. Then $T_T=T(z)T(z')>0$, and we use the following facts about the canonical section, valid on $A$.
\begin{enumerate}[label=(\alph*)]
\item (Datum constancy, \lref{Formalism.FourierDecayAlt.P.CutMP.JA.SSn.Rep.Core.TC.PR.datum_constancy_holds}) $\nu_0,\nu'_0,
  \iota_0$ are constant on $A$; and the fibres $\mathrm{Fib}(\zeta)$ do not depend on $\zeta\in A$
  (\lref{Formalism.FourierDecayAlt.P.TheSecStab.theSection_csec}).
\item ($\kappa$-line, \lref{Formalism.FourierDecayAlt.P.CutMP.JA.SSn.Rep.Core.TC.PR.kappa_line_holds}) There are projections
  $Q_\sigma$ on $W_q$, attached to the types and forming a $\kappa$-point of $d$, with $\mathrm{Pr}_\sigma(\nu_0(\zeta))X_f=
  X_fQ_\sigma$ for $X_f\in\mathrm{Fib}(\zeta)$; and by \lref{Formalism.FourierDecayAlt.P.CutMP.Gauge.GC.gauge_kappa_line_holds},
  at a point $\zeta_1\in A$, for every type $\sigma$ there is $E_\sigma\in\mathrm{Fib}(\zeta_1)$ with $E_\sigma Q_\sigma=
  E_\sigma$ and $\|E_\sigma\|_{\mathrm{HS}}=1$ which spans $\{X\in\mathrm{Fib}(\zeta_1):XQ_\sigma=X\}$.
\item (Fibre stability, \lref{Formalism.FourierDecayAlt.P.CutMP.JA.SSn.Rep.Core.TC.PR.fibre_stability_holds}) For
  $X_f\in\mathrm{Fib}(\zeta)$, $\mathrm{Pr}_\sigma(\nu_0(\zeta))X_f$ and $\theta_\zeta X_f$ lie in $\mathrm{Fib}(\zeta)$.
\item (Algebraic shift identity,
  \lref{Formalism.FourierDecayAlt.P.CutMP.JA.SSn.Rep.Core.TC.PR.algebraic_shift_holds}) In the normal-form algebra,
  $\kappa(e_\rho)\,QQ_{10}(Y)\,\kappa(e_\pi)=x_{\mathrm{Ex}}(\{l\},\pi,Y)\,S_l\,\kappa(e_\pi)$ for all $Y\in\C$, where
  $QQ_{10}$ is the $(1,0)$-entry of $QQ$ and $e_\sigma$ is the indicator of the type $\sigma$.
\item (Core polynomiality,
  \lref{Formalism.FourierDecayAlt.P.CutMP.JA.SSn.Rep.Core.TC.PR.corePoly_of}) If a radial function $\Phi$ takes a constant
  value $X_0$ at all points $(c,\mu(\zeta))$, then for every $c$ and $Y$ the function $\zeta\mapsto(t_L(Y)\Phi)(c,\mu(\zeta))$
  is a finite sum $\sum_ia_i(\zeta)v_i$ with fixed vectors $v_i$ and functions $a_i$ satisfying $g^{m_i}a_i=p_i(\zeta,\bar\zeta,
  \varsigma)$ on $A$ for some $m_i$ and polynomials $p_i$ (notation of \cref{pf:thmS-carrier:phaseextract}); here $t_L(Y)$ is
  the radial operator in the last variable.
\item (Phase invariance,
  \lref{Formalism.FourierDecayAlt.P.CutMP.JA.SSn.Rep.Core.TC.PR.phase_invariance_holds}) For every $\theta\in\R$ there is
  $c_\theta\in\C$ such that for all $\zeta\in A$ and all $X_f\in\mathrm{Fib}(\zeta)$ with $\mathrm{Pr}_\pi X_f=X_f$, both
  $\mathrm{Pr}_\rho S(e^{i\theta}\zeta)X_f=c_\theta\mathrm{Pr}_\rho S(\zeta)X_f$ and $\mathrm{Pr}_\rho\theta_{e^{i\theta}\zeta}X_f
  =c_\theta\mathrm{Pr}_\rho\theta_\zeta X_f$ (projections at $\nu_0$ of the respective points).
\item (Non-vanishing,
  \lref{Formalism.FourierDecayAlt.P.CutMP.JA.SSn.Rep.Core.TC.PR.word_nonvanishing_holds}) If $0\ne X_f\in\mathrm{Fib}(\zeta)$
  and $\mathrm{Pr}_\pi X_f=X_f$, then $\mathrm{Pr}_\rho\theta_\zeta X_f\ne0$.
\end{enumerate}
Here (e) is deduced in \lref{Formalism.FourierDecayAlt.P.CutMP.JA.SSn.Rep.Core.TC.PR.corePoly_of} from the radial
decomposition and the polynomiality of word vectors along cuts.

\emph{A line through $E_\pi$.} Fix $\zeta_1=T_T+1\in A$. By (a), $\mathrm{Fib}(\zeta)=\mathrm{Fib}(\zeta_1)$ for all
$\zeta\in A$; so $E_\pi,E_\rho\in\mathrm{Fib}(\zeta)$ and, by (b), $\mathrm{Pr}_\pi(\nu_0(\zeta))E_\pi=E_\pi Q_\pi=E_\pi$.
Since $E_\pi$ has norm $1$, $E_\pi\ne0$. By the surjectivity of the radical reduction
(\lref{Formalism.FourierDecayAlt.P.radical_reduction_surjective}) applied to the constant section $E_\pi$, there is a radial
function $\Phi$ with $\Phi(c,\mu(\zeta))=E_\pi$ for all $c\in(\C^\times)^r$ and $\zeta\in A$; since the $\kappa$-point
$(Q_\sigma)$ acts on radial functions pointwise by right multiplication, $\kappa(e_\pi)\Phi$ has the same values, and by
injectivity of the radical reduction (\lref{Formalism.FourierDecayAlt.P.radical_reduction_injective}),
$\kappa(e_\pi)\Phi=\Phi$.

\emph{The shift of $E_\pi$.} Choose $Y=1+\sum_i|\rho_\pi(i)|\in\R$; then each factor $Y+\rho_\pi(i)$ of
$x_{\mathrm{Ex}}(\{l\},\pi,Y)$ is positive, so $x:=x_{\mathrm{Ex}}(\{l\},\pi,Y)\ne0$. Apply (d) to $\Phi$; since
$X.QQ(Y)=d.QQ(Y)$ and $\kappa(e_\pi)\Phi=\Phi$, this gives $\kappa(e_\rho)(X.t(Y)\Phi)=x\,S_l\Phi$, where $X.t(Y)$ is the
$(1,0)$-entry of $X.QQ(Y)$. Evaluating at $(c_0,\mu(\zeta))$ with $c_0=(1,\dots,1)$, and using that $S$ is a pointwise shift,
that $X.t(Y)$ acts on radial functions by the radial operator $t_L(Y)$, and that $\kappa(e_\rho)$ acts by right
multiplication with $Q_\rho$,
\[
x\,S(\zeta)E_\pi=\bigl(t_L(Y)\Phi\bigr)(c_0,\mu(\zeta))\,Q_\rho .
\]
The values of radial functions lie in the fibres, so $S(\zeta)E_\pi\in\mathrm{Fib}(\zeta)$, and by (b) and $Q_\rho^2=Q_\rho$,
$\mathrm{Pr}_\rho S(\zeta)E_\pi=x^{-1}(t_L(Y)\Phi)(c_0,\mu(\zeta))Q_\rho$. By (c), $\mathrm{Pr}_\rho S(\zeta)E_\pi$ lies in the
line spanned by $E_\rho$, so $\mathrm{Pr}_\rho S(\zeta)E_\pi=F(\zeta)E_\rho$ with
$F(\zeta)=\langle\mathrm{Pr}_\rho S(\zeta)E_\pi,E_\rho\rangle_{\mathrm{HS}}$. Applying the linear functional
$X\mapsto x^{-1}\langle XQ_\rho,E_\rho\rangle_{\mathrm{HS}}$ to the finite sum of (e) shows that $F$ is a linear combination of
functions $a_i$ as in (e); since $g=\zeta\bar\zeta-T_T$ is a polynomial in $\zeta,\bar\zeta$, such linear combinations again
satisfy $g^mF=p(\zeta,\bar\zeta,\varsigma)$ for some $m$ and $p$.

\emph{The theta term.} Let $T_1=\mathrm{Pr}_\rho(\nu_0(\zeta_1))\theta_{\zeta_1}E_\pi$. By (g), $T_1\ne0$; by (c) and (b),
$T_1=t_0E_\rho$ with $t_0=\langle T_1,E_\rho\rangle_{\mathrm{HS}}\ne0$. By (a), $\mathrm{Pr}_\rho(\nu_0(\zeta))\theta_\zeta
E_\pi=T_1$ for all $\zeta\in A$.

\emph{Rotation invariance.} Let $\theta\in\R$ and $c_\theta$ be as in (f). For $\zeta_1$ and $X_f=E_\pi$, the second identity
of (f) and the constancy (a) give $T_1=c_\theta T_1$, hence $c_\theta=1$. The first identity of (f) then gives
$\mathrm{Pr}_\rho S(e^{i\theta}\zeta)E_\pi=\mathrm{Pr}_\rho S(\zeta)E_\pi$, so $F(e^{i\theta}\zeta)=F(\zeta)$ on $A$.

\emph{Conclusion.} By \cref{pf:thmS-carrier:phaseextract} there are $m$ and $q_1$ with $(|\zeta|^2-T_T)^mF(\zeta)=
q_1(\sqrt{|\zeta|^2-T_T})$ on $A$. Put $n=m$ and $q=t_0^{-1}q_1$. Let $\zeta\in A$ and $X_f\in\mathrm{Fib}(\zeta)$ with
$\mathrm{Pr}_\pi X_f=X_f$. Then $X_fQ_\pi=X_f$ by (b), and $X_f\in\mathrm{Fib}(\zeta_1)$, so $X_f=\xi_0E_\pi$ for some $\xi_0\in\C$ by
(b). Therefore
\[
(|\zeta|^2-T_T)^n\mathrm{Pr}_\rho S(\zeta)X_f=\xi_0\,q_1(\varsigma(\zeta))E_\rho=q(\varsigma(\zeta))\,\xi_0\,T_1
=q(\varsigma(\zeta))\,\mathrm{Pr}_\rho\theta_\zeta X_f,
\]
which is \eqref{eq:pf:thmS-carrier:polyradial}.
\end{proof}

\subsubsection{Proof of the model-place theorem}

\begin{proof}[Proof of \cref{pf:thmS-carrier:modelplace}]
By \cref{pf:thmS-carrier:openQQ} it suffices to verify (T1)--(T6). The formal proof does this through a chain of one-step
reductions, which we summarise.
\begin{itemize}
\item (T1) is \lref{Formalism.FourierDecayAlt.P.CutMP.jfSmoothBuilt_of_chartEqs}, which deduces the smoothness of the
  Bessel transforms on $\Omega_v$ at built places from the chart equations
  \lref{Formalism.FourierDecayAlt.P.CutMP.SmoothTransferRegularity.chart_equations_holds}.
\item (T2): \lref{Formalism.FourierDecayAlt.P.CutMP.Gauge.GC.genTransfer_of_rotK} derives a general transfer statement
  from the rotation $K$-invariance of the transfer
  \lref{Formalism.FourierDecayAlt.P.CutMP.Gauge.GC.rotation_K_invariance_holds}, and
  \lref{Formalism.FourierDecayAlt.P.CutMP.Gauge.GC.gCore_of_genTransfer} derives (T2) from it.
\item (T3) and (T4) are \lref{Formalism.FourierDecayAlt.P.CutMP.energy_trace_identity_holds} and
  \lref{Formalism.FourierDecayAlt.P.CutMP.energy_reality_holds}.
\item (T5) is deduced from \cref{pf:thmS-carrier:hatreal} by
  \lref{Formalism.FourierDecayAlt.P.CutMP.jDecayTrQQ_of_open}: the realisation of the hat operators bounds the products
  of hat operators applied to $\mathrm{tr}_e(J_{f_v})$ by transfers of derivatives of $J_{f_v}$.
\item (T6) is deduced by \lref{Formalism.FourierDecayAlt.P.CutMP.JA.jAdjTrQQ_of_three} from \cref{pf:thmS-carrier:hatreal},
  from \cref{pf:thmS-carrier:startransfer}, and from the realisation of every element of $\mathcal B(L)$ by a single operator
  of order $0$ preserving the positive cone together with its adjoint. The last statement is
  \lref{Formalism.FourierDecayAlt.P.CutMP.JA.jABddOpQQ_of_ssnRep}, which derives it from the representation of the
  shift operators $S_l$ by fibre maps (the second-shift representation); this representation is
  \lref{Formalism.FourierDecayAlt.P.CutMP.JA.SSn.Rep.ssnRep_of_polyS} applied to \cref{pf:thmS-carrier:polyradial}.
\end{itemize}
\end{proof}

% ---------------------------------------------------------------------------
\subsection{The energy limit}\label{sec:pf:thmS-carrier:LE}

This subsection proves the model lemma (\cref{pf:thmS-carrier:lemmamodel}). It is deduced from an abstract statement about
regularised energies (\cref{pf:thmS-carrier:LE}), which we prove in full.

\subsubsection{Energy systems}
\begin{definition}[Bases, logarithmic derivatives, cuts]\label{pf:thmS-carrier:base}
A \emph{base} $B=(\Omega,\mu,G)$ is a measure space with a measurable function $G\colon\Omega\to(0,\infty)$. For a finite set
$A$ let $\mathrm{Fn}(A)$ be the space of functions $u\colon\Omega\times A\to\C$. For $\varphi\colon\R\to\R$ put
$(\varphi\cdot u)(\omega,a)=\varphi(G(\omega))u(\omega,a)$; for weights $d\colon\Omega\times A\to\R$ put
\[
\mathrm{sq}_d(u)(\omega)=\sum_ad(\omega,a)|u(\omega,a)|^2,\qquad
\mathrm{pr}_d(u,u')(\omega)=\sum_ad(\omega,a)\,u(\omega,a)\overline{u'(\omega,a)} .
\]
The \emph{logarithmic derivatives} of $\varphi$ are $\thmScarrierld_p\varphi(x)=\frac{d^p}{dy^p}\varphi(e^y)|_{y=\log x}$
($x>0$); $\varphi$ is \emph{log-smooth} if $y\mapsto\varphi(e^y)$ is $C^\infty$. For an operator $T\colon\mathrm{Fn}(A)\to
\mathrm{Fn}(A')$ and $u\in\mathrm{Fn}(A)$ put $T_s(u)=G^{-s}\,T(G^s\cdot u)$ and $d^T_p(u)=\partial_s^pT_s(u)|_{s=0}$. $T$ is
\emph{jet-local} at $u$ if there are $m$ and $c_0,\dots,c_m\in\mathrm{Fn}(A')$ with $T(\varphi\cdot u)=\sum_{q\le m}
\thmScarrierld_q\varphi(G)\,c_q$ for all log-smooth $\varphi$, and \emph{polynomial of degree $\le d$} at $u$ if
$T_s(u)=\sum_{q\le d}s^qc'_q$ for all $s\in\R$. A \emph{cut} is a smooth $\chi\colon\R\to[0,1]$ with $\chi=0$ on
$(-\infty,\frac12]$ and $\chi=1$ on $[1,\infty)$; put $\chi_\eps(x)=\chi(x/\eps)$. The \emph{regularisers} are
\[
\theta_{\eps,\sigma}(x)=\Bigl(\frac{x}{x+\eps}\Bigr)^{\sigma},\qquad\text{so that }\theta_{\eps,\sigma}(e^y)=
(1+\eps e^{-y})^{-\sigma}.
\]
\formal{Formalism.FourierDecayAlt.LE.Base}
\formal{Formalism.FourierDecayAlt.LE.Cut}
\end{definition}

For $\sigma\ge0$, $\eps>0$ and $x>0$ one has $0<\theta_{\eps,\sigma}(x)\le1$; $\theta_{\eps,\sigma}$ is log-smooth;
$\theta_{\eps,\sigma}(x)$ increases to $1$ as $\eps\downarrow0$; and $\thmScarrierld_p\theta_{\eps,\sigma}(x)\to0$ for
$p\ge1$.

\begin{definition}[Steps]\label{pf:thmS-carrier:step}
A \emph{step} $S$ over $B$ consists of a finite set $\iota$, an integer $r$, $X\in\mathrm{Fn}(\iota)$, weights $d_w$ on
$\Omega\times\iota$, for each $Y\in\R$ operators $g_Y\colon\mathrm{Fn}(\iota)\to\mathrm{Fn}(\{0,1\}\times\iota)$ and $g_Y^*$
in the opposite direction, weights $d_Y$ on $\Omega\times\{0,1\}\times\iota$ and matrices $e_Y(\omega)\in M_\iota(\C)$, a
function $P\colon\R\to\R$, and numbers $a,L$; let $J=[a,a+L]$. Write $g_{Y,s}=(g_Y)_s(X)$, $\delta_{Y,p}=d^{g_Y}_p(X)$, and
\begin{gather*}
\mathrm{En}_\eps(Y)=\int\theta_{\eps,\sigma}(G)\,\mathrm{sq}_{d_Y}(g_YX)\,d\mu,\qquad
\mathrm{Ens}_\eps(Y,s)=\int\theta_{\eps,\sigma}(G)\,\mathrm{sq}_{d_Y}(g_{Y,s})\,d\mu,\\
\mathrm{Dn}_\eps(Y,p)=\int\theta_{\eps,\sigma}(G)\,\mathrm{sq}_{d_Y}(\delta_{Y,p})\,d\mu,\qquad
\mathrm{Esup}_\eps=\sup_{Y\in J}\frac{\mathrm{En}_\eps(Y)}{|P(Y)|},\\
A_0=\int\mathrm{sq}_{d_w}(X)\,d\mu,\qquad A_\eps=\int\theta_{\eps,\sigma}(G)\,\mathrm{sq}_{d_w}(X)\,d\mu,\\
\mathrm{term}_\eps(Y,p)=\frac1{p!}\int\thmScarrierld_p\theta_{\eps,\sigma}(G)\,\mathrm{pr}_{d_Y}(g_YX,\delta_{Y,p})\,d\mu
\end{gather*}
(the dependence on $\sigma$ is suppressed). For a cut $\chi$ and $\sigma,C_R,C_d,C_P\ge0$, the \emph{step hypotheses} with
constants $(\chi,\sigma,C_R,C_d,C_P)$ are: measurability of all functions involved; $d_w>0$, and $d_Y>0$ for $Y\in J$; $P$
continuous and non-vanishing on $J$; and, for $Y\in J$ (and $Y+s\in J$ where it occurs):
\begin{enumerate}[label=(S\arabic*)]
\item\label{pf:thmS-carrier:S3} $g_Y^*(d_Y\,g_YX)=|P(Y)|\,d_w\,X$;
\item\label{pf:thmS-carrier:SAdj} $\int\mathrm{pr}_1(g_Y^*(d_Yg_YX),\varphi\cdot X)\,d\mu=\int\mathrm{pr}_{d_Y}(g_YX,g_Y(\varphi\cdot
  X))\,d\mu$ for $\varphi=\chi_{\eps_2}\theta_{\eps,\sigma}$, $\eps,\eps_2>0$;
\item\label{pf:thmS-carrier:SJet} $g_Y$ is jet-local at $X$ and polynomial of degree $\le r+1$ at $X$;
\item\label{pf:thmS-carrier:SInt} $\theta_{\eps,\sigma}(G)\mathrm{sq}_{d_Y}(g_{Y,s})$ is integrable for all $s$, and
  $Y\mapsto\mathrm{En}_\eps(Y)$ is continuous on $J$;
\item\label{pf:thmS-carrier:SShift} $g_{Y,s}(0,i)=g_{Y+s}X(0,i)$ and $g_{Y,s}(1,i)=g_{Y+s}X(1,i)+\sum_j(e_Y-e_{Y+s})_{ij}\,
  g_{Y+s}X(0,j)$;
\item\label{pf:thmS-carrier:SRes} $\sum_id_Y(1,i)\bigl|\sum_j(e_Y-e_{Y+s})_{ij}v_j\bigr|^2\le C_R^2\sum_id_Y(0,i)|v_i|^2$ for
  all $v\in\C^\iota$;
\item\label{pf:thmS-carrier:SRat} $d_Y\le C_d\,d_{Y+s}$ and $|P(Y+s)|\le C_P|P(Y)|$.
\end{enumerate}
Put $C_7=C_dC_P(1+C_R)^2$.
\formal{Formalism.FourierDecayAlt.LE.Step}
\end{definition}

\begin{definition}[Families and systems]\label{pf:thmS-carrier:sys}
A \emph{family} $F$ in $N$ variables consists of finite sets $R,\iota$, an integer $r$, $u\in\mathrm{Fn}(R)$, operators $g^F_{Y}$
($Y\in\R^N$), a weight $W_t>0$ on $\Omega$, numbers $c_Y(i)$, monomial vectors $\mathrm{mon}(Y)\in\mathrm{Fn}(\iota)$,
matrices $B_Y(\omega)$ and sets $J_a\subset\R$; put $d^F_Y(\omega,i)=W_t(\omega)^Nc_Y(i)$. The \emph{family hypotheses} are:
measurability; $c_Y>0$ for $Y\in\prod J_a$; the factorisation $(g^F_Y)_s(u)=B_Y\,\mathrm{mon}(Y+s)$; boundedness of $B_Y$ and
of an inverse; $\mathrm{mon}$ polynomial of degree $\le r+1$ in each variable; jet-locality of $g^F_Y$ at $u$.

A \emph{system} $\mathcal S$ consists of a base $B$, a cut $\chi$, integers $k,r$, numbers $\lambda_1,\dots,\lambda_{k+2r+1}$,
an interval $[\alpha,\beta]$, an integer $N_{\mathrm{tot}}$, a \emph{box length} $L$, an \emph{energy function}
$\thmScarrierEn(\eps,w)$ of words $w$ (lists of letters $(b,Y)\in\{0,1\}\times\R$), steps $S_w$ for lists $w\in\R^*$ and
families $F_N$, such that every $S_w$ has the same $r$, $P(Y)=\prod_s(Y+\lambda_s)$, $a_w=\alpha+\frac{\beta-\alpha}3
+2\len(w)L$ and length $L$; the families have the same $r$ and $J_a=\mathrm{box}(N-1-a)$, where
\[
\mathrm{box}(a)=\Bigl[\alpha+\tfrac{\beta-\alpha}3+2aL,\ \alpha+\tfrac{\beta-\alpha}3+2aL+L\Bigr];
\]
and the densities match: $\mathrm{sq}_{d_w}(X)$ for $S_{(Y)\cdot w}$ equals $\mathrm{sq}_{d_Y}(g_YX)$ for $S_w$, and
$\mathrm{sq}_{d_w}(X)$ for $S_{(Y_1,\dots,Y_N)}$ equals $\mathrm{sq}_{d^F_Y}(g^F_Yu)$ for $F_N$.

For a system put $\mathrm{dens}(w)=\mathrm{sq}_{d_w}(X)$ for $S_w$, $j_{\mathcal S}=\int\mathrm{dens}(\emptyset)\,d\mu$,
$\mathrm{size}(w)=\prod_{(b,Y)\in w}|P(Y)|$ and $\mathrm{plus}(Y_1,\dots,Y_N)=((1,Y_1),\dots,(1,Y_N))$. A list
$(Y_1,\dots,Y_N)$ is \emph{boxed} if $Y_j\in\mathrm{box}(N-j)$ for all $j$; a word is \emph{admitted} if its letters lie in
$[\alpha,\beta]$ and are pairwise at distance $>2$. The conditions on a system are:
\begin{description}
\item[(E1) roots outside] $-\lambda_s\notin[\alpha,\beta]$ for all $s$;
\item[(E2) polynomiality] for $\eps>0$ and $N\le N_{\mathrm{tot}}$ there is a polynomial $Q_\eps\in\R[Y_1,\dots,Y_N]$ of
  degree $\le k+2r+1$ in each variable with $\thmScarrierEn(\eps,((b_j,Y_j))_j)=Q_\eps(Y)$ for every admitted word, whatever
  the signs $b_j$;
\item[(E3) base integrability] $\mathrm{dens}(\emptyset)$ is integrable;
\item[(E4) representation] for boxed $Y$ of length $N\le N_{\mathrm{tot}}$ and $\eps>0$,
  $\thmScarrierEn(\eps,\mathrm{plus}(Y))=\int\mathrm{sq}_{d^F_Y}(g^F_Y(\chi_\eps\cdot u))\,d\mu$ for $F_N$;
\item[(E5) step hypotheses with constants $(\sigma,C_R,C_d,C_P)$] the steps $S_w$ of boxed $w$ with $\len(w)<N_{\mathrm{tot}}$
  satisfy the step hypotheses with constants $(\chi,\sigma,C_R,C_d,C_P)$;
\item[(E6) family hypotheses] the families $F_N$, $N\le N_{\mathrm{tot}}$, satisfy the family hypotheses;
\item[(EL) energy limit] for every admitted $w$ with $\len(w)\le N_{\mathrm{tot}}$, $\thmScarrierEn(\eps,w)\to\mathrm{size}(w)\,
  j_{\mathcal S}$ as $\eps\to0^+$.
\end{description}
\formal{Formalism.FourierDecayAlt.LE.Fam}
\formal{Formalism.FourierDecayAlt.LE.Sys}
\formal{Formalism.FourierDecayAlt.LE.Sys.RootsOut}
\formal{Formalism.FourierDecayAlt.LE.Sys.EnergyLimit}
\end{definition}

\begin{theorem}[Abstract energy limit]\label{pf:thmS-carrier:LE}
Let $C_R,C_d,C_P\ge0$, $C_2\ge64(1+2e\sqrt{C_7})$ and $C_6\ge6C_2$. Let $\mathcal S$ be a system with $r+1\le N_{\mathrm{tot}}$,
$L=C_2N_{\mathrm{tot}}(r+1)^2$ and $C_6\bigl(kN_{\mathrm{tot}}(r+1)^2+N_{\mathrm{tot}}^2(r+1)^2\bigr)\le\beta-\alpha$, satisfying
(E1)--(E4), (E5) with constants $(4N_{\mathrm{tot}},C_R,C_d,C_P)$, and (E6). Then the energy limit (EL) holds.
\formal{Formalism.FourierDecayAlt.LE.LemmaE}
\formallabel{Formalism.FourierDecayAlt.P.le_E_proof}
\formallabel{Formalism.FourierDecayAlt.P.LEx2.lemmaE_closed}
\formallabel{Formalism.FourierDecayAlt.P.LEx2.lemmaE}
\formallabel{Formalism.FourierDecayAlt.P.constants_choice}
\end{theorem}

\subsubsection{Two analytic inputs}
\begin{lemma}[Cauchy bound]\label{pf:thmS-carrier:cauchy}
For $\sigma\ge4$, $\eps>0$, $p\in\N$ and $y\in\R$,
\[
\Bigl|\frac{d^p}{dy^p}(1+\eps e^{-y})^{-\sigma}\Bigr|\le e\,p!\,(4\sigma)^p\,(1+\eps e^{-y})^{-\sigma};
\]
equivalently $|\thmScarrierld_p\theta_{\eps,\sigma}(x)|\le e\,p!\,(4\sigma)^p\theta_{\eps,\sigma}(x)$ for $x>0$.
\formallabel{Formalism.FourierDecayAlt.P.cauchy_theta_identity}
\end{lemma}

\begin{proof}
Let $\rho=1/(4\sigma)\le\frac1{16}$ and $F(w)=(1+\eps e^{-w})^{-\sigma}$ (principal branch) on the open set
$U=\{w\in\C:\Re(1+\eps e^{-w})>0\}$, which contains $\R$; $F$ is holomorphic on $U$ and restricts to the given function on $\R$.
For $|w-y|\le\rho$ we have $|e^{-(w-y)}-1|\le2\rho$, hence $\Re(1+\eps e^{-w})\ge1+\eps e^{-y}(1-2\rho)\ge(1+\eps
e^{-y})(1-2\rho)>0$; so the closed disc of radius $\rho$ about $y$ lies in $U$, and on its boundary
\[
|F(w)|=|1+\eps e^{-w}|^{-\sigma}\le(1-2\rho)^{-\sigma}(1+\eps e^{-y})^{-\sigma}\le e\,(1+\eps e^{-y})^{-\sigma},
\]
because $(1-\frac1{2\sigma})^{-\sigma}\le e$ for $\sigma\ge4$. Cauchy's estimate $|F^{(p)}(y)|\le p!\rho^{-p}\max_{|w-y|=\rho}
|F(w)|$ gives the claim.
\end{proof}

\begin{littheorem}[A.~A.~Markov's inequality {\cite[Chapter~5]{BorweinErdelyi1995}}]\label{pf:thmS-carrier:markov1}
If $p$ is a real polynomial of degree $\le d$ with $|p|\le M$ on $[-1,1]$, then $|p'|\le d^2M$ on $[-1,1]$.
\formallabel{Formalism.FourierDecayAlt.P.markov_inequality}
\end{littheorem}

\begin{lemma}[Markov's inequality for vector-valued polynomials]\label{pf:thmS-carrier:markov}
Let $H$ be a real normed space, $q(t)=\sum_{k\le d}t^ka_k$ with $a_k\in H$, $L>0$ and $M\ge0$ with $\|q(t)\|\le M$ for
$t\in[t_0,t_0+L]$. Then $\|q^{(j)}(t)\|\le(2d^2/L)^jM$ for all $j\in\N$ and $t\in[t_0,t_0+L]$.
\formallabel{Formalism.FourierDecayAlt.P.markov_inequality_vector}
\end{lemma}

\begin{proof}
The affine substitution $t=t_0+\frac L2(x+1)$ and \cref{pf:thmS-carrier:markov1} give $|p'|\le(2d^2/L)M$ on $[t_0,t_0+L]$ for
real polynomials of degree $\le d$ bounded by $M$ there. For $H$-valued $q$ and a continuous linear functional $\ell$ of norm
$\le1$, $\ell\circ q$ is such a polynomial, so $|\ell(q'(t))|\le(2d^2/L)M$; by the Hahn--Banach theorem
$\|q'(t)\|\le(2d^2/L)M$. Since $q'$ again has degree $\le d$, induction on $j$ gives the claim.
\end{proof}

\begin{lemma}[Leibniz formula]\label{pf:thmS-carrier:leibniz}
If $T$ is jet-local and polynomial of degree $\le d$ at $u$, then for every log-smooth $\varphi$
\[
T(\varphi\cdot u)=\sum_{p\le d}\thmScarrierld_p\varphi(G)\,\frac{d^T_p(u)}{p!}.
\]
\formallabel{Formalism.FourierDecayAlt.P.LEx2.leibniz}
\end{lemma}

\begin{proof}
For $\varphi_s(x)=x^s$ we have $\varphi_s(e^y)=e^{sy}$, so $\thmScarrierld_q\varphi_s(x)=s^qx^s$. Jet-locality gives
$T_s(u)=G^{-s}\sum_{q\le m}s^qG^sc_q=\sum_{q\le m}s^qc_q$. Pointwise in $(\omega,a)$ this polynomial in $s$ coincides with
$\sum_{q\le d}s^qc'_q$; comparing coefficients, $c_q=c'_q$ (with $c'_q=0$ for $q>d$ and $c_q=0$ for $q>m$), and
$d^T_p(u)=p!\,c'_p$. Hence $T(\varphi\cdot u)=\sum_q\thmScarrierld_q\varphi(G)c_q=\sum_{p\le d}\thmScarrierld_p\varphi(G)
d^T_p(u)/p!$.
\end{proof}

\subsubsection{The step claim}
In this subsection $S$ is a step satisfying the step hypotheses with constants $(\chi,\sigma,C_R,C_d,C_P)$, where $\sigma\ge4$,
$C_R,C_d,C_P\ge0$ and $L>0$; $\mathrm{sq}_{d_w}(X)$ is integrable; $\eps>0$ and $Y\in J$. Note that $\mathrm{Esup}_\eps<\infty$,
since $\mathrm{En}_\eps$ is continuous and $P$ is continuous and non-vanishing on the compact interval $J$.

\begin{lemma}[Energy identity]\label{pf:thmS-carrier:star}
For every $p$ the function $\theta_{\eps,\sigma}(G)\mathrm{sq}_{d_Y}(\delta_{Y,p})$ is integrable, and
$|P(Y)|\,A_\eps=\sum_{p=0}^{r+1}\mathrm{term}_\eps(Y,p)$, with $\mathrm{term}_\eps(Y,0)=\mathrm{En}_\eps(Y)$.
\formallabel{Formalism.FourierDecayAlt.P.LEx2.adj}
\formallabel{Formalism.FourierDecayAlt.P.LEx2.coefIntegrable}
\formallabel{Formalism.FourierDecayAlt.P.LEx2.cutTheta}
\end{lemma}

\begin{proof}
By polynomiality \ref{pf:thmS-carrier:SJet} the coefficients $\delta_{Y,p}$ are linear combinations of finitely many values
$g_{Y,s}$ (Lagrange interpolation in $s$), so the integrability follows from \ref{pf:thmS-carrier:SInt}. The identity for
$\mathrm{term}_\eps(Y,0)$ holds because $\thmScarrierld_0\theta=\theta$ and $\delta_{Y,0}=g_YX$. For $\eps_2>0$ let
$\varphi_{\eps_2}=\chi_{\eps_2}\theta_{\eps,\sigma}$, a log-smooth function. By \ref{pf:thmS-carrier:S3} the left side of
\ref{pf:thmS-carrier:SAdj} is $\int\varphi_{\eps_2}(G)|P(Y)|\mathrm{sq}_{d_w}(X)\,d\mu$, and by \cref{pf:thmS-carrier:leibniz}
(with $d=r+1$) the right side is $\sum_{p\le r+1}\frac1{p!}\int\thmScarrierld_p\varphi_{\eps_2}(G)\,\mathrm{pr}_{d_Y}(g_YX,
\delta_{Y,p})\,d\mu$. Let $\eps_2\to0^+$. For fixed $x>0$, $\varphi_{\eps_2}=\theta_{\eps,\sigma}$ near $x$ once $\eps_2<x$, so
$\thmScarrierld_p\varphi_{\eps_2}(x)\to\thmScarrierld_p\theta_{\eps,\sigma}(x)$. Moreover
$|\thmScarrierld_p\varphi_{\eps_2}(x)|\le C\,\theta_{\eps,\sigma}(x)$ with $C$ independent of $\eps_2$: the logarithmic
derivatives of $\chi_{\eps_2}$ are those of $\chi$ at $x/\eps_2$, hence bounded uniformly, those of $\theta$ are bounded by
\cref{pf:thmS-carrier:cauchy}, and the Leibniz rule for $\frac{d^p}{dy^p}$ applies. Since $\theta\,\mathrm{sq}_{d_w}(X)$ and
$\theta\,|\mathrm{pr}_{d_Y}(g_YX,\delta_{Y,p})|\le\frac12\theta(\mathrm{sq}_{d_Y}(g_YX)+\mathrm{sq}_{d_Y}(\delta_{Y,p}))$ are
integrable, dominated convergence gives the identity in the limit.
\end{proof}

\begin{lemma}[Bounds for the terms]\label{pf:thmS-carrier:terms}
\begin{enumerate}
\item $|\mathrm{term}_\eps(Y,p)|\le e\,(4\sigma)^p\sqrt{\mathrm{En}_\eps(Y)}\sqrt{\mathrm{Dn}_\eps(Y,p)}$.
\item If $Y+s\in J$, then $\mathrm{Ens}_\eps(Y,s)\le C_7\,\mathrm{Esup}_\eps\,|P(Y)|$.
\item If $\mathrm{Ens}_\eps(Y,s)\le M$ whenever $Y+s\in J$, then $\mathrm{Dn}_\eps(Y,p)\le(2(r+1)^2/L)^{2p}M$.
\item $|\mathrm{term}_\eps(Y,p)|\le e\sqrt{C_7}\,\mathrm{Esup}_\eps\,|P(Y)|\,(8\sigma(r+1)^2/L)^p$.
\end{enumerate}
\formallabel{Formalism.FourierDecayAlt.P.LEx2.csTerm}
\formallabel{Formalism.FourierDecayAlt.P.LEx2.A.shiftBound}
\formallabel{Formalism.FourierDecayAlt.P.LEx2.A.shift_pt}
\formallabel{Formalism.FourierDecayAlt.P.LEx2.markovTerm}
\formallabel{Formalism.FourierDecayAlt.P.LEx2.A.termBound}
\end{lemma}

\begin{proof}
(1) By \cref{pf:thmS-carrier:cauchy}, $|\thmScarrierld_p\theta|\le e\,p!(4\sigma)^p\theta$; pointwise
$|\mathrm{pr}_d(u,v)|\le\mathrm{sq}_d(u)^{1/2}\mathrm{sq}_d(v)^{1/2}$ for $d\ge0$; then apply the Cauchy--Schwarz inequality in
$L^2(\theta(G)\,d\mu)$.

(2) Fix $\omega$ and write $u_b(i)=g_{Y+s}X(\omega,(b,i))$, $t_i=\sum_j(e_Y-e_{Y+s})_{ij}u_0(j)$, $A=\sum_id_Y(0,i)|u_0(i)|^2$,
$B=\sum_id_Y(1,i)|u_1(i)|^2$, $T=\sum_id_Y(1,i)|t_i|^2$ and $V=\sum_id_Y(1,i)|u_1(i)+t_i|^2$. By \ref{pf:thmS-carrier:SShift},
$\mathrm{sq}_{d_Y}(g_{Y,s})(\omega)=V+A$, and $T\le C_R^2A$ by \ref{pf:thmS-carrier:SRes}. If $C_R=0$ then $t=0$ (the weights
are positive) and $V=B$. If $C_R>0$, the inequality $|x+y|^2\le(1+C_R)|x|^2+(1+C_R^{-1})|y|^2$ gives
$V\le(1+C_R)B+(1+C_R)C_RA$. In both cases $V+A\le(1+C_R)B+((1+C_R)C_R+1)A\le(1+C_R)^2(B+A)$, and by
\ref{pf:thmS-carrier:SRat} $B+A\le C_d\,\mathrm{sq}_{d_{Y+s}}(g_{Y+s}X)(\omega)$. Multiplying by $\theta(G)$ and integrating,
$\mathrm{Ens}_\eps(Y,s)\le(1+C_R)^2C_d\,\mathrm{En}_\eps(Y+s)\le(1+C_R)^2C_d\,\mathrm{Esup}_\eps|P(Y+s)|\le
C_7\mathrm{Esup}_\eps|P(Y)|$.

(3) Consider the Hilbert space $\mathcal H=L^2(\Omega\times\{0,1\}\times\iota,\theta(G)d_Y\,d\mu)$. By
\ref{pf:thmS-carrier:SJet} and \cref{pf:thmS-carrier:star}, $s\mapsto g_{Y,s}$ is a polynomial of degree $\le r+1$ with
coefficients in $\mathcal H$, with $\|g_{Y,s}\|^2_{\mathcal H}=\mathrm{Ens}_\eps(Y,s)\le M$ on the interval $\{s:Y+s\in J\}$ of
length $L$, which contains $0$; its $p$-th derivative at $0$ is $\delta_{Y,p}$. \Cref{pf:thmS-carrier:markov} gives
$\|\delta_{Y,p}\|_{\mathcal H}\le(2(r+1)^2/L)^p\sqrt M$.

(4) By definition of $\mathrm{Esup}$, $\mathrm{En}_\eps(Y)\le\mathrm{Esup}_\eps|P(Y)|$; by (2) and (3),
$\mathrm{Dn}_\eps(Y,p)\le(2(r+1)^2/L)^{2p}C_7\mathrm{Esup}_\eps|P(Y)|$; insert both into (1).
\end{proof}

\begin{lemma}[Absorption]\label{pf:thmS-carrier:absorb}
Assume $L\ge16\sigma(r+1)^2(1+2e\sqrt{C_7})$. Then $\mathrm{En}_\eps(Y)/|P(Y)|\le A_0+\frac12\mathrm{Esup}_\eps$ for $Y\in J$;
hence $\mathrm{Esup}_\eps\le2A_0$, and
\[
\mathrm{Dn}_\eps(Y,p)\le2C_7\,(2(r+1)^2/L)^{2p}\,|P(Y)|\,A_0\qquad(Y\in J,\ p\in\N,\ \eps>0).
\]
\formallabel{Formalism.FourierDecayAlt.P.LEx2.A.energyStep}
\formallabel{Formalism.FourierDecayAlt.P.LEx2.A.absorb}
\formallabel{Formalism.FourierDecayAlt.P.LEx2.A.uniformBound}
\formallabel{Formalism.FourierDecayAlt.P.LEx2.A.geomSum}
\end{lemma}

\begin{proof}
By \cref{pf:thmS-carrier:star}, $\mathrm{En}_\eps(Y)=|P(Y)|A_\eps-\Re\sum_{p=1}^{r+1}\mathrm{term}_\eps(Y,p)$, and $A_\eps\le
A_0$ since $\theta\le1$. Let $x=8\sigma(r+1)^2/L$, so $x\le\frac1{2(1+2e\sqrt{C_7})}\le\frac12$. Then $\sum_{p\ge1}x^p\le2x$
and $2xe\sqrt{C_7}\le\frac12$, so $\sum_{p\ge1}e\sqrt{C_7}x^p\le\frac12$, and \cref{pf:thmS-carrier:terms}(4) gives
$\mathrm{En}_\eps(Y)\le|P(Y)|A_0+\frac12\mathrm{Esup}_\eps|P(Y)|$. Taking the supremum over $Y\in J$ and using
$\mathrm{Esup}_\eps<\infty$ gives $\mathrm{Esup}_\eps\le2A_0$. The last bound follows from \cref{pf:thmS-carrier:terms}(2),(3)
with $M=2C_7A_0|P(Y)|$.
\end{proof}

\begin{theorem}[Step claim]\label{pf:thmS-carrier:stepclaim}
Let $S$ be a step satisfying the step hypotheses with constants $(\chi,\sigma,C_R,C_d,C_P)$, where $\sigma\ge4$,
$C_R,C_d,C_P\ge0$, $L>0$ and $L\ge16\sigma(r+1)^2(1+2e\sqrt{C_7})$, and with $\mathrm{sq}_{d_w}(X)$ integrable. Then for every
$Y\in J$ the function $\mathrm{sq}_{d_Y}(g_YX)$ is integrable and $\int\mathrm{sq}_{d_Y}(g_YX)\,d\mu=|P(Y)|\,A_0$.
\formallabel{Formalism.FourierDecayAlt.P.LEx2.stepClaim}
\formallabel{Formalism.FourierDecayAlt.P.LEx2.A.stepClaimFromSteps}
\formallabel{Formalism.FourierDecayAlt.P.LEx2.fatouStep}
\formallabel{Formalism.FourierDecayAlt.P.LEx2.domConv}
\end{theorem}

\begin{proof}
By \cref{pf:thmS-carrier:absorb}, $\int\theta_{\eps,\sigma}(G)\mathrm{sq}_{d_Y}(\delta_{Y,p})\le K_p$ for all $\eps>0$, with $K_p$
independent of $\eps$. Since $\theta_{\eps,\sigma}\uparrow1$ as $\eps\downarrow0$, Fatou's lemma shows that
$\mathrm{sq}_{d_Y}(\delta_{Y,p})$ is integrable for every $p$; for $p=0$ this is the integrability of
$\mathrm{sq}_{d_Y}(g_YX)$. Now let $\eps\to0^+$ in \cref{pf:thmS-carrier:star}: $A_\eps\to A_0$ and
$\mathrm{term}_\eps(Y,0)=\mathrm{En}_\eps(Y)\to\int\mathrm{sq}_{d_Y}(g_YX)$ by monotone convergence; for $p\ge1$,
$\thmScarrierld_p\theta_{\eps,\sigma}\to0$ pointwise with $|\thmScarrierld_p\theta_{\eps,\sigma}|\le e\,p!(4\sigma)^p$
(\cref{pf:thmS-carrier:cauchy}), and $|\mathrm{pr}_{d_Y}(g_YX,\delta_{Y,p})|$ is integrable, so $\mathrm{term}_\eps(Y,p)\to0$ by
dominated convergence. Hence $|P(Y)|A_0=\int\mathrm{sq}_{d_Y}(g_YX)\,d\mu$.
\end{proof}

\subsubsection{Proof of the abstract energy limit}
\begin{lemma}[Induction over boxed words]\label{pf:thmS-carrier:claims}
Under the hypotheses of \cref{pf:thmS-carrier:LE}, for every boxed $w$ with $\len(w)\le N_{\mathrm{tot}}$ the density
$\mathrm{dens}(w)$ is integrable and $\int\mathrm{dens}(w)\,d\mu=\mathrm{size}(\mathrm{plus}(w))\,j_{\mathcal S}$.
\formallabel{Formalism.FourierDecayAlt.P.LEx2.le_claimAll_proof}
\formallabel{Formalism.FourierDecayAlt.P.LEx2.le_claimInduction_proof}
\end{lemma}

\begin{proof}
Induction on $\len(w)$; the empty word is (E3) together with the definition of $j_{\mathcal S}$. Let $(Y)\cdot w$ be boxed with
$\len(w)<N_{\mathrm{tot}}$; then $w$ is boxed and $Y\in\mathrm{box}(\len(w))=[a_w,a_w+L]=J(S_w)$. With $\sigma=4N_{\mathrm{tot}}
\ge4$ we have $16\sigma(r+1)^2(1+2e\sqrt{C_7})=64N_{\mathrm{tot}}(r+1)^2(1+2e\sqrt{C_7})\le C_2N_{\mathrm{tot}}(r+1)^2=L$, so
\cref{pf:thmS-carrier:stepclaim} applies to $S_w$, whose base density $\mathrm{dens}(w)$ is integrable by induction. Since
$\mathrm{dens}((Y)\cdot w)=\mathrm{sq}_{d_Y}(g_YX)$ for $S_w$, it is integrable with
$\int\mathrm{dens}((Y)\cdot w)=|P(Y)|\int\mathrm{dens}(w)=|P(Y)|\,\mathrm{size}(\mathrm{plus}(w))\,j_{\mathcal S}=
\mathrm{size}(\mathrm{plus}((Y)\cdot w))\,j_{\mathcal S}$.
\end{proof}

\begin{lemma}[Limits on boxes]\label{pf:thmS-carrier:boxlimit}
Under the hypotheses of \cref{pf:thmS-carrier:LE}, for boxed $Y$ of length $N\le N_{\mathrm{tot}}$,
$\thmScarrierEn(\eps,\mathrm{plus}(Y))\to\mathrm{size}(\mathrm{plus}(Y))\,j_{\mathcal S}$ as $\eps\to0^+$.
\formallabel{Formalism.FourierDecayAlt.P.LEx2.le_boxLimitFromClaim_proof}
\formallabel{Formalism.FourierDecayAlt.P.LEx2.limitB}
\formallabel{Formalism.FourierDecayAlt.P.LEx2.coefficient_L2}
\formallabel{Formalism.FourierDecayAlt.P.LEx2.limit_L2}
\formallabel{Formalism.FourierDecayAlt.P.LEx2.pairLimit}
\end{lemma}

\begin{proof}
By (E4), $\thmScarrierEn(\eps,\mathrm{plus}(Y))=\int\mathrm{sq}_{d^F_Y}(g^F_Y(\chi_\eps\cdot u))$ for $F=F_N$. The lists
$Y'\in\prod_aJ_a$ are exactly the boxed lists of length $N$, and for them $\mathrm{sq}_{d^F_{Y'}}(g^F_{Y'}u)=\mathrm{dens}(Y')$ is
integrable by \cref{pf:thmS-carrier:claims}. We show $\int\mathrm{sq}_{d^F_Y}(g^F_Y(\chi_\eps\cdot u))\to\int\mathrm{sq}_{d^F_Y}
(g^F_Yu)$; the limit is $\int\mathrm{dens}(Y)=\mathrm{size}(\mathrm{plus}(Y))j_{\mathcal S}$ by \cref{pf:thmS-carrier:claims}.

First, by \lref{Formalism.FourierDecayAlt.P.LEx2.limitA}, $g^F_Y$ is polynomial at $u$ of degree $\le N(r+1)$ and all
$s$-derivatives $d_p=d^{g^F_Y}_p(u)$ satisfy $\mathrm{sq}_{d^F_Y}(d_p)\in L^1$. This uses the factorisation
$(g^F_Y)_s(u)=B_Y\,\mathrm{mon}(Y+s)$, the boundedness of $B_Y$ and of its inverse, and the following interpolation fact: the
coefficients of a family $\mathrm{mon}(Y')=\sum_\kappa Y'^\kappa c_\kappa$, polynomial in each variable, which is weighted-$L^2$
at all points of a product of infinite sets, are weighted-$L^2$ (each coefficient is a fixed linear combination of values on a
finite product grid, by Lagrange interpolation in each variable). By \cref{pf:thmS-carrier:leibniz},
$g^F_Y(\chi_\eps\cdot u)-g^F_Yu=\sum_p(\thmScarrierld_p\chi_\eps(G)-\delta_{p0})d_p/p!$. Here
$\thmScarrierld_p\chi_\eps(x)=(\thmScarrierld_p\chi)(x/\eps)$ is bounded uniformly in $\eps$ and equals $\delta_{p0}$ for
$x>\eps$, so the difference tends to $0$ in the weighted $L^2$-norm by dominated convergence; convergence in $L^2$ implies
convergence of the norms.
\end{proof}

\begin{lemma}[Reduction to boxes]\label{pf:thmS-carrier:reduce}
Assume $L>2$, $6N_{\mathrm{tot}}L\le\beta-\alpha$, (E1), (E2), and the conclusion of \cref{pf:thmS-carrier:boxlimit}. Then the
energy limit (EL) holds.
\formallabel{Formalism.FourierDecayAlt.P.LEx2.le_reduceToBoxes_proof}
\formallabel{Formalism.FourierDecayAlt.P.LEx2.le_boxesFit_proof}
\formallabel{Formalism.FourierDecayAlt.P.LEx2.le_constSign_proof}
\formallabel{Formalism.FourierDecayAlt.P.LEx2.gridConvergence}
\end{lemma}

\begin{proof}
Since $6N_{\mathrm{tot}}L\le\beta-\alpha$, the boxes $\mathrm{box}(a)$, $a<N_{\mathrm{tot}}$, lie in the middle third of
$[\alpha,\beta]$, and for $a<b$ every point of $\mathrm{box}(b)$ exceeds every point of $\mathrm{box}(a)$ by at least $L>2$;
hence $\mathrm{plus}(Y)$ is admitted for every boxed $Y$. By (E1) each factor $Y+\lambda_s$ has constant sign on $[\alpha,
\beta]$, so there is $\mathrm{sg}\in\{\pm1\}$ with $|P(Y)|=\mathrm{sg}\,P(Y)$ there. Let $w=((b_j,Y_j))_{j\le N}$ be admitted,
$N\le N_{\mathrm{tot}}$, and $Q_0=\mathrm{sg}^Nj_{\mathcal S}\prod_j\prod_s(Y_j+\lambda_s)$; then $\mathrm{size}(w')j_{\mathcal
S}=Q_0(Y')$ for every admitted $w'$ with letters $Y'$, and $Q_0$ and the polynomials $Q_\eps$ of (E2) have degree $\le
n:=k+2r+1$ in each variable. Choose $n+1$ distinct points in each $\mathrm{box}(N-j)$. On the resulting grid the lists are
boxed, so by (E2) and \cref{pf:thmS-carrier:boxlimit}, $Q_\eps\to Q_0$ at every grid point. A polynomial of degree $\le n$ in
each of $N$ variables is a fixed linear combination of its values on such a grid (Lagrange interpolation in each variable), so
$Q_\eps(Y)\to Q_0(Y)$ for every $Y\in\R^N$. At the letters of $w$ this reads $\thmScarrierEn(\eps,w)\to\mathrm{size}(w)
j_{\mathcal S}$.
\end{proof}

\begin{proof}[Proof of \cref{pf:thmS-carrier:LE}]
From $r+1\le N_{\mathrm{tot}}$, $C_6\ge6C_2$ and the bound on $\beta-\alpha$ we get $6N_{\mathrm{tot}}L=6C_2N_{\mathrm{tot}}^2
(r+1)^2\le C_6N_{\mathrm{tot}}^2(r+1)^2\le\beta-\alpha$, and $L\ge64>2$. \Cref{pf:thmS-carrier:claims} (which uses
\cref{pf:thmS-carrier:stepclaim}), \cref{pf:thmS-carrier:boxlimit} and \cref{pf:thmS-carrier:reduce} give (EL). The analytic
inputs \cref{pf:thmS-carrier:cauchy,pf:thmS-carrier:markov} enter through \cref{pf:thmS-carrier:stepclaim}.
\end{proof}

\subsubsection{The model lemma}
\begin{definition}[The model system]\label{pf:thmS-carrier:modelsys}
For a parameter $P=(k,r,T,T',u,u')$ let $T_T\ge0$ be its threshold and $Q=Q_{r,T_T}$, with the measure $\mu_0$ of density
$\prod_l|x_l|^{-2}$ and the function $G_q(x,z)=\prod_l|x_l|^2(|z|^2-T_T)>0$; this is the \emph{model base}. The \emph{model
operator algebra} consists of matrices, indexed by the types $\rho$ of $P$, of pairs of operators on the algebra of smooth
functions on $Q$ generated by multiplication by the coordinates $x_l,z$, their conjugates and inverses, and Wirtinger-type
derivatives; it contains the operators $C_0(Y)$ ($Y\in\R$) of the model frame at level $l$. For model data $(L,j,\lambda)$, a
cut $\chi$ and $\eps>0$, the \emph{cut-off energy form} of a word $w$ is
\[
\thmScarrierEn^{\mathrm{mod}}(\eps,w)=\Re\Bigl\langle\prod_{(b,Y)\in w}C_0(Y)\,\chi_\eps(G_q)j,\ \chi_\eps(G_q)j
\Bigr\rangle_{L^2(\mu_0)},
\]
the product over the letters taken in reverse order. With explicit steps and families it is the energy function of a system
$\mathcal S^{\mathrm{mod}}=\mathcal S^{\mathrm{mod}}(P,l,L,j,\lambda,\alpha,\beta,L_b,\chi,N_{\mathrm{tot}})$ with base
$(Q,\mu_0,G_q)$, $P(Y)=\prod_s(Y+\lambda_s)$ and box length $L_b$.
\formal{Formalism.FourierDecayAlt.M.Par}
\formal{Formalism.FourierDecayAlt.MA.baseQ}
\formal{Formalism.FourierDecayAlt.MA.A}
\formal{Formalism.FourierDecayAlt.MA.cutoffEnergyForm}
\formal{Formalism.FourierDecayAlt.MA.sysQ}
\end{definition}

We call the hypotheses on the data in the following three results the \emph{model hypotheses}: the model Hida conditions hold
on $[\alpha,\beta]$ at level $l$, $j$ satisfies the decay and adjointness conditions of the model to order $2N_{\mathrm{tot}}+2$
(\cref{pf:thmS-carrier:hatdef}) and the eigenfunction condition, $r+1\le N_{\mathrm{tot}}\le L_b$, $L_b>2$,
$(k+2r+1)L_b\le\beta-\alpha$ and $6N_{\mathrm{tot}}L_b\le\beta-\alpha$.

\begin{theorem}[Step hypotheses of the model system]\label{pf:thmS-carrier:stepimports}
There are $C_R,C_d\ge0$ such that, with $C_P=e^3$, for every parameter $P$, all data $(l,L,j,\lambda,\alpha,\beta,L_b,\chi,
N_{\mathrm{tot}})$ satisfying the model hypotheses, and every boxed list $w$ with $\len(w)<N_{\mathrm{tot}}$, the step $S_w$ of
$\mathcal S^{\mathrm{mod}}$ satisfies the step hypotheses with constants $(\chi,4N_{\mathrm{tot}},C_R,C_d,C_P)$.
\formallabel{Formalism.FourierDecayAlt.P.induction_imports}
\formallabel{Formalism.FourierDecayAlt.P.induction_imports_of_parts}
\end{theorem}

\begin{proof}
The step hypotheses (\cref{pf:thmS-carrier:step}) consist of twelve conditions, and each is supplied, for all parameters and
all data satisfying the model hypotheses, by one statement about the model steps:
\begin{itemize}
\item measurability: \lref{Formalism.FourierDecayAlt.P.measure_regular};
\item $d_w>0$ and $d_Y>0$ on $J$: \lref{Formalism.FourierDecayAlt.P.weight_positive};
\item $P$ continuous and non-vanishing on $J$: \lref{Formalism.FourierDecayAlt.P.polynomial_regular};
\item the absolute identity \ref{pf:thmS-carrier:S3}: \lref{Formalism.FourierDecayAlt.P.third_identity};
\item the adjointness \ref{pf:thmS-carrier:SAdj} for the cut $\chi$ and $\sigma=4N_{\mathrm{tot}}$:
  \lref{Formalism.FourierDecayAlt.P.adjoint_bounded};
\item jet-locality of $g_Y$ at $X$: \lref{Formalism.FourierDecayAlt.P.jet_module_closed};
\item polynomiality of $g_Y$ of degree $\le r+1$ at $X$: \lref{Formalism.FourierDecayAlt.P.fourth_identity};
\item the integrability and continuity \ref{pf:thmS-carrier:SInt} for $\sigma=4N_{\mathrm{tot}}$:
  \lref{Formalism.FourierDecayAlt.P.fourth_identity_dual};
\item the shift formula \ref{pf:thmS-carrier:SShift}: \lref{Formalism.FourierDecayAlt.P.fifth_identity_first};
\item the residue bound \ref{pf:thmS-carrier:SRes} with a constant $C_R\ge0$: \lref{Formalism.FourierDecayAlt.P.residuesAlt};
\item the weight ratio $d_Y\le C_dd_{Y+s}$ with a constant $C_d\ge0$: \lref{Formalism.FourierDecayAlt.P.weight_ratio};
\item the ratio $|P(Y+s)|\le e^3|P(Y)|$: \lref{Formalism.FourierDecayAlt.P.probability_ratio}.
\end{itemize}
The constants $C_R$ and $C_d$ of the residue and weight-ratio statements do not depend on the data; together with $C_P=e^3$
they give the theorem.
\end{proof}

\begin{theorem}[Family hypotheses of the model system]\label{pf:thmS-carrier:famimports}
For every parameter $P$, all data satisfying the model hypotheses and every cut $\chi$, the families $F_N$
($N\le N_{\mathrm{tot}}$) of $\mathcal S^{\mathrm{mod}}$ satisfy the family hypotheses (E6).
\formallabel{Formalism.FourierDecayAlt.P.family_imports}
\formallabel{Formalism.FourierDecayAlt.P.family_imports_of_parts}
\end{theorem}

\begin{proof}
The family $F_N$ of the model system is defined through a choice of a monomial vector $\mathrm{mon}$ and a matrix $B$ (a
\emph{factorisation datum}) for which the pre-family is measurable, factorises as $(g^F_Y)_s(u)=B_Y\mathrm{mon}(Y+s)$, has $B_Y$
and an inverse bounded, and has $\mathrm{mon}$ polynomial of degree $\le r+1$ in each variable; such a datum exists for every
$N\le N_{\mathrm{tot}}$ by \lref{Formalism.FourierDecayAlt.P.factor_exists}, and the families of the model system are
defined with such a chosen datum, so they have these four properties. The remaining two family hypotheses, positivity $c_Y>0$ on
the boxes and jet-locality of $g^F_Y$ at $u$, are \lref{Formalism.FourierDecayAlt.P.family_restriction}.
\end{proof}

\begin{theorem}[Theorem P for the model system]\label{pf:thmS-carrier:thmPmodel}
For every parameter $P$, all data satisfying the model hypotheses and every cut $\chi$, the model system satisfies (E4): for
every boxed $Y=(Y_1,\dots,Y_N)$ with $N\le N_{\mathrm{tot}}$ and every $\eps>0$,
\[
\thmScarrierEn^{\mathrm{mod}}(\eps,\mathrm{plus}(Y))=\int\mathrm{sq}_{d^F_Y}\bigl(g^F_Y(\chi_\eps\cdot u)\bigr)\,d\mu_0 .
\]
\formallabel{Formalism.FourierDecayAlt.P.theorem_P}
\formallabel{Formalism.FourierDecayAlt.P.theorem_P_of_parts}
\end{theorem}

\begin{proof}
Write $J_\eps=\chi_\eps(G_q)j$, and let $g_e$ ($e$ in the index set of $F_N$) and $d_{ee'}$ be the components of the family
operators and of the weights, as elements of the model operator algebra. The energy of $\mathrm{plus}(Y)$ is the cut-off
energy form (\cref{pf:thmS-carrier:modelsys}). The proof is a chain of three identities, each valid for boxed $Y$ and
$\eps>0$:
\begin{align*}
\thmScarrierEn^{\mathrm{mod}}(\eps,\mathrm{plus}(Y))
&=\Re\Bigl\langle\Bigl(\sum_{e,e'}g_e^*d_{ee'}g_{e'}\Bigr)J_\eps,J_\eps\Bigr\rangle_Q
&&\text{(\lref{Formalism.FourierDecayAlt.P.theorem_P_form})}\\
&=\Re\sum_e\Bigl\langle\sum_{e'}d_{ee'}g_{e'}J_\eps,\ g_eJ_\eps\Bigr\rangle_Q
&&\text{(\lref{Formalism.FourierDecayAlt.P.theorem_P_adjoint})}\\
&=\int\mathrm{sq}_{d^F_Y}\bigl(g^F_Y(\chi_\eps\cdot u)\bigr)\,d\mu_0
&&\text{(\lref{Formalism.FourierDecayAlt.P.theorem_P_weight})}.
\end{align*}
The first identity rewrites the product of the operators $C_0(Y_i)$ in terms of the family operators; the second moves $g_e^*$
to the other side of the model inner product (an adjointness statement for the cut functions); the third identifies the
weighted sum with the integral of the weighted square, which requires the integrability of the terms.
\end{proof}

\begin{theorem}[The model lemma]\label{pf:thmS-carrier:lemmamodel}
There are $C_2,C_6>0$ such that for every parameter $P$, level $l$, level datum $L$, family $j$, numbers $\lambda_s$, interval
$[\alpha,\beta]$, cut $\chi$ and integer $N_{\mathrm{tot}}$, if the model Hida conditions hold on $[\alpha,\beta]$, $j$ satisfies
the decay and adjointness conditions of the model to order $2N_{\mathrm{tot}}+2$ and the eigenfunction condition (for every
$Y\in\R$ the model operator $C_0(Y)$ acts on $j$ as multiplication by $\prod_s(Y+\lambda_s)$), $r+1\le N_{\mathrm{tot}}$ and
$C_6\bigl(kN_{\mathrm{tot}}(r+1)^2+N_{\mathrm{tot}}^2(r+1)^2\bigr)\le\beta-\alpha$, then for every word $w$ admitted by
$\mathcal S^{\mathrm{mod}}$ with box length $L_b=C_2N_{\mathrm{tot}}(r+1)^2$ and $\len(w)\le N_{\mathrm{tot}}$,
\[
\thmScarrierEn^{\mathrm{mod}}(\eps,w)\longrightarrow\mathrm{size}(w)\int_Q\sum_\rho|j_\rho|^2\,d\mu_0\qquad(\eps\to0^+).
\]
\formal{Formalism.FourierDecayAlt.MA.LemmaEModel}
\formallabel{Formalism.FourierDecayAlt.P.lemma_model}
\formallabel{Formalism.FourierDecayAlt.P.lemma_model_of_parts}
\formallabel{Formalism.FourierDecayAlt.P.jacobi_norm}
\end{theorem}

\begin{proof}
By \cref{pf:thmS-carrier:stepimports} there are $C_R,C_d\ge0$ such that, with $C_P=e^3$, for all data satisfying the model
hypotheses the steps of boxed lists of length $<N_{\mathrm{tot}}$ satisfy the step hypotheses with constants
$(\chi,4N_{\mathrm{tot}},C_R,C_d,C_P)$. Put $C_2=64(1+2e\sqrt{C_7})$ and $C_6=6C_2$. Given the data,
$L_b=C_2N_{\mathrm{tot}}(r+1)^2\ge N_{\mathrm{tot}}$ and $L_b\ge64>2$, and elementary inequalities give $(k+2r+1)L_b\le\beta-
\alpha$ and $6N_{\mathrm{tot}}L_b\le\beta-\alpha$; so the model hypotheses hold. \Cref{pf:thmS-carrier:LE}, applied to
$\mathcal S^{\mathrm{mod}}$, needs:
(E1), which follows from the interlacing of the $\lambda_s$ (\lref{Formalism.FourierDecayAlt.P.roots_outside}); (E2), the
polynomiality of the energy in the letters (\lref{Formalism.FourierDecayAlt.P.first_identity}); (E3), the integrability of
the base density (\lref{Formalism.FourierDecayAlt.P.jacobi_square_integrable}); (E4), the representation of the energy of
boxed words as cut-off family energies (\cref{pf:thmS-carrier:thmPmodel}); (E5) as above; and (E6), the family hypotheses for
the model families (\cref{pf:thmS-carrier:famimports}). It gives
$\thmScarrierEn^{\mathrm{mod}}(\eps,w)\to\mathrm{size}(w)j_{\mathcal S^{\mathrm{mod}}}$, and
$j_{\mathcal S^{\mathrm{mod}}}=\int\sum_\rho|j_\rho|^2d\mu_0$, because the weight of the empty word is $1$, so that the base
density of the model system is $\sum_\rho|j_\rho|^2$.
\end{proof}

% ---------------------------------------------------------------------------
\subsection{Remarks on the formal statements}\label{sec:pf:thmS-carrier:remarks}
\begin{enumerate}
\item The real inputs do not involve the kernels $\mathcal K_i$, $\Theta_{i,s,t}$ and the numbers $e_i$ of the energy datum
  (\cref{pf:thmS-carrier:rem-statement}); \cref{pf:thmS-carrier:compact} proves them for all such data at once.
\item In \cref{pf:thmS-carrier:gram} the grid condition is not used, and in \cref{pf:thmS-carrier:xpins} the uniqueness of
  Bessel models is not used; the latter enters the proof of the main theorem through \cref{pf:thmS-carrier:xleaves}.
\item The energy constant in the conclusion of $\sfname{CarrierP}$ does not depend on the auxiliary data
  (\cref{pf:thmS-carrier:intN}).
\item Four auxiliary statements are supplied as typeclass instances (\cref{pf:thmS-carrier:rem-facts}); since they are found by
  instance resolution rather than by explicit reference, their proofs are not part of the part of the formal proof treated here.
\item The top lifts are produced with the explicit constant $C_b=11$ (\cref{pf:thmS-carrier:liftat3}).
\item At built places the model data can be taken with $C_m=\sigma=1$ (\cref{pf:thmS-carrier:openQQ}); the scaling constants
  allowed in the definition of model places are not needed there.
\item In \cref{pf:thmS-carrier:hatreal} the gauge condition is among the hypotheses but is not used; only the frame property, the
  radial gauge condition and the model level enter.
\item At the top level the germ rate holds with the trivial unimodular data $\eps'=1$, $\phi=1$
  (\cref{pf:thmS-carrier:germrate}).
\end{enumerate}

